% !TEX program = lualatex
\documentclass[11pt,twoside,openany]{book}
\usepackage{optitab-livre}
\usepackage{figures-geometrie}
\usepackage{relecture-livre}
\usepackage{edition-france}
\newcommand{\ligneSommaire}[2]{\par\vspace{3.5mm}\hyperref[#2]{#1}\hfill{\small p.\,\pageref*{#2}}\par}
\newcommand{\otDecoalgorithmique}{\otDecoboucles}
\newcommand{\otDecoprobabilites}{\otDecoconditionnelles}
\newcommand{\otDecogeometrie}{\otDecoGrille
\draw[white,opacity=.25,line width=1pt](113mm,24mm)--(194mm,24mm)--(194mm,100mm)--cycle;
\draw[white,opacity=.25,line width=.8pt](186mm,24mm)--(186mm,32mm)--(194mm,32mm);}
\hypersetup{pdftitle={Mathématiques - Troisième - Nouveau programme - Brevet 2029 - Livre complet},
pdfsubject={France, édition 2028-2029, programme 2026 : 18 chapitres, 288 exercices corrigés, deux brevets blancs},pdfauthor={OptiTAB}}
\newcommand{\otDecoseconddegre}{%
  \otDecoGrille
  \draw[white,opacity=.25] (132mm,47mm)--(210mm,47mm);
  \draw[white,opacity=.25] (171mm,18mm)--(171mm,118mm);
  \draw[otGreen,line width=1.8pt] plot[domain=-1.8:1.8,samples=75]
    ({(171+20*\x)*1mm},{(47+16*\x*\x)*1mm});
  \fill[white] (171mm,47mm) circle (1.5mm);}
\newcommand{\otDecovecteursplan}{%
  \otDecoGrille
  \draw[white,opacity=.7,line width=1.5pt,-{Stealth[length=3mm]}] (130mm,47mm)--(199mm,61mm);
  \draw[otGreen,line width=1.8pt,-{Stealth[length=3mm]}] (130mm,47mm)--(157mm,105mm);
  \draw[white,opacity=.35,dashed] (157mm,105mm)--(168mm,55mm);
  \fill[white] (130mm,47mm) circle (1.5mm);}
\newcommand{\otDecocercleplan}{%
  \otDecoGrille
  \draw[white,opacity=.25] (130mm,73mm)--(210mm,73mm) (170mm,36mm)--(170mm,116mm);
  \draw[otGreen,line width=1.8pt] (170mm,73mm) circle (30mm);
  \draw[white,opacity=.75,line width=1.2pt] (162mm,116.5mm)--(210mm,80.5mm);
  \draw[white,opacity=.45,dashed] (170mm,73mm)--(188mm,97mm);
  \fill[white] (188mm,97mm) circle (1.6mm);}
\newcommand{\otDecoensembles}{%
  \otDecoGrille
  \fill[white,opacity=.10] (155mm,76mm) circle (22mm);
  \fill[otGreen,opacity=.2] (181mm,76mm) circle (22mm);
  \draw[white,opacity=.8,line width=1.3pt] (155mm,76mm) circle (22mm);
  \draw[otGreen,line width=1.7pt] (181mm,76mm) circle (22mm);}

\begin{document}
\pagenumbering{alph}
%% =====================================================================
%%  1re de couverture
%% =====================================================================
\thispagestyle{empty}
\begin{tikzpicture}[remember picture,overlay]
\begin{scope}[shift={(current page.south west)}]
  % Fond
  \fill[otNavy] (0,0) rectangle (210mm,297mm);
  \shade[bottom color=otNavyDark,top color=otNavy] (0,0) rectangle (210mm,160mm);
  \draw[white,opacity=0.06,line width=0.3pt,step=7mm] (0,0) grid (210mm,297mm);
  % Parallélogramme et somme de deux vecteurs (relation de Chasles)
\draw[white,opacity=.22,line width=.6pt](98mm,32mm)--(160mm,32mm)--(194mm,88mm)--(132mm,88mm)--cycle;
\draw[white,opacity=.85,line width=1.6pt,-{Stealth[length=4.2mm,width=3.2mm]}](98mm,32mm)--(158.5mm,32mm);
\draw[white,opacity=.85,line width=1.6pt,-{Stealth[length=4.2mm,width=3.2mm]}](160mm,32mm)--(193.2mm,86.7mm);
\draw[otGreen,line width=3pt,-{Stealth[length=6mm,width=4.6mm]}](98mm,32mm)--(192.4mm,87.1mm);
\fill[white](98mm,32mm)circle(1.8mm);\fill[otGreen](98mm,32mm)circle(.9mm);
  % Logo
  \node[anchor=north west,inner sep=0pt] at (22mm,279mm)
    {\includegraphics[width=60mm]{logo-blanc}};

  % Titre
  \node[anchor=base west,inner sep=0pt,text=otGreen,
    font=\otSemi\fontsize{11}{13}\selectfont] at (24mm,236mm) {\otCaps[22]{\otNiveau}};
  \node[anchor=base west,inner sep=0pt,text=white,
    font=\otBold\fontsize{50}{54}\selectfont] at (21.5mm,212mm) {\otTitreLivre};
  \node[anchor=base west,inner sep=0pt,text=white,
    font=\otLight\fontsize{32}{36}\selectfont] at (23mm,194mm) {\otSousTitre};
  \fill[otGreen] (24mm,181mm) rectangle (54mm,182.4mm);
  \node[anchor=base west,inner sep=0pt,text=white,opacity=0.85,
    font=\otRegular\fontsize{11}{14}\selectfont] at (24mm,168mm)
    {Édition 2028-2029\enspace·\enspace Nouveau programme\enspace·\enspace Brevet 2029};

  % Points forts
  \node[anchor=north west,inner sep=0pt,text width=100mm,text=white,
    font=\otRegular\fontsize{10.5}{14}\selectfont] at (24mm,146mm)
    {\otCheck\hspace{3mm}Exercices classés par difficulté\par\vspace{3.2mm}
     \otCheck\hspace{3mm}La méthode expliquée à chaque étape\par\vspace{3.2mm}
     \otCheck\hspace{3mm}18 défis avec indications progressives\par\vspace{3.2mm}
     \otCheck\hspace{3mm}Deux brevets blancs originaux};

  % Pastille
  \fill[otGreen] (146mm,116mm) circle (18mm);
  \node[text=white,align=center,inner sep=0pt] at (146mm,116mm)
    {{\otMedium\fontsize{7}{9}\selectfont Programme 2026}\\[0.6mm]
     {\otBold\fontsize{30}{30}\selectfont 288}\\[1.4mm]
     {\otSemi\fontsize{7.5}{9}\selectfont\otCaps[10]{exercices}}\\
     {\otSemi\fontsize{7.5}{9}\selectfont\otCaps[10]{corrigés}}};

  % Pied
  \node[anchor=base west,inner sep=0pt,text=white,opacity=0.75,
    font=\otMedium\fontsize{10}{12}\selectfont] at (24mm,16mm) {\otSite};
\end{scope}
\end{tikzpicture}
\clearpage

\pagenumbering{arabic}
%% =====================================================================
%%  Page de titre
%% =====================================================================
\thispagestyle{empty}
\begin{tikzpicture}[remember picture,overlay]
\begin{scope}[shift={(current page.south west)}]
  \node[anchor=north,inner sep=0pt] at (105mm,262mm)
    {\includegraphics[width=64mm]{logo-couleur}};
  \node[anchor=base,inner sep=0pt,text=otGreenDark,
    font=\otSemi\fontsize{10}{12}\selectfont] at (105mm,182mm) {\otCaps[22]{\otNiveau}};
  \node[anchor=base,inner sep=0pt,text=otNavy,
    font=\otBold\fontsize{42}{46}\selectfont] at (105mm,163mm) {\otTitreLivre};
  \node[anchor=base,inner sep=0pt,text=otInk,
    font=\otLight\fontsize{24}{28}\selectfont] at (105mm,148mm) {\otSousTitre};
  \fill[otGreen] (93mm,138mm) rectangle (117mm,139.2mm);
  \node[anchor=base,inner sep=0pt,text=otMuted,
    font=\itshape\fontsize{12}{15}\selectfont] at (105mm,126mm)
    {288 exercices progressifs et deux brevets blancs corrigés};
  \node[anchor=base,inner sep=0pt,text=otGreenDark,
    font=\otSemi\fontsize{10.5}{13}\selectfont] at (105mm,112mm)
    {Édition 2028-2029\enspace·\enspace Nouveau programme\enspace·\enspace Brevet 2029};
  \node[anchor=base,inner sep=0pt,text=otInk,
    font=\otMedium\fontsize{10}{12}\selectfont] at (105mm,42mm)
    {Collection OptiTAB};
  \node[anchor=base,inner sep=0pt,text=otMuted,
    font=\otLight\fontsize{9}{11}\selectfont] at (105mm,35mm) {\otSite};
\end{scope}
\end{tikzpicture}
\clearpage

\thispagestyle{empty}\vspace*{\fill}
\begin{minipage}{128mm}\fontsize{9}{13}\selectfont\color{otMuted}
{\otSemi\color{otInk}Mathématiques · Troisième\par Édition 2028-2029 · Nouveau programme · Brevet 2029}\par\medskip
Conception et mise en page : collection OptiTAB.\par
270 exercices de chapitre, 18 défis et deux brevets blancs originaux, corrigés.\par
Charte graphique de la collection OptiTAB.\par\bigskip
Édition du 23 octobre 2026 · France · Année scolaire 2028-2029.\par\bigskip
Cette édition suit le programme de mathématiques du cycle 4 fixé par l’arrêté
du 18 février 2026 (BO n° 10 du 5 mars 2026), applicable en Troisième à la
rentrée 2028. Elle prépare la session 2029 du diplôme national du brevet.\par\bigskip
Au 23 octobre 2026, aucun texte propre à la session 2029 n’était publié. Les
brevets blancs suivent la définition de l’épreuve de mathématiques en vigueur
à compter de la session 2027 (note de service du 11 septembre 2026, BO spécial
n° 4 du 17 septembre 2026). Vérifier avant l’examen qu’aucun texte plus récent
ne l’a modifiée.\par\bigskip
L’ordre des chapitres, les exercices et les barèmes d’entraînement sont des
choix éditoriaux. Les défis sont facultatifs. Les situations et données
numériques sont construites pour l’apprentissage ; elles ne constituent pas
des données économiques ou scientifiques relevées.\par\bigskip
Les schémas de principe ne permettent pas de mesurer une longueur inconnue.
Les coordonnées et graduations explicites font foi pour les graphiques.
Sauf demande d’arrondi, conserver les résultats exacts.\par\bigskip
© OptiTAB, 2026.\par\otSite
\end{minipage}\clearpage

\clearpage\thispagestyle{otFront}\phantomsection\label{avant-propos}\pdfbookmark[0]{Apprendre et comprendre}{bm-avant-propos}\otPageTitre{Apprendre et comprendre}{Un livre pour toute l’année de Troisième}\fontsize{11}{15}\selectfont

Ce livre accompagne la classe de Troisième et la préparation du brevet 2029,
avec le nouveau programme de mathématiques du cycle 4. Il aide à consolider les
calculs, à comprendre les nouvelles notions, puis à résoudre des problèmes en
reliant plusieurs outils. Il complète le travail de classe, dans l’ordre adapté aux besoins.\par\medskip
\begin{savoirfaire}[Une progression en trois temps]Chaque chapitre propose deux pages de cours, cinq exercices de découverte, cinq d’entraînement et cinq problèmes. Les corrections explicitent la méthode, les calculs et l’interprétation. Une figure utile est reprise avec le corrigé ; une construction cherchée n’apparaît qu’à cet endroit.\end{savoirfaire}
\begin{savoirfaire}[Chercher avant de consulter]Écrire une première démarche même incomplète. Après correction, fermer le livre et refaire les étapes décisives. Quelques jours plus tard, reprendre un exercice proche sans aide : reconnaître un résultat n’est pas encore savoir le retrouver.\end{savoirfaire}
\begin{savoirfaire}[Aller plus loin]Les dix-huit défis demandent une initiative supplémentaire, avec deux indications séparées et une solution rédigée. Deux brevets blancs permettent un entraînement chronométré. Les fiches et la matrice finale aident à repérer la notion à revoir.\end{savoirfaire}
Les six compétences travaillées sont chercher, modéliser, représenter,
raisonner, calculer et communiquer. Un dessin, un tableau ou un script sont
des moyens de penser autant que des supports de réponse.

\clearpage

\clearpage\thispagestyle{otFront}\phantomsection\label{mode-emploi}\pdfbookmark[0]{Mode d’emploi}{bm-mode-emploi}\otPageTitre{Mode d’emploi}{Des repères stables dans tout le livre}\fontsize{11}{15}\selectfont

\begin{tabularx}{\linewidth}{@{}lY@{}}\toprule
Une pastille & Découvrir : appliquer une propriété, lire une représentation ou faire un calcul direct.\\[3mm]
Deux pastilles & S’entraîner : choisir une méthode, combiner deux étapes ou interpréter un résultat.\\[3mm]
Trois pastilles & Résoudre : organiser une démarche plus longue, justifier et contrôler.\\\bottomrule
\end{tabularx}\par\medskip
Ces niveaux sont des repères pédagogiques, pas un classement des élèves.
Un exercice simple peut révéler une notion à consolider ; un problème peut
être réussi avec une représentation bien choisie.\par\medskip
\begin{exemple}[Une correction utile]Pour résoudre $3(x-2)=x+4$, on développe : $3x-6=x+4$, puis $2x=10$, donc $x=5$. Le contrôle dans l’énoncé donne $3(5-2)=9=5+4$. L’égalité obtenue vérifie la solution.\end{exemple}
\begin{attention}Lire tout l’énoncé, y compris les unités, les mots « exactement », « au moins », « et », « ou » et la précision demandée. En géométrie, un angle droit dessiné sans hypothèse n’est pas une preuve.\end{attention}
Les liens du sommaire et les numéros de page conduisent aux chapitres et
aux corrigés. Le tableau de bord permet de noter ce qui est réussi seul,
réussi avec aide ou à reprendre. Pour les scripts, les blocs indentés
appartiennent au bloc qui les contient ; les lignes « fin » en indiquent la fermeture.

\clearpage

\clearpage\thispagestyle{otFront}\phantomsection\label{parcours}\pdfbookmark[0]{Construire son parcours}{bm-parcours}\otPageTitre{Construire son parcours}{Avancer avec des retours réguliers}\fontsize{11}{15}\selectfont

\begin{tabularx}{\linewidth}{@{}lY@{}}\toprule
Étape & Travail conseillé\\\midrule
Consolider & Chapitres 1 à 5 : nombres, puissances, calcul littéral, équations et inéquations.\\[3mm]
Modéliser & Chapitres 6 à 11 : proportionnalité, fonctions, statistiques et probabilités.\\[3mm]
Démontrer & Chapitres 12 à 16 : théorèmes du triangle, vecteurs, espace.\\[3mm]
Programmer & Chapitre 17, à reprendre pendant toute l’année.\\[3mm]
Synthétiser & Chapitre 18, défis et brevets blancs.\\\bottomrule
\end{tabularx}
\begin{savoirfaire}[Une séance autonome de 35 à 50 minutes]Relire une propriété et son exemple ; traiter deux exercices sans correction ; comparer la démarche ; refaire une question ratée. Terminer par une phrase : « La prochaine fois, je dois penser à… ».\end{savoirfaire}
\begin{savoirfaire}[Une révision espacée]Revenir sur une notion quelques jours après, puis quelques semaines plus tard. Mélanger un calcul, une lecture graphique et une justification géométrique évite de dépendre seulement de l’ordre du chapitre.\end{savoirfaire}
Pour un défi, prévoir 30 à 60 minutes, éventuellement en deux séances.
Pour un brevet blanc, utiliser deux copies et respecter la remise de la
première au bout de 20 minutes. Refaire ensuite les erreurs par notion,
sans attendre le prochain sujet complet.

\clearpage

\clearpage\thispagestyle{otFront}\phantomsection\label{nouveautes}\pdfbookmark[0]{Nouveau programme}{bm-nouveautes}\otPageTitre{Nouveau programme}{Ce qui change en Troisième à la rentrée 2028}\fontsize{11}{15}\selectfont

Le programme de mathématiques du cycle 4 publié en 2026 s’applique en Troisième
à la rentrée 2028. Cette édition a été réécrite en conséquence ; voici les
principaux changements par rapport au programme précédent.\par\medskip
\fontsize{10}{13.5}\selectfont
\begin{tabularx}{\linewidth}{@{}>{\raggedright\arraybackslash}p{27mm}Y@{}}\toprule
Nouveau en Troisième & Translations et vecteurs : vecteurs égaux, somme, relation de Chasles (ch. 15). Effectifs cumulés, quartiles et boîtes à moustaches (ch. 10). Réunion, intersection et relation $P(A\cup B)+P(A\cap B)=P(A)+P(B)$ (ch. 11). Inéquations $ax\ge b$ par le calcul et graphiquement (ch. 5 et 9). Trois identités remarquables et raisonnement par analyse-synthèse (ch. 4 et 5). Fonction carré et équation $x^2=a$ (ch. 2 et 8). Sphère, boule, grands cercles et sections (ch. 16). Conditions composées, boucles conditionnelles et blocs personnalisés (ch. 17).\\[3mm]
Prolongements & Notations $\mathbb N$, $\mathbb Z$, $\mathbb D$, $\mathbb Q$ (ch. 1) ; préfixes de nano à giga (ch. 2) ; identité de Sophie Germain et cylindre d’Archimède (défis et ch. 16).\\[3mm]
Réactivé & Pythagore, sa réciproque et sa contraposée, cercle circonscrit au triangle rectangle, droite des milieux, grandeurs quotients et expériences à deux épreuves : étudiés en Quatrième, ils sont attendus en automatismes et dans les problèmes (ch. 7, 11 et 12).\\[3mm]
Retiré & Homothéties et rotations, triangles semblables, effet d’un agrandissement sur les aires et les volumes, repérage sur la sphère et dans un pavé, histogrammes : ces notions ne figurent plus dans le programme du cycle 4 et ne sont plus traitées.\\\bottomrule
\end{tabularx}
\par\medskip\fontsize{11}{15}\selectfont
Le vocabulaire suit le programme : quotient, fraction et nombre rationnel ;
demi-tour pour la symétrie centrale ; théorème, réciproque et contraposée ;
évènement, réunion et intersection. La notation $x\mapsto f(x)$ est utilisée
pour les fonctions. La matrice de fin de livre relie chaque rubrique du
programme aux chapitres et aux exercices.

\clearpage

\thispagestyle{otFront}\phantomsection\pdfbookmark[0]{Sommaire}{bm-sommaire}
\otPageTitre{Sommaire}{Troisième · Nouveau programme · Brevet 2029}
\vspace{-4mm}\sommaire\clearpage
\thispagestyle{otFront}\otPageTitre{Parcours et ressources}{Cours, défis, sujets et corrections}
\ligneSommaire{Nouveau programme : ce qui change}{nouveautes}
\ligneSommaire{Dix-huit défis d’approfondissement}{defis-debut}
\ligneSommaire{Indications progressives}{indices-debut}
\ligneSommaire{Organisation des deux brevets blancs}{bilans-guide}
\ligneSommaire{Brevet blanc A}{bilan-A-auto}
\ligneSommaire{Brevet blanc B}{bilan-B-auto}
\ligneSommaire{Comprendre quelques démonstrations}{preuves-debut}
\ligneSommaire{Corrigés des 270 exercices}{corriges-debut}
\ligneSommaire{Corrigés des 18 défis}{defi-cor-1}
\ligneSommaire{Corrigés des brevets blancs}{bilan-A-cor-auto}
\ligneSommaire{Fiches outils}{outils-debut}
\ligneSommaire{Correspondance avec le programme}{couverture-programme}
\ligneSommaire{Sources officielles et périmètre de l’édition}{references}
\clearpage

%% =====================================================================
%%  Mon tableau de bord (suivi de progression)
%% =====================================================================
\thispagestyle{otFront}
\phantomsection\pdfbookmark[0]{Mon tableau de bord}{bm-tableau-de-bord}
\otPageTitre{Mon tableau de bord}{}
\vspace{-4mm}
{\fontsize{10}{13.5}\selectfont
Après chaque exercice, colorie sa pastille :\enspace
\tikz[baseline=-0.6ex]\fill[otGreen] (0,0) circle (1.75mm);\ réussi seul\quad
\tikz[baseline=-0.6ex]\fill[otOrange] (0,0) circle (1.75mm);\ réussi avec le corrigé\quad
\tikz[baseline=-0.6ex]\fill[otRed] (0,0) circle (1.75mm);\ à refaire\par}
\vspace{1mm}
\tableauDeBord

\clearpage

\pagestyle{otMain}
\begin{ouverture}[courbe]{1}{Nombres rationnels et fractions}{Partie 1 · Nombres et calculs}
{Vocabulaire · Priorités · Fractions irréductibles · Problèmes}
\objectif{Calculer en respectant les priorités et les signes.}
\objectif{Rendre une fraction irréductible et calculer exactement.}
\objectif{Résoudre un problème faisant appel à des fractions.}
\prerequis{Tables de multiplication, nombres relatifs et opérations sur les fractions.}
\end{ouverture}
\pagestyle{otMain}\essentiel\fontsize{11}{15}\selectfont
\essSection{1}{Vocabulaire et priorités}
\begin{definition}Un quotient est le résultat d’une division, ou l’écriture de cette division. Une fraction est le quotient $a/b$ de deux entiers, avec $b\ne0$ : elle désigne à la fois un nombre et une écriture. Un nombre rationnel est un nombre égal au quotient de deux entiers.\end{definition}
\begin{propriete}Les parenthèses sont prioritaires, puis les puissances, les produits et quotients, enfin les sommes et différences. À priorité égale, on calcule de gauche à droite.\end{propriete}
\begin{exemple}$-4+3\times(5-7)=-4+3\times(-2)=-10$. Pour $(-3)^2$, le carré porte sur le nombre $-3$ ; dans $-3^2$, le signe moins est extérieur au carré.\end{exemple}
\begin{savoirfaire}[Gérer les signes]Un produit ou un quotient de deux nombres de même signe est positif ; de signes contraires, il est négatif. Pour une somme, ces règles ne s’appliquent pas : comparer les valeurs et les signes.\end{savoirfaire}
\begin{definition}[Ensembles de nombres]$\mathbb N$ désigne les entiers naturels, $\mathbb Z$ les entiers relatifs, $\mathbb D$ les nombres décimaux et $\mathbb Q$ les nombres rationnels. Chacun est contenu dans le suivant. Ainsi $-3\in\mathbb Z$, $2{,}5\in\mathbb D$ et $1/3\in\mathbb Q$, mais $1/3$ n’est pas décimal.\end{definition}
\begin{attention}On ne divise jamais par zéro. Une égalité doit relier des nombres égaux : ne pas enchaîner un résultat partiel avec tout le calcul initial.\end{attention}
\clearpage
\essSection{2}{Calculer avec des fractions}
\begin{propriete}Pour additionner, utiliser un dénominateur commun. Pour multiplier, multiplier les numérateurs et les dénominateurs. Pour diviser par une fraction non nulle, multiplier par son inverse.\end{propriete}
\begin{exemple}
\[\frac34-\frac56=\frac9{12}-\frac{10}{12}=-\frac1{12},
\qquad \frac23\div\frac57=\frac23\times\frac75=\frac{14}{15}.\]
\end{exemple}
\begin{definition}Une fraction est irréductible lorsque son numérateur et son dénominateur n’ont pas d’autre diviseur commun positif que 1.\end{definition}
\begin{savoirfaire}[Rendre irréductible]Simplifier successivement par des diviseurs communs, en s’aidant des critères de divisibilité, ou décomposer les deux termes en facteurs premiers. $\dfrac{84}{126}=\dfrac{42}{63}=\dfrac{14}{21}=\dfrac23$ : on a simplifié par 2, par 3, puis par 7.\end{savoirfaire}
\begin{savoirfaire}[Une fraction d’une quantité]Prendre $3/5$ de 40 revient à calculer $40\times3/5=24$. « Les $2/3$ du reste » impose de calculer d’abord le reste : le tout de référence a changé. Un schéma en barres aide à suivre ces parts.\end{savoirfaire}
\begin{attention}On simplifie des facteurs dans un produit, jamais des termes isolés d’une somme. Ainsi $(6+3)/3=3$, et non $6+1$.\end{attention}
\clearpage\rubrique{Découverte}
\begin{blocLivre}{ex-1-1}\begin{exercice}[1]{Une température}
Le matin, il fait $-6\,^\circ$C. La température augmente de 11 degrés puis baisse de 4 degrés. Quelle est la température finale ?
\end{exercice}\end{blocLivre}
\begin{blocLivre}{ex-1-2}\begin{exercice}[1]{Respecter les priorités}
Calculer $A=18-3(2+5)$ et $B=(18-3)(2+5)$.
\end{exercice}\end{blocLivre}
\begin{blocLivre}{ex-1-3}\begin{exercice}[1]{Deux écritures, un nombre}
Compléter $\dfrac5{8}=\dfrac{\,?\,}{40}$ et $\dfrac{42}{63}=\dfrac2{\,?\,}$.
\end{exercice}\end{blocLivre}
\begin{blocLivre}{ex-1-4}\begin{exercice}[1]{Une somme exacte}
Calculer $\dfrac7{12}+\dfrac5{18}$ et donner une fraction irréductible.
\end{exercice}\end{blocLivre}
\begin{blocLivre}{ex-1-5}\begin{exercice}[1]{Produit et quotient}
Calculer $\dfrac{-9}{14}\times\dfrac7{15}$ puis $\dfrac35\div\dfrac{-9}{10}$.
\end{exercice}\end{blocLivre}

\clearpage\rubrique{Entraînement}
\begin{blocLivre}{ex-1-6}\begin{exercice}[2]{Simplifications successives}
Rendre irréductibles $\dfrac{150}{225}$ et $\dfrac{126}{198}$ par simplifications successives. Les deux fractions sont-elles égales ?
\end{exercice}\end{blocLivre}
\begin{blocLivre}{ex-1-7}\begin{exercice}[2]{Comparer sans arrondir}
Ranger $-\dfrac34$, $-\dfrac56$ et $-\dfrac7{12}$ dans l’ordre croissant.
\end{exercice}\end{blocLivre}
\begin{blocLivre}{ex-1-8}\begin{exercice}[2]{Une expression à étages}
Calculer exactement $C=\dfrac{\frac23+\frac16}{\frac54-\frac12}$.
\end{exercice}\end{blocLivre}
\begin{blocLivre}{ex-1-9}\begin{exercice}[2]{Des ensembles de nombres}
Pour chacun des nombres $-7$ ; $0$ ; $2{,}5$ ; $\dfrac{12}{4}$ ; $-\dfrac53$, donner le plus petit des ensembles $\mathbb N$, $\mathbb Z$, $\mathbb D$, $\mathbb Q$ qui le contient.
\end{exercice}\end{blocLivre}
\begin{blocLivre}{ex-1-10}\begin{exercice}[2]{Une réduction trompeuse}
Un élève écrit $(8+4)/4=8+1$. Expliquer son erreur et calculer correctement.
\end{exercice}\end{blocLivre}

\clearpage\rubrique{Problèmes}
\begin{blocLivre}{ex-1-11}\begin{exercice}[3]{Deux étapes de lecture}
Un lecteur lit $2/5$ d’un livre le lundi, puis $1/3$ de ce qui reste le mardi. Quelle fraction du livre reste-t-il ? Combien de pages pour un livre de 270 pages ?
\end{exercice}\end{blocLivre}
\begin{blocLivre}{ex-1-12}\begin{exercice}[3]{Des bouquets vendus}
Une fleuriste vend les $2/3$ de ses bouquets le matin, puis les $3/4$ du reste l’après-midi. Le soir, il lui reste 9 bouquets. Combien en avait-elle au début ?
\end{exercice}\end{blocLivre}
\begin{blocLivre}{ex-1-13}\begin{exercice}[3]{Une réserve d’eau}
Une cuve est remplie aux $7/8$. On consomme $2/7$ de l’eau présente, puis on ajoute 30 L. La cuve est alors remplie aux $3/4$. Quelle est sa capacité ?
\end{exercice}\end{blocLivre}
\begin{blocLivre}{ex-1-14}\begin{exercice}[3]{Des votes et des absents}
Dans une classe, $5/6$ des élèves sont présents. Parmi eux, $3/5$ choisissent le projet A. Le projet A reçoit 15 voix. Combien y a-t-il d’élèves et d’absents ?
\end{exercice}\end{blocLivre}
\begin{blocLivre}{ex-1-15}\begin{exercice}[3]{Trouver toutes les possibilités}
On cherche les entiers naturels $n$ tels que $n/12$ soit strictement supérieur à $2/3$ et strictement inférieur à $5/4$. Les donner tous.
\end{exercice}\end{blocLivre}

\clearpage

\begin{ouverture}[courbe]{2}{Puissances, racines et notation scientifique}{Partie 1 · Nombres et calculs}
{Exposants · Préfixes · Ordres de grandeur · Équation $x^2=a$}
\objectif{Utiliser les puissances d’exposant entier.}
\objectif{Écrire et calculer en notation scientifique.}
\objectif{Résoudre une équation $x^2=a$ et un problème avec une racine carrée.}
\prerequis{Produits, quotients et carrés des entiers de 0 à 12.}
\end{ouverture}
\pagestyle{otMain}\essentiel\fontsize{11}{15}\selectfont
\essSection{1}{Puissances et notation scientifique}
\begin{definition}Pour $n$ entier positif, $a^n$ est le produit de $n$ facteurs égaux à $a$. Pour $a\ne0$, on pose $a^0=1$ et $a^{-n}=1/a^n$ : en divisant $10^3$ successivement par 10, on obtient $10^2$, $10^1$, $10^0=1$, puis $10^{-1}=0{,}1$.\end{definition}
\begin{propriete}Pour $a$ et $b$ non nuls et $m$, $n$ entiers relatifs :
\[a^m\times a^n=a^{m+n},\qquad \frac{a^m}{a^n}=a^{m-n},\qquad a^n\times b^n=(ab)^n.\]\end{propriete}
\begin{definition}La notation scientifique d’un nombre positif est $c\times10^p$, avec $1\le c<10$ et $p$ entier relatif. Par exemple $0{,}000056=5{,}6\times10^{-5}$.\end{definition}
\begin{savoirfaire}[Calculer en notation scientifique]Regrouper les coefficients, puis les puissances de 10. $(3\times10^5)(4\times10^{-2})=12\times10^3=1{,}2\times10^4$.\end{savoirfaire}
\begin{propriete}[Préfixes]nano : $10^{-9}$ ; micro (µ) : $10^{-6}$ ; milli : $10^{-3}$ ; kilo : $10^{3}$ ; méga : $10^{6}$ ; giga : $10^{9}$. Ainsi $1\text{ nm}=10^{-9}\text{ m}$ et $1\text{ Go}=10^9$ octets.\end{propriete}
\begin{attention}$a^m+a^n$ ne se transforme pas en $a^{m+n}$. Les règles sur les exposants concernent des produits et des quotients.\end{attention}
\clearpage
\essSection{2}{Racine carrée et équation $x^2=a$}
\begin{definition}Pour $a\ge0$, $\sqrt a$ est l’unique nombre positif ou nul dont le carré vaut $a$. Ainsi $(\sqrt a)^2=a$. La racine carrée d’un nombre négatif n’est pas définie.\end{definition}
\begin{propriete}[Résoudre $x^2=a$]Si $a>0$, l’équation a deux solutions, $-\sqrt a$ et $\sqrt a$ : l’ensemble des solutions est $\{-\sqrt a\,;\sqrt a\}$. Si $a=0$, la seule solution est 0. Si $a<0$, il n’y a aucune solution : l’ensemble des solutions est vide, noté $\varnothing$.\end{propriete}
\begin{exemple}Pour $x^2=9$, on écrit $x^2-9=(x-3)(x+3)=0$ : les solutions sont $-3$ et 3. Graphiquement, la droite d’équation $y=9$ coupe la courbe de la fonction carré en deux points d’abscisses $-3$ et 3.\end{exemple}
\begin{savoirfaire}[Arrondir une mesure]Un carré d’aire 45 cm$^2$ a pour côté $\sqrt{45}$ cm, environ $6{,}71$ cm : une longueur est positive, on ne garde pas $-\sqrt{45}$. Conserver les valeurs exactes pendant les calculs et arrondir à la fin, à la précision demandée.\end{savoirfaire}
\begin{attention}$\sqrt{(-4)^2}=4$ : une racine carrée est toujours positive ou nulle. Et $\sqrt{a+b}$ n’est pas, en général, égal à $\sqrt a+\sqrt b$.\end{attention}
\clearpage\rubrique{Découverte}
\begin{blocLivre}{ex-2-1}\begin{exercice}[1]{Comprendre un exposant}
Écrire $7\times7\times7\times7$ comme une puissance. Calculer $2^{-3}$ et $9^0$.
\end{exercice}\end{blocLivre}
\begin{blocLivre}{ex-2-2}\begin{exercice}[1]{Parenthèses et signe}
Calculer $(-5)^2$, $-5^2$, $(-2)^3$ et $(-2)^4$.
\end{exercice}\end{blocLivre}
\begin{blocLivre}{ex-2-3}\begin{exercice}[1]{Règles de puissances}
Écrire sous la forme d’une seule puissance : $3^4\times3^2$, $\dfrac{5^7}{5^3}$, $2^3\times5^3$ et $\dfrac{10^{-2}}{10^{4}}$.
\end{exercice}\end{blocLivre}
\begin{blocLivre}{ex-2-4}\begin{exercice}[1]{Écriture scientifique}
Écrire $74000000$ et $0{,}000092$ en notation scientifique.
\end{exercice}\end{blocLivre}
\begin{blocLivre}{ex-2-5}\begin{exercice}[1]{Quelques racines}
Calculer $\sqrt{144}$, $\sqrt{0{,}81}$ et $\sqrt{49/64}$.
\end{exercice}\end{blocLivre}

\clearpage\rubrique{Entraînement}
\begin{blocLivre}{ex-2-6}\begin{exercice}[2]{Produit scientifique}
Calculer $(6\times10^8)(4\times10^{-5})$ en notation scientifique.
\end{exercice}\end{blocLivre}
\begin{blocLivre}{ex-2-7}\begin{exercice}[2]{Quotient scientifique}
Calculer $\dfrac{8{,}4\times10^{-3}}{2{,}1\times10^5}$.
\end{exercice}\end{blocLivre}
\begin{blocLivre}{ex-2-8}\begin{exercice}[2]{Des préfixes}
Écrire en notation scientifique, dans l’unité indiquée : $4{,}7$ Go en octets ; 250 nm en mètres ; $0{,}6$ µs en secondes.
\end{exercice}\end{blocLivre}
\begin{blocLivre}{ex-2-9}\begin{exercice}[2]{Une équation $x^2=a$}
Résoudre $x^2=49$, $x^2=7$, $x^2=0$ et $x^2=-4$. Donner chaque ensemble de solutions.
\end{exercice}\end{blocLivre}
\begin{blocLivre}{ex-2-10}\begin{exercice}[2]{Réfuter une règle}
Réfuter $\sqrt{a+b}=\sqrt a+\sqrt b$ en choisissant deux nombres positifs simples.
\end{exercice}\end{blocLivre}

\clearpage\rubrique{Problèmes}
\begin{blocLivre}{ex-2-11}\begin{exercice}[3]{Une population qui double}
Une culture contient 800 cellules. Leur nombre double toutes les 20 minutes, selon un modèle valable pendant l’expérience. Combien de cellules après 2 heures ? Écrire le résultat en notation scientifique.
\end{exercice}\end{blocLivre}
\begin{blocLivre}{ex-2-12}\begin{exercice}[3]{Une masse minuscule}
Un lot de $4\times10^{12}$ particules a une masse totale de $1{,}8$ g. Dans ce modèle, toutes ont la même masse. Calculer la masse d’une particule en grammes, puis en kilogrammes.
\end{exercice}\end{blocLivre}
\begin{blocLivre}{ex-2-13}\begin{exercice}[3]{Une distance astronomique}
Un signal parcourt $3\times10^5$ km par seconde. Il voyage pendant 8 min 20 s. Quelle distance parcourt-il ? Donner aussi cette distance en mètres.
\end{exercice}\end{blocLivre}
\begin{blocLivre}{ex-2-14}\begin{exercice}[3]{Choisir le côté d’une dalle}
Une dalle carrée a une aire de 2 m$^2$. Donner son côté exact en mètres, puis l’arrondir au millimètre. Peut-on le remplacer par 1,4 m pour calculer une aire exacte ?
\end{exercice}\end{blocLivre}
\begin{blocLivre}{ex-2-15}\begin{exercice}[3]{Un cercle et trois droites}
Dans un repère orthonormé, le cercle de centre $O$ et de rayon 5 contient les points $M(x\,;y)$ tels que $OM=5$, c’est-à-dire $x^2+y^2=25$ d’après Pythagore. On le coupe par les droites horizontales d’équations $y=3$, $y=5$ et $y=6$. Pour chacune, résoudre l’équation obtenue en $x$ et conclure.
\label{fig-ex-2-15}\figureGeo{cercle-droites}
\end{exercice}\end{blocLivre}

\clearpage

\begin{ouverture}[courbe]{3}{Multiples, diviseurs et nombres premiers}{Partie 1 · Nombres et calculs}
{Divisibilité · Décomposition · Fractions irréductibles · Périodicités}
\objectif{Décomposer un entier en facteurs premiers.}
\objectif{Simplifier une fraction et résoudre des partages.}
\objectif{Organiser une recherche de multiples communs.}
\prerequis{Division euclidienne et critères de divisibilité par 2, 3, 5 et 9.}
\end{ouverture}
\pagestyle{otMain}\essentiel\fontsize{11}{15}\selectfont
\essSection{1}{Divisibilité et nombres premiers}
\begin{definition}Un entier $a$ est un multiple d’un entier non nul $b$ lorsque $a=bk$ pour un entier $k$. Alors $b$ est un diviseur de $a$. Un nombre premier est un entier supérieur ou égal à 2 ayant exactement deux diviseurs positifs : 1 et lui-même.\end{definition}
\begin{exemple}2, 3, 5, 7, 11 et 13 sont premiers. Le nombre 1 ne l’est pas. Le nombre 91 ne l’est pas non plus : $91=7\times13$.\end{exemple}
\begin{propriete}Un entier est divisible par 2 si son dernier chiffre est pair, par 5 s’il se termine par 0 ou 5, par 3 ou par 9 si la somme de ses chiffres l’est. Un nombre peut vérifier plusieurs critères.\end{propriete}
\begin{savoirfaire}[Décomposer]Diviser successivement par des nombres premiers. Par exemple $60=2^2\times3\times5$ et $360=2^3\times3^2\times5$. Continuer jusqu’à ce que tous les facteurs soient premiers.\end{savoirfaire}
\begin{attention}Pour prouver qu’un nombre est composé, un seul diviseur autre que 1 et lui-même suffit. Pour prouver qu’il est premier, il faut exclure tous les diviseurs nécessaires, pas seulement 2 et 3.\end{attention}
\clearpage
\essSection{2}{Simplifier et organiser des groupes}
\begin{definition}Une fraction d’entiers est irréductible lorsque numérateur et dénominateur n’ont aucun diviseur commun positif autre que 1.\end{definition}
\begin{exemple}$84/150=(2^2\times3\times7)/(2\times3\times5^2)=14/25$. Les facteurs communs ont été simplifiés ; 14 et 25 n’en ont plus.\end{exemple}
\begin{savoirfaire}[Le plus grand nombre de groupes identiques]Le nombre de groupes doit diviser chacun des stocks. Chercher leur plus grand diviseur commun, à l’aide des décompositions ou de la liste des diviseurs.\end{savoirfaire}
\begin{savoirfaire}[La première rencontre de phénomènes périodiques]Le temps écoulé doit être un multiple de chaque période. Chercher le plus petit multiple commun strictement positif. Calendriers, engrenages et cycles d’éclosion d’espèces animales relèvent de ce raisonnement.\end{savoirfaire}
\begin{attention}Le mot « plus grand » ou « premier » demande une justification de l’optimalité. Donner un partage possible ou une rencontre possible ne suffit pas.\end{attention}
\otIntertitre{Avec un tableur}
Une colonne d’entiers et une colonne de restes permettent de tester les diviseurs. Contrôler la liste obtenue avec une décomposition en facteurs premiers.
\clearpage\rubrique{Découverte}
\begin{blocLivre}{ex-3-1}\begin{exercice}[1]{Tester des diviseurs}
Parmi 2, 3, 5, 9 et 10, lesquels divisent 2340 ?
\end{exercice}\end{blocLivre}
\begin{blocLivre}{ex-3-2}\begin{exercice}[1]{Premier ou composé ?}
Parmi 1, 2, 17, 21 et 49, relever les nombres premiers et justifier les exclusions.
\end{exercice}\end{blocLivre}
\begin{blocLivre}{ex-3-3}\begin{exercice}[1]{Décomposer un entier}
Décomposer 420 en facteurs premiers.
\end{exercice}\end{blocLivre}
\begin{blocLivre}{ex-3-4}\begin{exercice}[1]{Une fraction irréductible}
Rendre $198/252$ irréductible.
\end{exercice}\end{blocLivre}
\begin{blocLivre}{ex-3-5}\begin{exercice}[1]{Écrire des multiples}
Lister les multiples de 14 compris entre 70 et 130, bornes comprises.
\end{exercice}\end{blocLivre}

\clearpage\rubrique{Entraînement}
\begin{blocLivre}{ex-3-6}\begin{exercice}[2]{Tous les diviseurs}
Donner tous les diviseurs positifs de 72 et expliquer pourquoi la liste est complète.
\end{exercice}\end{blocLivre}
\begin{blocLivre}{ex-3-7}\begin{exercice}[2]{Un chiffre inconnu}
Déterminer les chiffres $a$ pour lesquels le nombre $45a$ est divisible par 3 et par 2.
\end{exercice}\end{blocLivre}
\begin{blocLivre}{ex-3-8}\begin{exercice}[2]{Une grande fraction}
Réduire $\dfrac{2^4\times3^2\times7}{2^2\times3^3\times5}$.
\end{exercice}\end{blocLivre}
\begin{blocLivre}{ex-3-9}\begin{exercice}[2]{Un carré parfait ?}
Le nombre $2^4\times3^3$ est-il un carré d’entier ? Par quel plus petit entier positif le multiplier pour obtenir un carré ?
\end{exercice}\end{blocLivre}
\begin{blocLivre}{ex-3-10}\begin{exercice}[2]{Un calcul par tableur}
Une colonne A contient les entiers de 1 à 20. La cellule B2 contient la formule \code{=MOD(84;A2)}. Que signifie un résultat nul ? Quels diviseurs de 84 inférieurs ou égaux à 20 obtient-on ?
\end{exercice}\end{blocLivre}

\clearpage\rubrique{Problèmes}
\begin{blocLivre}{ex-3-11}\begin{exercice}[3]{Préparer des lots}
On possède 156 cartes rouges et 204 cartes bleues. Former le plus grand nombre de lots identiques sans reste, avec les deux couleurs dans chaque lot. Donner le nombre de lots et leur contenu.
\end{exercice}\end{blocLivre}
\begin{blocLivre}{ex-3-12}\begin{exercice}[3]{Des cigales périodiques}
Dans un modèle simplifié, une espèce de cigales sort de terre tous les 13 ans et une autre tous les 17 ans. Les deux sont sorties la même année en 2024. En quelle année cette coïncidence se reproduira-t-elle pour la première fois ?
\end{exercice}\end{blocLivre}
\begin{blocLivre}{ex-3-13}\begin{exercice}[3]{Daller un rectangle}
Un rectangle de 180 cm sur 252 cm doit être recouvert de carrés identiques, sans découpe. Les côtés des carrés sont des nombres entiers de centimètres. Trouver le plus grand côté possible et le nombre de carrés.
\label{fig-ex-3-13}\figureGeo{dallage}
\end{exercice}\end{blocLivre}
\begin{blocLivre}{ex-3-14}\begin{exercice}[3]{Un stock à retrouver}
Un stock compte entre 500 et 600 objets. Il se range sans reste par paquets de 8, de 9 ou de 15. Déterminer toutes les possibilités.
\end{exercice}\end{blocLivre}
\begin{blocLivre}{ex-3-15}\begin{exercice}[3]{Des engrenages}
Deux roues de 24 et 40 dents engrènent. Deux dents marquées sont en contact au départ. Après combien de passages de dents retrouve-t-on ce contact ? Combien de tours fait chaque roue ?
\label{fig-ex-3-15}\figureGeo{engrenages}
\end{exercice}\end{blocLivre}

\clearpage

\begin{ouverture}[courbe]{4}{Calcul littéral et identités remarquables}{Partie 1 · Nombres et calculs}
{Expressions · Distributivité · Identités remarquables · Preuves}
\objectif{Transformer une expression sans changer sa valeur.}
\objectif{Développer et factoriser avec les identités remarquables.}
\objectif{Démontrer une propriété générale avec une lettre.}
\prerequis{Produits, sommes, puissances et calculs avec les relatifs.}
\end{ouverture}
\pagestyle{otMain}\essentiel\fontsize{11}{15}\selectfont
\essSection{1}{Lire, simplifier et développer}
\begin{definition}Une expression littérale décrit un calcul avec des lettres. Sa nature est donnée par la dernière opération effectuée : $3x+2$ est une somme, $5(x+4)$ est un produit. Développer transforme un produit en somme ; factoriser transforme une somme en produit.\end{definition}
\begin{propriete}La distributivité donne $k(a+b)=ka+kb$. En particulier, $-(a-b)=-a+b$ : prendre l’opposé d’une expression change le signe de chacun de ses termes. Pour deux parenthèses :
\[(a+b)(c+d)=ac+ad+bc+bd.\]\end{propriete}
\begin{exemple}$(2x-3)(x+4)=2x^2+8x-3x-12=2x^2+5x-12$. Le terme $2x^2$ est un carré multiplié par 2, pas un terme en $x$.\end{exemple}
\begin{savoirfaire}[Simplifier un produit ou un rapport]Regrouper les facteurs de même nature : $(3a^2b)\times(2bc)=6a^2b^2c$. Dans un rapport, simplifier les facteurs communs non nuls : $\dfrac{6ab^2}{9a^2b}=\dfrac{2b}{3a}$.\end{savoirfaire}
\begin{attention}$3x+2x=5x$, mais $3x+2$ ne se réduit pas à $5x$. Un test numérique peut détecter une erreur ; il ne démontre pas une égalité pour tous les nombres.\end{attention}
\clearpage
\essSection{2}{Identités remarquables et factorisation}
\begin{propriete}[Identités remarquables]Pour tous nombres $a$ et $b$ :
\[(a+b)^2=a^2+2ab+b^2,\qquad (a-b)^2=a^2-2ab+b^2,\qquad (a-b)(a+b)=a^2-b^2.\]
Lues de gauche à droite, elles développent ; lues de droite à gauche, elles factorisent.\end{propriete}
\begin{exemple}
\[4x^2-25=(2x-5)(2x+5),\qquad x^2+10x+25=(x+5)^2,\qquad 99^2=(100-1)^2=10000-200+1=9801.\]
\end{exemple}
\begin{savoirfaire}[Factoriser]Chercher d’abord un facteur commun apparent : $6x^2-9x=3x(2x-3)$ et $3x(x-2)+5(x-2)=(x-2)(3x+5)$. Sinon, reconnaître une identité remarquable en identifiant $a$ et $b$.\end{savoirfaire}
\begin{savoirfaire}[Prouver avec un programme de calcul]Nommer le nombre choisi $x$, traduire chaque étape, puis transformer l’expression. Pour « ajouter 4, multiplier par 3, retirer le triple du nombre initial », on obtient $3(x+4)-3x=12$ : le résultat est indépendant de $x$.\end{savoirfaire}
\begin{attention}$(a+b)^2$ n’est généralement pas $a^2+b^2$. Il faut les deux produits croisés, dont la somme vaut $2ab$.\end{attention}
\clearpage\rubrique{Découverte}
\begin{blocLivre}{ex-4-1}\begin{exercice}[1]{Remplacer une lettre}
Calculer $A=2x^2-3x+1$ pour $x=-2$ et pour $x=3$.
\end{exercice}\end{blocLivre}
\begin{blocLivre}{ex-4-2}\begin{exercice}[1]{Somme ou produit ?}
Donner la nature de chaque expression : $A=3x+2$ ; $B=5(x+4)$ ; $C=(x-1)(x+3)$ ; $D=x^2-9$ ; $E=2x(x+1)+3$.
\end{exercice}\end{blocLivre}
\begin{blocLivre}{ex-4-3}\begin{exercice}[1]{Prendre l’opposé}
Supprimer les parenthèses et réduire $-(4x-7)+2(3-x)$.
\end{exercice}\end{blocLivre}
\begin{blocLivre}{ex-4-4}\begin{exercice}[1]{Double distributivité}
Développer et réduire $(3x+2)(x-5)$.
\end{exercice}\end{blocLivre}
\begin{blocLivre}{ex-4-5}\begin{exercice}[1]{Un facteur commun}
Factoriser $14x+21$ et $12x^2-8x$.
\end{exercice}\end{blocLivre}

\clearpage\rubrique{Entraînement}
\begin{blocLivre}{ex-4-6}\begin{exercice}[2]{Simplifier des produits et des rapports}
Simplifier, pour $a$, $b$, $c$ et $R$ non nuls : $(3a^2b)\times(2bc)$ ; $\dfrac{6ab^2}{9a^2b}$ ; $\dfrac{\frac43\pi R^3}{\pi R^2}$.
\end{exercice}\end{blocLivre}
\begin{blocLivre}{ex-4-7}\begin{exercice}[2]{Différences de carrés}
Factoriser $x^2-81$ et $25x^2-4$.
\end{exercice}\end{blocLivre}
\begin{blocLivre}{ex-4-8}\begin{exercice}[2]{Deux carrés de sommes}
Développer $(x+6)^2$ et $(2x-3)^2$.
\end{exercice}\end{blocLivre}
\begin{blocLivre}{ex-4-9}\begin{exercice}[2]{Reconnaître un carré}
Factoriser $x^2+10x+25$, puis $9x^2-12x+4$.
\end{exercice}\end{blocLivre}
\begin{blocLivre}{ex-4-10}\begin{exercice}[2]{Factoriser en deux temps}
Factoriser $3x(x-2)+5(x-2)$, puis $9x^2-36$.
\end{exercice}\end{blocLivre}

\clearpage\rubrique{Problèmes}
\begin{blocLivre}{ex-4-11}\begin{exercice}[3]{Deux programmes équivalents}
Programme A : ajouter 5 au nombre choisi, multiplier par 4, retirer 8. Programme B : multiplier le nombre choisi par 4, ajouter 12. Montrer qu’ils sont équivalents.
\end{exercice}\end{blocLivre}
\begin{blocLivre}{ex-4-12}\begin{exercice}[3]{Une bordure}
Un carré de côté $x$ cm est entouré d’une bordure de largeur 2 cm à l’extérieur. Exprimer l’aire de la bordure. Calculer pour $x=7$.
\label{fig-ex-4-12}\figureGeo{bordure}
\end{exercice}\end{blocLivre}
\begin{blocLivre}{ex-4-13}\begin{exercice}[3]{Des entiers et leur parité}
Montrer que le carré d’un entier impair est impair. Le carré d’un entier pair est-il toujours multiple de 4 ?
\end{exercice}\end{blocLivre}
\begin{blocLivre}{ex-4-14}\begin{exercice}[3]{Une égalité impossible ?}
Deux carrés ont pour côtés positifs $x+3$ et $x$. Leur différence d’aires peut-elle valoir 9 cm$^2$ ? Peut-elle valoir 39 cm$^2$ ?
\end{exercice}\end{blocLivre}
\begin{blocLivre}{ex-4-15}\begin{exercice}[3]{Une différence toujours positive}
Comparer, pour tout nombre $x$, $A=(x+1)^2$ et $B=x(x+2)$. Démontrer la réponse.
\end{exercice}\end{blocLivre}

\clearpage

\begin{ouverture}[courbe]{5}{Équations, inéquations et produits nuls}{Partie 1 · Nombres et calculs}
{Équations équivalentes · Inéquations · Produits nuls · Analyse-synthèse}
\objectif{Résoudre une équation ou une inéquation du premier degré.}
\objectif{Factoriser pour résoudre une équation produit nul.}
\objectif{Mettre un problème en équation et valider sa solution.}
\prerequis{Calcul littéral, identités remarquables et opérations inverses.}
\end{ouverture}
\pagestyle{otMain}\essentiel\fontsize{11}{15}\selectfont
\essSection{1}{Transformer une équation}
\begin{definition}Résoudre une équation, c’est trouver toutes les valeurs de l’inconnue qui rendent l’égalité vraie. Deux équations sont équivalentes lorsqu’elles ont les mêmes solutions.\end{definition}
\begin{propriete}Ajouter ou soustraire une même quantité aux deux membres donne une équation équivalente. Multiplier ou diviser les deux membres par un même nombre non nul aussi.\end{propriete}
\begin{exemple}$5x-7=2x+8$ équivaut à $3x=15$, donc à $x=5$. Contrôle : les deux membres initiaux valent 18.\end{exemple}
\begin{savoirfaire}[Analyse et synthèse]Analyse : on suppose qu’un nombre est solution et on en déduit les seules valeurs possibles. Synthèse : on vérifie que ces valeurs conviennent, dans l’équation et dans le contexte. Cette vérification est indispensable lorsque toutes les étapes ne sont pas des équivalences.\end{savoirfaire}
\begin{attention}Une équation peut avoir une seule solution, aucune solution, ou tous les nombres pour solutions. $2x+3=2x+5$ est impossible ; $2(x+3)=2x+6$ est toujours vraie. Diviser par $x$ peut perdre la solution $x=0$.\end{attention}
\clearpage
\essSection{2}{Inéquations et produits nuls}
\begin{propriete}Ajouter un même nombre aux deux membres d’une inégalité conserve son sens. Multiplier ou diviser par un nombre strictement positif le conserve aussi ; par un nombre strictement négatif, il change le sens.\end{propriete}
\begin{exemple}$-2x\ge8$ équivaut à $x\le-4$, car on divise par $-2$. Sur une droite graduée, les solutions forment la demi-droite à gauche de $-4$, le crochet tourné vers les solutions indiquant que $-4$ est inclus. On peut noter cet ensemble $]-\infty\,;-4]$.\end{exemple}
\begin{savoirfaire}[Résoudre graphiquement $ax\ge b$]Tracer la droite d’équation $y=ax$ et la droite horizontale $y=b$. Les solutions sont les abscisses des points de la première droite situés sur ou au-dessus de la seconde.\end{savoirfaire}
\begin{propriete}Un produit est nul si et seulement si l’un au moins de ses facteurs est nul. Ainsi $(x-4)(2x+3)=0$ donne $x=4$ ou $x=-3/2$.\end{propriete}
\begin{savoirfaire}[Se ramener à un produit nul]Tout passer dans un membre, puis factoriser : $x^2=5x$ équivaut à $x(x-5)=0$ ; $(x+1)^2-9=0$ équivaut à $(x-2)(x+4)=0$ grâce à $a^2-b^2$.\end{savoirfaire}
\clearpage\rubrique{Découverte}
\begin{blocLivre}{ex-5-1}\begin{exercice}[1]{Une opération inverse}
Résoudre $x+13=4$ et $-4x=18$.
\end{exercice}\end{blocLivre}
\begin{blocLivre}{ex-5-2}\begin{exercice}[1]{Deux étapes}
Résoudre $7x-5=30$.
\end{exercice}\end{blocLivre}
\begin{blocLivre}{ex-5-3}\begin{exercice}[1]{L’inconnue des deux côtés}
Résoudre $8x+3=5x-12$.
\end{exercice}\end{blocLivre}
\begin{blocLivre}{ex-5-4}\begin{exercice}[1]{Avec une parenthèse}
Résoudre $3(2x-5)=4x+9$.
\end{exercice}\end{blocLivre}
\begin{blocLivre}{ex-5-5}\begin{exercice}[1]{Une première inéquation}
Résoudre $3x\ge12$, puis $-2x\ge8$. Représenter chaque ensemble de solutions sur une droite graduée.
\label{fig-ex-5-5}\figureGeo{inequation-droites}
\end{exercice}\end{blocLivre}

\clearpage\rubrique{Entraînement}
\begin{blocLivre}{ex-5-6}\begin{exercice}[2]{Une équation fractionnaire}
Résoudre $\dfrac{x}{3}+\dfrac{x}{4}=14$.
\end{exercice}\end{blocLivre}
\begin{blocLivre}{ex-5-7}\begin{exercice}[2]{Tout ou rien}
Résoudre $4(x-2)=4x-8$, puis $4(x-2)=4x+1$.
\end{exercice}\end{blocLivre}
\begin{blocLivre}{ex-5-8}\begin{exercice}[2]{Un produit nul}
Résoudre $(2x-7)(x+4)=0$.
\end{exercice}\end{blocLivre}
\begin{blocLivre}{ex-5-9}\begin{exercice}[2]{Factoriser pour résoudre}
Résoudre $x^2-6x=0$, puis $(x+3)^2-16=0$.
\end{exercice}\end{blocLivre}
\begin{blocLivre}{ex-5-10}\begin{exercice}[2]{Lire une inéquation sur un graphique}
La droite représente la fonction $x\mapsto1{,}5x$. Résoudre graphiquement $1{,}5x\ge6$, puis confirmer par le calcul.
\label{fig-ex-5-10}\figureGeo{inequation-graphique}
\end{exercice}\end{blocLivre}

\clearpage\rubrique{Problèmes}
\begin{blocLivre}{ex-5-11}\begin{exercice}[3]{Des billets de spectacle}
Un groupe achète 28 billets : les billets jeunes coûtent 6 €, les autres 10 €. La dépense totale est 232 €. Combien de billets jeunes ?
\end{exercice}\end{blocLivre}
\begin{blocLivre}{ex-5-12}\begin{exercice}[3]{Un rectangle à retrouver}
Un rectangle a un périmètre de 58 cm. Sa longueur dépasse sa largeur de 7 cm. Déterminer ses dimensions et son aire.
\label{fig-ex-5-12}\figureGeo{rectangle-equation}
\end{exercice}\end{blocLivre}
\begin{blocLivre}{ex-5-13}\begin{exercice}[3]{Un achat rentable}
Une gourde filtrante coûte 38 € et ses recharges 12 € par an. Elle remplace un achat de bouteilles estimé à 30 € par an. Écrire une inéquation traduisant « après $n$ années, la dépense avec la gourde est inférieure ou égale à celle des bouteilles », puis la résoudre.
\end{exercice}\end{blocLivre}
\begin{blocLivre}{ex-5-14}\begin{exercice}[3]{Deux aires égales}
Un carré de côté $x+2$ cm et un rectangle de côtés $x$ cm et $x+5$ cm ont la même aire. Les dimensions sont positives. Déterminer $x$ et l’aire commune.
\label{fig-ex-5-14}\figureGeo{aires-egales}
\end{exercice}\end{blocLivre}
\begin{blocLivre}{ex-5-15}\begin{exercice}[3]{Analyse et synthèse}
Un rectangle a une longueur qui dépasse sa largeur $x$ de 2 cm et une aire de 24 cm$^2$. Vérifier que $x(x+2)-24=(x+1)^2-25$, puis déterminer ses dimensions par une analyse suivie d’une synthèse.
\end{exercice}\end{blocLivre}

\clearpage

\begin{ouverture}[courbe]{6}{Proportionnalité, ratios et pourcentages}{Partie 2 · Proportionnalité et fonctions}
{Coefficients · Ratios · Échelles · Évolutions · Lien avec Thalès}
\objectif{Reconnaître et modéliser une proportionnalité.}
\objectif{Partager selon un ratio et utiliser une échelle.}
\objectif{Appliquer une évolution par un coefficient multiplicateur.}
\prerequis{Fractions, produits et repérage dans le plan.}
\end{ouverture}
\pagestyle{otMain}\essentiel\fontsize{11}{15}\selectfont
\essSection{1}{Reconnaître une proportionnalité}
\begin{definition}Deux grandeurs sont proportionnelles si l’on obtient les valeurs de l’une en multipliant celles de l’autre par un même coefficient $a$. La situation est modélisée par la fonction linéaire $x\mapsto ax$ ; sa représentation est une droite passant par l’origine.\end{definition}
\begin{propriete}[Lien avec Thalès]Dans un repère, des points de coordonnées positives sont alignés avec l’origine si et seulement si leurs coordonnées sont proportionnelles. Le théorème de Thalès, appliqué aux triangles formés avec l’axe des abscisses, justifie ce résultat.\end{propriete}
\begin{savoirfaire}[Choisir une méthode]Utiliser le passage à l’unité, un coefficient, la linéarité ou un produit en croix. Si 3 kg coûtent 7,20 €, 1 kg coûte 2,40 €, et 5 kg coûtent 12 €.\end{savoirfaire}
\begin{definition}Partager selon le ratio $a:b$, c’est partager en $a+b$ parts égales et en attribuer $a$ au premier, $b$ au second. Pour un ratio $2:3$ sur 35 L, les volumes sont 14 L et 21 L. Un ratio peut avoir trois termes.\end{definition}
\begin{definition}Une échelle est le rapport longueur sur le plan sur longueur réelle, dans la même unité. À l’échelle $1/500$, 1 cm sur le plan représente 500 cm, soit 5 m.\end{definition}
\clearpage
\essSection{2}{Appliquer et comprendre un pourcentage}
\begin{propriete}Calculer $t\,\%$ d’une quantité revient à la multiplier par $t/100$. Augmenter de $t\,\%$ revient à multiplier par le coefficient multiplicateur $1+t/100$ ; diminuer de $t\,\%$, par $1-t/100$.\end{propriete}
\begin{exemple}Après une baisse de 15\,\%, un prix de 80 € devient $80\times0{,}85=68$ €. Retrouver le prix initial exige de diviser par 0,85.\end{exemple}
\begin{savoirfaire}[Évolutions successives]Multiplier les coefficients, sans additionner les taux. Une hausse de 20\,\% suivie d’une baisse de 20\,\% donne $1{,}2\times0{,}8=0{,}96$ : la baisse totale est de 4\,\%.\end{savoirfaire}
\begin{propriete}Le taux d’évolution de $V_i>0$ à $V_f$ vaut $(V_f-V_i)/V_i$. Préciser la valeur initiale choisie comme référence.\end{propriete}
\begin{attention}Passer de 30\,\% à 36\,\% représente 6 points de pourcentage, mais une augmentation relative de $6/30=20\,\%$. Un tarif avec un forfait fixe n’est pas proportionnel au nombre d’unités.\end{attention}
\clearpage\rubrique{Découverte}
\begin{blocLivre}{ex-6-1}\begin{exercice}[1]{Passer par l’unité}
Quatre cahiers coûtent 10,80 €. Le prix est proportionnel au nombre de cahiers. Quel prix pour sept cahiers ?
\end{exercice}\end{blocLivre}
\begin{blocLivre}{ex-6-2}\begin{exercice}[1]{Un tableau proportionnel ?}
Les couples quantité-prix sont $(2\,;7)$, $(5\,;17{,}5)$ et $(8\,;28)$. Le tableau est-il proportionnel ? Donner la fonction linéaire qui le modélise.
\end{exercice}\end{blocLivre}
\begin{blocLivre}{ex-6-3}\begin{exercice}[1]{Appliquer une réduction}
Un sac coûte 64 €. Son prix baisse de 25 \%. Calculer le nouveau prix.
\end{exercice}\end{blocLivre}
\begin{blocLivre}{ex-6-4}\begin{exercice}[1]{Calculer un taux}
Un effectif passe de 160 à 184. Quel est son taux d’augmentation ?
\end{exercice}\end{blocLivre}
\begin{blocLivre}{ex-6-5}\begin{exercice}[1]{Partager selon un ratio}
Partager 56 billes selon le ratio $3:5$.
\end{exercice}\end{blocLivre}

\clearpage\rubrique{Entraînement}
\begin{blocLivre}{ex-6-6}\begin{exercice}[2]{Retrouver un prix initial}
Après une hausse de 8 \%, un abonnement coûte 81 €. Quel était son prix ?
\end{exercice}\end{blocLivre}
\begin{blocLivre}{ex-6-7}\begin{exercice}[2]{Deux baisses}
Un prix de 250 € baisse de 10 \%, puis de 20 \%. Calculer le prix final et la réduction globale.
\end{exercice}\end{blocLivre}
\begin{blocLivre}{ex-6-8}\begin{exercice}[2]{Une échelle}
Deux lieux distants de 7,2 km sont séparés de 9 cm sur une carte. Déterminer l’échelle. Quelle distance réelle représentent 4 cm sur cette carte ?
\end{exercice}\end{blocLivre}
\begin{blocLivre}{ex-6-9}\begin{exercice}[2]{Une masse en trois parts}
Un mélange de 4,5 kg est préparé selon le ratio farine : sucre : beurre $=5:3:1$. Calculer chaque masse.
\end{exercice}\end{blocLivre}
\begin{blocLivre}{ex-6-10}\begin{exercice}[2]{Proportionnellement à l’âge}
Deux sœurs de 12 ans et 18 ans se partagent 450 € proportionnellement à leur âge. Combien reçoit chacune ?
\end{exercice}\end{blocLivre}

\clearpage\rubrique{Problèmes}
\begin{blocLivre}{ex-6-11}\begin{exercice}[3]{Comparer deux promotions}
Un article coûte 120 €. Offre A : réduction de 30 €. Offre B : réduction de 20 \%, puis bon d’achat de 10 € utilisable immédiatement sur cet achat. Comparer les paiements.
\end{exercice}\end{blocLivre}
\begin{blocLivre}{ex-6-12}\begin{exercice}[3]{Un mélange à corriger}
On mélange 2 L de sirop avec 8 L d’eau. On veut que le sirop représente 16 \% du mélange sans en retirer. Quel volume d’eau ajouter ?
\end{exercice}\end{blocLivre}
\begin{blocLivre}{ex-6-13}\begin{exercice}[3]{Deux pourcentages emboîtés}
Dans un collège de 720 élèves, 45 \% sont inscrits à l’association sportive. Parmi les inscrits, $2/9$ pratiquent le badminton. Quel est leur effectif et leur part dans le collège ?
\end{exercice}\end{blocLivre}
\begin{blocLivre}{ex-6-14}\begin{exercice}[3]{Lire une annonce}
La part des élèves venant à vélo passe de 12 \% à 15 \%. Une annonce parle de « 3 \% de hausse ». Expliquer précisément l’évolution. Sur 800 élèves à effectif constant, combien de cyclistes supplémentaires ?
\end{exercice}\end{blocLivre}
\begin{blocLivre}{ex-6-15}\begin{exercice}[3]{Proportionnalité et alignement}
Dans un repère, $O$ est l’origine, $A(4\,;3)$ et $B(10\,;7{,}5)$ sont alignés avec $O$. On note $H(4\,;0)$ et $K(10\,;0)$. Justifier que $(AH)$ et $(BK)$ sont parallèles, puis utiliser le théorème de Thalès pour montrer que les coordonnées de $A$ et de $B$ sont proportionnelles.
\label{fig-ex-6-15}\figureGeo{thales-repere}
\end{exercice}\end{blocLivre}

\clearpage

\begin{ouverture}[courbe]{7}{Grandeurs quotients et conversions}{Partie 2 · Proportionnalité et fonctions}
{Durées · Vitesses · Débits · Masses volumiques · Unités}
\objectif{Convertir les unités adaptées à une situation.}
\objectif{Calculer avec une vitesse, un débit ou une masse volumique.}
\objectif{Contrôler les unités et l’ordre de grandeur d’un résultat.}
\prerequis{Grandeurs quotients rencontrées en Quatrième, proportionnalité et unités.}
\end{ouverture}
\pagestyle{otMain}\essentiel\fontsize{11}{15}\selectfont
\essSection{1}{Choisir et convertir les unités}
\begin{propriete}Pour les longueurs, $1\text{ m}=100\text{ cm}$. Pour les aires, $1\text{ m}^2=10000\text{ cm}^2$ ; pour les volumes, $1\text{ m}^3=10^6\text{ cm}^3$. Le facteur de conversion est mis au carré ou au cube.\end{propriete}
\begin{propriete}$1\text{ L}=1\text{ dm}^3$, $1\text{ mL}=1\text{ cm}^3$ et $1\text{ m}^3=1000\text{ L}$. Une tonne vaut 1000 kg. Une heure vaut 60 min, soit 3600 s.\end{propriete}
\begin{exemple}$2{,}4\text{ h}=2\text{ h}+0{,}4\times60\text{ min}=2\text{ h }24\text{ min}$. À l’inverse, 2 h 15 min valent $2+15/60=2{,}25$ h.\end{exemple}
\begin{attention}1 h 30 min n’est pas 1,30 h. De même, multiplier par 100 pour convertir une aire en m$^2$ vers des cm$^2$ est incorrect : le facteur est 10000.\end{attention}
\begin{savoirfaire}[Contrôler la cohérence]Nommer la grandeur et son unité à chaque étape. Une aire ne s’exprime pas en mètres, une durée ne s’exprime pas en kilomètres. Donner un ordre de grandeur avant de conclure.\end{savoirfaire}
\clearpage
\essSection{2}{Des grandeurs quotients}
\begin{definition}Une grandeur quotient est le quotient de deux grandeurs. Une vitesse moyenne vaut $v=d/t$, avec $t>0$. Un débit volumique est un volume par unité de temps : $q=V/t$. Une masse volumique est une masse par unité de volume : $\rho=m/V$.\end{definition}
\begin{exemple}$72\text{ km/h}=72000/3600=20\text{ m/s}$. Pour convertir des km/h en m/s, diviser par 3,6 ; pour l’opération inverse, multiplier par 3,6.\end{exemple}
\begin{savoirfaire}[Une chaîne de calculs]On peut utiliser $d=vt$ et $t=d/v$ pour $v>0$. Pour remplir 600 L à un débit de 15 L/min, la durée est $600/15=40$ min. À débit constant, doubler le volume double la durée : la durée est proportionnelle au volume.\end{savoirfaire}
\begin{attention}La moyenne de deux vitesses n’est pas généralement la vitesse moyenne d’un trajet. Les durées passées à chaque vitesse comptent : diviser la distance totale par la durée totale.\end{attention}
Ces grandeurs, étudiées en Quatrième, interviennent souvent dans les problèmes du brevet et dans les autres disciplines.
\clearpage\rubrique{Découverte}
\begin{blocLivre}{ex-7-1}\begin{exercice}[1]{Des longueurs et des aires}
Convertir 3,8 km en mètres, 0,42 m$^2$ en cm$^2$ et 6500 cm$^2$ en m$^2$.
\end{exercice}\end{blocLivre}
\begin{blocLivre}{ex-7-2}\begin{exercice}[1]{Volumes et contenances}
Convertir 2,7 m$^3$ en litres et 450 mL en cm$^3$ puis en litres.
\end{exercice}\end{blocLivre}
\begin{blocLivre}{ex-7-3}\begin{exercice}[1]{Une durée décimale}
Écrire 1,75 h en heures et minutes ; écrire 2 h 42 min en heures décimales.
\end{exercice}\end{blocLivre}
\begin{blocLivre}{ex-7-4}\begin{exercice}[1]{Une vitesse moyenne}
Un trajet de 150 km est parcouru en 2 h 30 min. Calculer la vitesse moyenne.
\end{exercice}\end{blocLivre}
\begin{blocLivre}{ex-7-5}\begin{exercice}[1]{Un débit}
Un robinet fournit 84 L en 7 minutes, à débit constant. Donner son débit et le volume en 25 minutes.
\end{exercice}\end{blocLivre}

\clearpage\rubrique{Entraînement}
\begin{blocLivre}{ex-7-6}\begin{exercice}[2]{Convertir une vitesse}
Convertir 90 km/h en m/s et 12 m/s en km/h.
\end{exercice}\end{blocLivre}
\begin{blocLivre}{ex-7-7}\begin{exercice}[2]{Le temps de réaction}
Une voiture roule à 54 km/h. Le conducteur met 0,8 s avant de commencer à freiner. Quelle distance parcourt-elle pendant cette réaction ?
\end{exercice}\end{blocLivre}
\begin{blocLivre}{ex-7-8}\begin{exercice}[2]{Une masse volumique}
Un matériau a une masse volumique de 2,5 g/cm$^3$. Un bloc occupe 360 cm$^3$. Calculer sa masse en grammes puis en kilogrammes.
\end{exercice}\end{blocLivre}
\begin{blocLivre}{ex-7-9}\begin{exercice}[2]{Une consommation électrique}
Un appareil de puissance constante 1,2 kW fonctionne pendant 45 minutes. Calculer l’énergie en kWh. Le kWh coûte 0,24 € dans le modèle proposé. Calculer le coût.
\end{exercice}\end{blocLivre}
\begin{blocLivre}{ex-7-10}\begin{exercice}[2]{Un téléchargement}
Un fichier de 2,4 Go est téléchargé à un débit constant de 40 Mo/s. Écrire ces données en octets avec des puissances de 10, puis calculer la durée du téléchargement.
\end{exercice}\end{blocLivre}

\clearpage\rubrique{Problèmes}
\begin{blocLivre}{ex-7-11}\begin{exercice}[3]{Aller et retour}
Un cycliste parcourt 12 km à 18 km/h, puis revient sur 12 km à 12 km/h. Calculer sa vitesse moyenne sur l’aller-retour.
\end{exercice}\end{blocLivre}
\begin{blocLivre}{ex-7-12}\begin{exercice}[3]{Deux robinets}
Un premier robinet remplit seul une cuve en 30 min, un second en 45 min. Ils fonctionnent ensemble avec des débits constants et indépendants. Quelle durée pour remplir la cuve ?
\end{exercice}\end{blocLivre}
\begin{blocLivre}{ex-7-13}\begin{exercice}[3]{Choisir assez de peinture}
Un mur rectangulaire mesure 5 m sur 2,6 m. Une ouverture de 2 m$^2$ n’est pas peinte. Deux couches sont prévues. Un litre couvre 8 m$^2$ par couche. Les pots contiennent 2,5 L. Combien de pots acheter ?
\end{exercice}\end{blocLivre}
\begin{blocLivre}{ex-7-14}\begin{exercice}[3]{Une fuite discrète}
Une fuite laisse tomber 2 gouttes par seconde. On modélise une goutte par 0,05 mL. Quel volume d’eau est perdu en une journée ? En 30 jours ?
\end{exercice}\end{blocLivre}
\begin{blocLivre}{ex-7-15}\begin{exercice}[3]{Une livraison de gravier}
Une allée de 12 m sur 1,5 m reçoit une couche uniforme de gravier de 6 cm. La masse volumique apparente est 1,6 tonne/m$^3$. Un camion peut emporter 1,5 tonne. Combien de voyages au minimum ?
\end{exercice}\end{blocLivre}

\clearpage

\begin{ouverture}[courbe]{8}{Fonctions, images et fonction carré}{Partie 2 · Proportionnalité et fonctions}
{Formules · Tableaux · Courbes · Images · Antécédents · $x\mapsto x^2$}
\objectif{Utiliser le vocabulaire et les notations des fonctions.}
\objectif{Passer d’une formule à un tableau ou à une courbe.}
\objectif{Représenter la fonction carré et résoudre graphiquement $x^2=a$.}
\prerequis{Repérage, substitution dans une expression et racine carrée.}
\end{ouverture}
\pagestyle{otMain}\essentiel\fontsize{11}{15}\selectfont
\essSection{1}{Une entrée et une image}
\begin{definition}Une fonction associe à chaque nombre autorisé une seule image. On écrit $f(x)$ pour l’image de $x$ par $f$, et $f:x\mapsto x^2+1$ pour décrire la fonction. Si $f(3)=8$, 8 est l’image de 3 et 3 est un antécédent de 8. Un nombre peut avoir plusieurs antécédents, ou aucun.\end{definition}
\begin{exemple}Pour $f(x)=x^2+1$, $f(-2)=5$. Les nombres $-2$ et 2 sont deux antécédents de 5. Aucun nombre n’a pour image 0, car $x^2+1\ge1$.\end{exemple}
\begin{savoirfaire}[Trois représentations]Une formule permet de calculer toute image ; un tableau de valeurs n’en donne que quelques-unes ; une courbe rassemble les points de coordonnées $(x\,;f(x))$. Un point appartient à la courbe lorsque son ordonnée est l’image de son abscisse.\end{savoirfaire}
\begin{attention}$f(2)$ ne signifie pas $f\times2$. Ne pas échanger l’image et l’antécédent dans une phrase. Un tableau partiel ne permet pas d’inventer l’image d’une entrée non indiquée.\end{attention}
\begin{propriete}Comparer deux fonctions graphiquement revient à comparer leurs ordonnées pour une même abscisse. Les intersections correspondent aux entrées donnant la même image.\end{propriete}
\clearpage
\essSection{2}{Lire une courbe et la fonction carré}
\begin{savoirfaire}[Lire une image ou des antécédents]Partir d’une abscisse, rejoindre la courbe verticalement, puis lire l’ordonnée. Pour chercher les antécédents d’un nombre, partir de cette ordonnée et repérer toutes les intersections horizontales avec la courbe. Une lecture graphique est souvent approchée.\end{savoirfaire}
\begin{definition}La fonction carré est la fonction $x\mapsto x^2$. Sa courbe, symétrique par rapport à l’axe des ordonnées, est une parabole qui passe par l’origine et reste au-dessus de l’axe des abscisses.\end{definition}
\begin{propriete}Résoudre graphiquement $x^2=a$, c’est chercher les abscisses des points d’intersection de la parabole avec la droite horizontale d’équation $y=a$ : deux points si $a>0$, un seul si $a=0$, aucun si $a<0$.\end{propriete}
\begin{savoirfaire}[Modéliser]Nommer la variable et son unité, exprimer la quantité cherchée, puis préciser les entrées possibles. Un nombre de billets est entier ; une longueur doit rester positive. Pour une cuve de base 12 m$^2$, le volume est $V(h)=12h$, pour $h$ compris entre 0 et la hauteur maximale.\end{savoirfaire}
\clearpage\rubrique{Découverte}
\begin{blocLivre}{ex-8-1}\begin{exercice}[1]{Calculer des images}
Pour $f:x\mapsto3x-4$, calculer $f(0)$, $f(2)$ et $f(-3)$.
\end{exercice}\end{blocLivre}
\begin{blocLivre}{ex-8-2}\begin{exercice}[1]{Lire un tableau}
À l’aide du tableau, donner $f(-2)$ et les antécédents de 1 qui y figurent.\begin{center}\renewcommand{\arraystretch}{1.35}\begin{tabular}{c|cccc}\toprule
$x$ & $-2$ & $0$ & $1$ & $3$\\\midrule
$f(x)$ & $5$ & $1$ & $-1$ & $1$\\
\bottomrule\end{tabular}\end{center}
\end{exercice}\end{blocLivre}
\begin{blocLivre}{ex-8-3}\begin{exercice}[1]{Traduire une phrase}
Écrire avec une égalité : « 7 est l’image de $-1$ par $g$ ». Dire ensuite ce que signifie $g(4)=7$. Ces deux informations sont-elles compatibles ?
\end{exercice}\end{blocLivre}
\begin{blocLivre}{ex-8-4}\begin{exercice}[1]{Un point sur une courbe}
Pour $f(x)=x^2-2$, les points $A(3\,;7)$ et $B(-2\,;1)$ appartiennent-ils à la courbe ?
\end{exercice}\end{blocLivre}
\begin{blocLivre}{ex-8-5}\begin{exercice}[1]{Un antécédent algébrique}
Pour $h(x)=5x+2$, déterminer l’antécédent de 17.
\end{exercice}\end{blocLivre}

\clearpage\rubrique{Entraînement}
\begin{blocLivre}{ex-8-6}\begin{exercice}[2]{Une formule et un programme}
On choisit un nombre, on ajoute 3, puis on élève le résultat au carré. Donner la fonction obtenue et les images de 1 et $-3$.
\end{exercice}\end{blocLivre}
\begin{blocLivre}{ex-8-7}\begin{exercice}[2]{Lire une ligne brisée}
La courbe est formée des segments dessinés, avec les sommets sur les graduations indiquées. Lire $f(-1)$ et tous les antécédents de 1.
\label{fig-ex-8-7}\figureGeo{fonction-brisee}
\end{exercice}\end{blocLivre}
\begin{blocLivre}{ex-8-8}\begin{exercice}[2]{Plusieurs ou aucun antécédent}
Pour $f(x)=x^2$, déterminer les antécédents de 16, de 0 et de $-1$.
\label{fig-ex-8-8}\figureGeo{fonction-carre}
\end{exercice}\end{blocLivre}
\begin{blocLivre}{ex-8-9}\begin{exercice}[2]{Comparer deux représentations}
Les fonctions $f(x)=x+2$ et $g(x)=2x-1$ sont représentées. Pour quelle valeur de $x$ leurs images sont-elles égales ? Vérifier par un calcul.
\label{fig-ex-8-9}\figureGeo{fonction-intersection}
\end{exercice}\end{blocLivre}
\begin{blocLivre}{ex-8-10}\begin{exercice}[2]{Représenter la fonction carré}
Compléter un tableau de valeurs de $x\mapsto x^2$ pour $x$ égal à $-3$, $-2$, $-1$, $-0{,}5$, 0, 0,5, 1, 2 et 3, puis tracer la courbe dans le repère.
\label{fig-ex-8-10}\figureGeo{carre-trace}
\end{exercice}\end{blocLivre}

\clearpage\rubrique{Problèmes}
\begin{blocLivre}{ex-8-11}\begin{exercice}[3]{Un réservoir}
Un réservoir cylindrique de base d’aire 0,8 m$^2$ a une hauteur maximale de 1,5 m. Exprimer son volume en fonction de la hauteur $h$, puis trouver la hauteur correspondant à 600 L.
\label{fig-ex-8-11}\figureGeo{cuve-hauteur}
\end{exercice}\end{blocLivre}
\begin{blocLivre}{ex-8-12}\begin{exercice}[3]{Résoudre graphiquement $x^2=a$}
À l’aide de la courbe de la fonction carré, lire des valeurs approchées des solutions de $x^2=5$, puis donner leurs valeurs exactes et leurs arrondis au dixième. Combien de solutions ont $x^2=0$ et $x^2=-2$ ?
\label{fig-ex-8-12}\figureGeo{carre-equation}
\end{exercice}\end{blocLivre}
\begin{blocLivre}{ex-8-13}\begin{exercice}[3]{Des aires possibles}
Un rectangle de périmètre 40 cm a un côté de longueur $x$. Exprimer son aire. Quelles sont les entrées possibles ? Donner les images de 3 et de 17.
\label{fig-ex-8-13}\figureGeo{rectangle-fonction}
\end{exercice}\end{blocLivre}
\begin{blocLivre}{ex-8-14}\begin{exercice}[3]{Une lecture de déplacement}
La courbe donne la distance au point de départ pendant une promenade. Les segments ont pour sommets $(0\,;0)$, $(10\,;600)$, $(15\,;600)$ et $(25\,;0)$, en minutes et mètres. Décrire les phases et calculer la distance totale parcourue sur cet aller-retour rectiligne.
\label{fig-ex-8-14}\figureGeo{promenade}
\end{exercice}\end{blocLivre}
\begin{blocLivre}{ex-8-15}\begin{exercice}[3]{Résoudre par un produit nul}
Pour $f(x)=(x-2)(3x+6)$, déterminer les antécédents de 0. Donner ensuite une image strictement négative.
\end{exercice}\end{blocLivre}

\clearpage

\begin{ouverture}[courbe]{9}{Fonctions linéaires et affines}{Partie 2 · Proportionnalité et fonctions}
{Proportionnalité · Droites · Coefficients · Inéquations · Tarifs}
\objectif{Reconnaître une fonction linéaire ou affine.}
\objectif{Tracer une droite et déterminer graphiquement ses coefficients.}
\objectif{Comparer des modèles et résoudre des inéquations par le calcul et le graphique.}
\prerequis{Fonctions, équations et inéquations du premier degré.}
\end{ouverture}
\pagestyle{otMain}\essentiel\fontsize{11}{15}\selectfont
\essSection{1}{Deux familles de fonctions}
\begin{definition}Une fonction linéaire est une fonction de la forme $f:x\mapsto ax$. Elle modélise une situation de proportionnalité, de coefficient $a$, et sa représentation est une droite passant par l’origine.\end{definition}
\begin{definition}Une fonction affine est une fonction de la forme $g:x\mapsto ax+b$. Sa représentation est une droite ; $b=g(0)$ est son ordonnée à l’origine. Le coefficient $a$, appelé coefficient directeur ou pente, indique de combien varie l’ordonnée lorsque l’abscisse augmente de 1.\end{definition}
\begin{exemple}Pour $g(x)=2x-3$, le point $(0\,;-3)$ est sur la droite. Quand l’abscisse augmente de 1, l’ordonnée augmente de 2. Les points $(0\,;-3)$ et $(2\,;1)$ suffisent pour tracer la droite.\end{exemple}
\begin{propriete}Un coefficient $a$ positif correspond à une droite montante de gauche à droite, un coefficient négatif à une droite descendante. Si $a=0$, la fonction est constante.\end{propriete}
\begin{attention}« Affine » n’implique pas « proportionnel ». Il faut en plus $b=0$ pour obtenir une fonction linéaire.\end{attention}
\clearpage
\essSection{2}{Retrouver une formule et comparer}
\begin{savoirfaire}[Coefficients sur un graphique]Lire $b$ à l’intersection avec l’axe des ordonnées. Pour $a$, partir d’un point de la droite, avancer d’une unité vers la droite et lire la variation verticale, positive en montant, négative en descendant.\end{savoirfaire}
\begin{savoirfaire}[À partir de deux points]Pour deux points de la droite $(x_1\,;y_1)$ et $(x_2\,;y_2)$ avec $x_1\ne x_2$ :
\[a=\frac{y_2-y_1}{x_2-x_1}.\]
Utiliser ensuite $y_1=ax_1+b$ pour retrouver $b$. Une fonction linéaire est déterminée par un seul nombre non nul et son image.\end{savoirfaire}
\begin{savoirfaire}[Inéquations et seuils]Pour comparer deux tarifs, résoudre leur égalité pour repérer le seuil, puis l’inéquation correspondante, par le calcul ou en comparant les positions des droites. Revenir au nombre entier d’unités lorsque la situation l’exige.\end{savoirfaire}
\begin{attention}Une droite verticale ne représente pas une fonction : la même entrée aurait plusieurs images. Une formule affine doit être justifiée par le modèle ou les données.\end{attention}
\clearpage\rubrique{Découverte}
\begin{blocLivre}{ex-9-1}\begin{exercice}[1]{Reconnaître les formes}
Parmi $f(x)=-3x$, $g(x)=4x+1$, $h(x)=x^2$ et $k(x)=7$, préciser les fonctions affines et les fonctions linéaires.
\end{exercice}\end{blocLivre}
\begin{blocLivre}{ex-9-2}\begin{exercice}[1]{Calculer avec une linéaire}
Pour $f(x)=-2{,}5x$, calculer $f(4)$ et l’antécédent de 15.
\end{exercice}\end{blocLivre}
\begin{blocLivre}{ex-9-3}\begin{exercice}[1]{Lire les coefficients}
Pour $g(x)=-3x+8$, donner le coefficient directeur, l’ordonnée à l’origine et $g(2)$.
\end{exercice}\end{blocLivre}
\begin{blocLivre}{ex-9-4}\begin{exercice}[1]{Tracer une linéaire}
Tracer la représentation de $f(x)=1{,}5x$ à l’aide de deux points.
\label{fig-ex-9-4}\figureGeo{affine-trace-lineaire}
\end{exercice}\end{blocLivre}
\begin{blocLivre}{ex-9-5}\begin{exercice}[1]{Tracer une affine}
Tracer $g(x)=-x+4$ à l’aide de deux points, puis lire son antécédent de 0.
\label{fig-ex-9-5}\figureGeo{affine-trace}
\end{exercice}\end{blocLivre}

\clearpage\rubrique{Entraînement}
\begin{blocLivre}{ex-9-6}\begin{exercice}[2]{Une linéaire à retrouver}
Une fonction linéaire $f$ vérifie $f(6)=21$. Donner son expression et $f(-2)$.
\end{exercice}\end{blocLivre}
\begin{blocLivre}{ex-9-7}\begin{exercice}[2]{Deux points, une formule}
Une fonction affine passe par $A(2\,;1)$ et $B(5\,;10)$. Déterminer son expression.
\label{fig-ex-9-7}\figureGeo{affine-deux-points}
\end{exercice}\end{blocLivre}
\begin{blocLivre}{ex-9-8}\begin{exercice}[2]{Retrouver un tarif}
Un service affine facture 14 € pour 2 heures et 26 € pour 5 heures. Déterminer le forfait fixe et le prix par heure.
\end{exercice}\end{blocLivre}
\begin{blocLivre}{ex-9-9}\begin{exercice}[2]{Coefficients lus sur un graphique}
Déterminer graphiquement les coefficients $a$ et $b$ de la fonction affine représentée, puis donner son expression.
\label{fig-ex-9-9}\figureGeo{affine-lecture}
\end{exercice}\end{blocLivre}
\begin{blocLivre}{ex-9-10}\begin{exercice}[2]{Comparer par une différence}
On donne $f(x)=4x-1$ et $g(x)=4x+6$. Comparer leurs images pour toute entrée.
\label{fig-ex-9-10}\figureGeo{affine-paralleles}
\end{exercice}\end{blocLivre}

\clearpage\rubrique{Problèmes}
\begin{blocLivre}{ex-9-11}\begin{exercice}[3]{Deux abonnements}
Pour $n$ séances, l’offre A coûte $20+3n$ euros et l’offre B $8+5n$ euros. Déterminer le nombre de séances donnant le même coût, puis résoudre l’inéquation $20+3n<8+5n$ et conclure.
\label{fig-ex-9-11}\figureGeo{affine-tarifs}
\end{exercice}\end{blocLivre}
\begin{blocLivre}{ex-9-12}\begin{exercice}[3]{Une vidange à débit constant}
Une cuve contient 900 L. Elle perd 35 L par minute. Donner le volume restant $V(t)$ avant qu’elle soit vide. Quand reste-t-il 200 L ? Après 30 minutes, la formule peut-elle encore représenter un volume ?
\end{exercice}\end{blocLivre}
\begin{blocLivre}{ex-9-13}\begin{exercice}[3]{Une vitesse et une inéquation}
Un mobile se déplace à vitesse constante de 5 m/s. On note $d(t)$ la distance parcourue, en mètres, après $t$ secondes. Exprimer $d$, puis déterminer à partir de quel instant la distance dépasse strictement 120 m. Retrouver le résultat graphiquement.
\end{exercice}\end{blocLivre}
\begin{blocLivre}{ex-9-14}\begin{exercice}[3]{Un point suffit-il ?}
On sait seulement qu’une fonction affine vérifie $f(2)=7$. Peut-on déterminer sa formule ? Donner deux formules différentes possibles.
\end{exercice}\end{blocLivre}
\begin{blocLivre}{ex-9-15}\begin{exercice}[3]{Un seuil avec contrainte}
Deux locations pour $d$ jours entiers coûtent $A(d)=45+12d$ et $B(d)=15+18d$. Un budget de 120 € est prévu. Déterminer la meilleure offre pour 6 jours et le nombre maximal de jours possible dans chaque offre.
\end{exercice}\end{blocLivre}

\clearpage

\begin{ouverture}[courbe]{10}{Statistiques et boîtes à moustaches}{Partie 3 · Données et probabilités}
{Effectifs cumulés · Médiane · Quartiles · Boîtes à moustaches · Tableur}
\objectif{Calculer et interpréter moyenne, médiane, quartiles et étendue.}
\objectif{Utiliser les effectifs cumulés croissants.}
\objectif{Construire et comparer des boîtes à moustaches.}
\prerequis{Moyenne, médiane, étendue et fréquences.}
\end{ouverture}
\pagestyle{otMain}\essentiel\fontsize{11}{15}\selectfont
\essSection{1}{Effectifs cumulés, médiane et quartiles}
\begin{definition}L’effectif cumulé croissant d’une valeur est le nombre d’observations inférieures ou égales à cette valeur. Le dernier effectif cumulé est l’effectif total $N$. La moyenne est la somme des valeurs, chacune comptée avec son effectif, divisée par $N$.\end{definition}
\begin{definition}Une médiane partage la série ordonnée en deux groupes de même effectif. Si $N$ est impair, c’est la valeur de rang $(N+1)/2$ ; si $N$ est pair, ce livre retient la moyenne des valeurs de rangs $N/2$ et $N/2+1$.\end{definition}
\begin{definition}Le premier quartile $Q_1$ est la plus petite valeur de la série telle qu’au moins 25\,\% des valeurs lui soient inférieures ou égales. Le troisième quartile $Q_3$ est la plus petite valeur telle qu’au moins 75\,\% des valeurs lui soient inférieures ou égales.\end{definition}
\begin{exemple}Pour $N=20$, $Q_1$ est la 5\textsuperscript{e} valeur, car $20\times0{,}25=5$, et $Q_3$ la 15\textsuperscript{e}, car $20\times0{,}75=15$. Pour $N=11$, $11\times0{,}25=2{,}75$ : $Q_1$ est la 3\textsuperscript{e} valeur. On arrondit le rang à l’entier supérieur.\end{exemple}
\begin{savoirfaire}[Lire les rangs]Ranger les valeurs ou calculer les effectifs cumulés croissants ; repérer la première valeur dont l’effectif cumulé atteint le rang cherché.\end{savoirfaire}
\clearpage
\essSection{2}{Boîtes à moustaches et comparaisons}
\begin{definition}Une boîte à moustaches représente, sur un axe gradué, le minimum, $Q_1$, la médiane, $Q_3$ et le maximum. La boîte va de $Q_1$ à $Q_3$ et contient environ la moitié centrale des valeurs ; les moustaches rejoignent les extrêmes.\end{definition}
\begin{savoirfaire}[Comparer deux séries]Comparer les médianes pour la position, puis l’étendue et la longueur $Q_3-Q_1$ de la boîte pour la dispersion. Une phrase précise s’appuie sur un pourcentage : « au moins 75\,\% des valeurs sont inférieures ou égales à $Q_3$ ».\end{savoirfaire}
\begin{savoirfaire}[Tableur]Les fonctions MOYENNE, MEDIANE, MIN et MAX calculent ces indicateurs sur une plage de cellules, par exemple \code{=MAX(A2:A21)-MIN(A2:A21)} pour l’étendue. La fonction QUARTILE d’un tableur suit souvent une autre définition et peut donner un résultat différent.\end{savoirfaire}
\begin{attention}Une moyenne identique n’impose ni une médiane identique ni une même dispersion. Une valeur extrême peut déplacer fortement la moyenne sans changer la médiane. Lire l’origine des axes avant de comparer des hauteurs.\end{attention}
\clearpage\rubrique{Découverte}
\begin{blocLivre}{ex-10-1}\begin{exercice}[1]{Trois indicateurs}
Pour la série $7\,;9\,;10\,;14\,;15$, calculer la moyenne, la médiane et l’étendue.
\end{exercice}\end{blocLivre}
\begin{blocLivre}{ex-10-2}\begin{exercice}[1]{Une fréquence}
Dans un groupe de 30 élèves, 12 viennent à vélo. Donner la fréquence et le pourcentage correspondants.
\end{exercice}\end{blocLivre}
\begin{blocLivre}{ex-10-3}\begin{exercice}[1]{Des effectifs cumulés}
Calculer les effectifs cumulés croissants de la série de notes, puis l’effectif total.\begin{center}\renewcommand{\arraystretch}{1.35}\begin{tabular}{c|ccccc}\toprule
Note & $6$ & $8$ & $10$ & $12$ & $16$\\\midrule
Effectif & $2$ & $3$ & $6$ & $5$ & $4$\\
\bottomrule\end{tabular}\end{center}
\end{exercice}\end{blocLivre}
\begin{blocLivre}{ex-10-4}\begin{exercice}[1]{Ordonner avant la médiane}
Déterminer la médiane de $13\,;4\,;9\,;7\,;11$, puis celle de $1\,;3\,;6\,;8\,;10\,;20$.
\end{exercice}\end{blocLivre}
\begin{blocLivre}{ex-10-5}\begin{exercice}[1]{Lire une boîte à moustaches}
La boîte à moustaches résume les durées, en minutes, d’un trajet. Lire le minimum, les quartiles, la médiane et le maximum. Calculer l’étendue et compléter : « au moins … \% des durées sont inférieures ou égales à 14 min ».
\label{fig-ex-10-5}\figureGeo{boite-lecture}
\end{exercice}\end{blocLivre}

\clearpage\rubrique{Entraînement}
\begin{blocLivre}{ex-10-6}\begin{exercice}[2]{Moyenne et quartiles d’un tableau}
Pour la série de notes de l’exercice 3, calculer la moyenne, la médiane et les deux quartiles.\begin{center}\renewcommand{\arraystretch}{1.35}\begin{tabular}{c|ccccc}\toprule
Note & $6$ & $8$ & $10$ & $12$ & $16$\\\midrule
Effectif & $2$ & $3$ & $6$ & $5$ & $4$\\
\bottomrule\end{tabular}\end{center}
\end{exercice}\end{blocLivre}
\begin{blocLivre}{ex-10-7}\begin{exercice}[2]{Quartiles d’une série brute}
Déterminer la médiane et les quartiles de la série ordonnée $3\,;5\,;6\,;8\,;9\,;10\,;12\,;13\,;15\,;18\,;20$.
\end{exercice}\end{blocLivre}
\begin{blocLivre}{ex-10-8}\begin{exercice}[2]{Construire une boîte}
Construire la boîte à moustaches de la série de l’exercice précédent sur l’axe gradué fourni.
\label{fig-ex-10-8}\figureGeo{boite-construire}
\end{exercice}\end{blocLivre}
\begin{blocLivre}{ex-10-9}\begin{exercice}[2]{Quartiles et diagramme en barres}
Le diagramme en barres donne le nombre de livres lus pendant les vacances par les élèves d’une classe. Déterminer l’effectif total, la médiane et les quartiles.
\label{fig-ex-10-9}\figureGeo{stat-barres}
\end{exercice}\end{blocLivre}
\begin{blocLivre}{ex-10-10}\begin{exercice}[2]{Des formules de tableur}
Les durées 12, 15, 9, 20, 15, 11, 18 et 14, en minutes, sont saisies dans les cellules A2 à A9. Écrire des formules donnant la moyenne, la médiane et l’étendue, puis donner les résultats.
\end{exercice}\end{blocLivre}

\clearpage\rubrique{Problèmes}
\begin{blocLivre}{ex-10-11}\begin{exercice}[3]{Réunir deux groupes}
Un groupe de 12 élèves a une moyenne de 13 ; un autre de 18 élèves a une moyenne de 11. Quelle est la moyenne des 30 élèves ?
\end{exercice}\end{blocLivre}
\begin{blocLivre}{ex-10-12}\begin{exercice}[3]{Comparer deux collèges}
Les boîtes à moustaches résument les durées de trajet, en minutes, des élèves de deux collèges A et B. Comparer les médianes et les dispersions. Dans quel collège au moins un quart des élèves mettent-ils 22 minutes ou plus ?
\label{fig-ex-10-12}\figureGeo{boites-comparer}
\end{exercice}\end{blocLivre}
\begin{blocLivre}{ex-10-13}\begin{exercice}[3]{Ajouter une observation}
Une série de 9 valeurs a une moyenne de 14. On ajoute 24. Calculer la nouvelle moyenne. Peut-on en déduire la nouvelle médiane ?
\end{exercice}\end{blocLivre}
\begin{blocLivre}{ex-10-14}\begin{exercice}[3]{Même moyenne, mêmes résultats ?}
Comparer les séries A : $9\,;10\,;10\,;11$ et B : $2\,;10\,;10\,;18$. Calculer moyenne, médiane et étendue, puis commenter.
\end{exercice}\end{blocLivre}
\begin{blocLivre}{ex-10-15}\begin{exercice}[3]{Un graphique trompeur}
Deux barres représentent 80 et 84 inscriptions. Un axe vertical commence à 78. Comparer la hausse réelle et les hauteurs visibles à partir de 78.
\label{fig-ex-10-15}\figureGeo{stat-axe-tronque}
\end{exercice}\end{blocLivre}

\clearpage

\begin{ouverture}[probabilites]{11}{Probabilités, évènements et simulations}{Partie 3 · Données et probabilités}
{Issues · Évènements · Réunion · Intersection · Tableaux · Fréquences}
\objectif{Décrire des évènements avec le vocabulaire des ensembles.}
\objectif{Utiliser la relation $P(A\cup B)+P(A\cap B)=P(A)+P(B)$.}
\objectif{Dénombrer les issues de deux épreuves indépendantes et interpréter une simulation.}
\prerequis{Fractions, probabilité d’un évènement et de son contraire.}
\end{ouverture}
\pagestyle{otMain}\essentiel\fontsize{11}{15}\selectfont
\essSection{1}{Évènements et opérations}
\begin{definition}Une issue est un résultat possible d’une expérience aléatoire. Un évènement est un ensemble d’issues. L’évènement contraire de $A$, noté $\overline A$, contient les issues qui ne réalisent pas $A$. L’évènement impossible correspond à l’ensemble vide $\varnothing$.\end{definition}
\begin{definition}L’intersection $A\cap B$ (« $A$ et $B$ ») contient les issues communes à $A$ et $B$. La réunion $A\cup B$ (« $A$ ou $B$ ») contient les issues de $A$, de $B$, ou des deux. $A$ et $B$ sont incompatibles si $A\cap B=\varnothing$.\end{definition}
\begin{propriete}Si les issues sont équiprobables, $P(A)$ est le nombre d’issues favorables divisé par le nombre total d’issues. De plus $P(\overline A)=1-P(A)$ et
\[P(A\cup B)+P(A\cap B)=P(A)+P(B).\]\end{propriete}
\begin{exemple}Avec un dé équilibré, $A$ : « pair » $=\{2;4;6\}$ et $B$ : « au moins 4 » $=\{4;5;6\}$. $A\cap B=\{4;6\}$ et $A\cup B=\{2;4;5;6\}$ : on vérifie $4/6+2/6=3/6+3/6$.\end{exemple}
\begin{attention}Le « ou » des mathématiques est inclusif. « $A$ ou $B$, mais pas les deux » se dit « ou exclusif » ; sa probabilité vaut $P(A\cup B)-P(A\cap B)$.\end{attention}
\clearpage
\essSection{2}{Deux épreuves et simulations}
\begin{savoirfaire}[Dénombrer]Pour deux épreuves indépendantes, un tableau à double entrée ou un arbre de dénombrement liste toutes les issues. Si chaque épreuve a des issues équiprobables, les couples obtenus le sont aussi : deux dés donnent $6\times6=36$ couples équiprobables.\end{savoirfaire}
\begin{exemple}Deux pièces équilibrées lancées indépendamment donnent quatre issues équiprobables : PP, PF, FP, FF. L’évènement « exactement une face F » contient PF et FP, donc a pour probabilité $2/4=1/2$.\end{exemple}
\begin{savoirfaire}[Numéroter pour compter]Dans un sac de 3 boules rouges et 2 bleues, numéroter les boules $R_1,R_2,R_3,B_1,B_2$. Avec remise, deux tirages donnent $5\times5=25$ couples équiprobables, dont $3\times3=9$ formés de deux rouges.\end{savoirfaire}
\begin{propriete}Simuler une expérience, c’est la reproduire par un dispositif ou un programme qui respecte le modèle, par exemple un nombre entier aléatoire de 1 à 6 pour un dé. Quand le nombre de répétitions augmente, la fréquence observée se stabilise généralement autour de la probabilité. Cela ne prédit ni le prochain résultat ni une égalité exacte.\end{propriete}
\clearpage\rubrique{Découverte}
\begin{blocLivre}{ex-11-1}\begin{exercice}[1]{Un dé}
On lance un dé équilibré à six faces numérotées de 1 à 6. Quelle est la probabilité d’obtenir un nombre pair ?
\end{exercice}\end{blocLivre}
\begin{blocLivre}{ex-11-2}\begin{exercice}[1]{Des boules colorées}
Un sac contient 5 boules rouges, 3 bleues et 2 vertes, indiscernables au toucher. On tire une boule au hasard. Donner les probabilités de rouge et de non-rouge.
\end{exercice}\end{blocLivre}
\begin{blocLivre}{ex-11-3}\begin{exercice}[1]{Écrire des évènements}
On lance un dé équilibré. $A$ est l’évènement « obtenir un nombre pair » et $B$ l’évènement « obtenir un multiple de 3 ». Écrire $A$, $B$, $A\cap B$, $A\cup B$ et $\overline A$ sous forme d’ensembles.
\end{exercice}\end{blocLivre}
\begin{blocLivre}{ex-11-4}\begin{exercice}[1]{Impossible ou certain}
Sur un dé classique, donner les probabilités d’obtenir 9 et d’obtenir moins de 7. Comment note-t-on l’évènement impossible ?
\end{exercice}\end{blocLivre}
\begin{blocLivre}{ex-11-5}\begin{exercice}[1]{Une fréquence expérimentale}
On obtient 46 fois face au cours de 100 lancers d’une pièce. Donner la fréquence de face. Est-ce une preuve que la pièce est déséquilibrée ?
\end{exercice}\end{blocLivre}

\clearpage\rubrique{Entraînement}
\begin{blocLivre}{ex-11-6}\begin{exercice}[2]{La relation de la réunion}
Avec les évènements de l’exercice 3, calculer $P(A)$, $P(B)$ et $P(A\cap B)$, puis en déduire $P(A\cup B)$ par la relation du cours. Contrôler par un dénombrement.
\end{exercice}\end{blocLivre}
\begin{blocLivre}{ex-11-7}\begin{exercice}[2]{Deux pièces}
Deux pièces équilibrées sont lancées indépendamment. Compléter l’arbre de dénombrement et calculer la probabilité d’obtenir exactement une face F, puis au moins une face F.
\label{fig-ex-11-7}\figureGeo{proba-pieces}
\end{exercice}\end{blocLivre}
\begin{blocLivre}{ex-11-8}\begin{exercice}[2]{La somme de deux dés}
Deux dés équilibrés sont lancés indépendamment. Utiliser le tableau des sommes pour calculer la probabilité d’une somme égale à 7.
\label{fig-ex-11-8}\figureGeo{proba-des}
\end{exercice}\end{blocLivre}
\begin{blocLivre}{ex-11-9}\begin{exercice}[2]{Un « ou » exclusif}
On lance un dé équilibré. $A$ : « pair » et $B$ : « au moins 4 ». Calculer la probabilité de l’évènement « $A$ ou $B$, mais pas les deux ».
\end{exercice}\end{blocLivre}
\begin{blocLivre}{ex-11-10}\begin{exercice}[2]{Au moins un six}
Deux dés équilibrés sont lancés. On note $A$ : « le premier dé donne 6 » et $B$ : « le second dé donne 6 ». Calculer la probabilité d’obtenir au moins un six.
\label{fig-ex-11-10}\figureGeo{proba-des}
\end{exercice}\end{blocLivre}

\clearpage\rubrique{Problèmes}
\begin{blocLivre}{ex-11-11}\begin{exercice}[3]{Deux tirages avec remise}
Un sac contient 3 boules rouges $R_1,R_2,R_3$ et 2 bleues $B_1,B_2$. On tire une boule, on la remet, puis on en tire une seconde. Utiliser le tableau pour calculer la probabilité de deux rouges, puis de deux boules de même couleur.
\label{fig-ex-11-11}\figureGeo{proba-tableau-boules}
\end{exercice}\end{blocLivre}
\begin{blocLivre}{ex-11-12}\begin{exercice}[3]{Un jeu à points}
On paie 2 points pour lancer un dé équilibré. Le 6 rapporte 8 points, le 5 rapporte 4 points, les autres 0. Sur un grand nombre de parties, le gain brut moyen théorique est la somme des gains multipliés par leurs probabilités. Le jeu équilibre-t-il mise et gain moyen ?
\end{exercice}\end{blocLivre}
\begin{blocLivre}{ex-11-13}\begin{exercice}[3]{Un effectif inconnu}
Un sac contient 6 rouges et $b$ bleues seulement. La probabilité de tirer rouge au hasard est $2/5$. Trouver $b$.
\end{exercice}\end{blocLivre}
\begin{blocLivre}{ex-11-14}\begin{exercice}[3]{Un code à deux chiffres}
Un générateur choisit indépendamment et au hasard deux chiffres parmi 0, 1, 2, 3, avec répétition possible. Quelle est la probabilité que ses deux chiffres soient distincts ? Que leur somme soit paire ou que le premier chiffre soit 0 ?
\end{exercice}\end{blocLivre}
\begin{blocLivre}{ex-11-15}\begin{exercice}[3]{Estimer sans promettre}
Un programme simule un évènement de probabilité $1/3$ à l’aide de nombres aléatoires. Il obtient 327 succès en 1000 essais, puis 3354 en 10000 essais. Comparer les écarts entre fréquences et probabilité. Peut-on garantir que tout essai supplémentaire rapproche la fréquence de $1/3$ ?
\end{exercice}\end{blocLivre}

\clearpage

\begin{ouverture}[geometrie]{12}{Pythagore, cercle et droite des milieux}{Partie 4 · Espace et géométrie}
{Pythagore · Réciproque · Contraposée · Cercle circonscrit · Droite des milieux}
\objectif{Calculer une longueur dans un triangle rectangle.}
\objectif{Prouver ou réfuter qu’un triangle est rectangle.}
\objectif{Utiliser le cercle circonscrit et les théorèmes des milieux.}
\prerequis{Carrés, racines carrées, médiatrices et parallélogrammes.}
\end{ouverture}
\pagestyle{otMain}\essentiel\fontsize{11}{15}\selectfont
\essSection{1}{Pythagore, réciproque et contraposée}
\begin{propriete}[Théorème de Pythagore]Si ABC est rectangle en A, alors $BC^2=AB^2+AC^2$. BC, opposé à l’angle droit, est l’hypoténuse et le plus long côté.\end{propriete}
\begin{savoirfaire}[Calculer une longueur]Pour $AB=6$ cm et $AC=8$ cm, $BC^2=36+64=100$, donc $BC=\sqrt{100}=10$ cm. Si $BC=13$ cm et $AB=5$ cm, $AC^2=169-25=144$, d’où $AC=12$ cm. Une longueur est positive.\end{savoirfaire}
\begin{propriete}[Réciproque]Dans un triangle, si le carré du plus long côté est égal à la somme des carrés des deux autres, alors le triangle est rectangle, et l’angle droit est opposé au plus long côté.\end{propriete}
\begin{definition}La réciproque de « si $P$, alors $Q$ » est « si $Q$, alors $P$ ». Sa contraposée est « si $Q$ est faux, alors $P$ est faux ». Une propriété vraie a une contraposée vraie ; sa réciproque doit être démontrée à part.\end{definition}
\begin{propriete}[Contraposée de Pythagore]Si le carré du plus long côté n’est pas égal à la somme des carrés des deux autres, alors le triangle n’est pas rectangle.\end{propriete}
\begin{attention}Le théorème exige un angle droit connu : une apparence de dessin ne suffit pas. Conserver la valeur exacte d’une racine, puis arrondir à la précision demandée.\end{attention}
\clearpage
\essSection{2}{Cercle circonscrit et droite des milieux}
\begin{propriete}Si ABC est rectangle en A, alors A appartient au cercle de diamètre [BC] : le centre du cercle circonscrit est le milieu de l’hypoténuse. Réciproquement, si A appartient au cercle de diamètre [BC], avec A distinct de B et de C, alors ABC est rectangle en A.\end{propriete}
\begin{propriete}[Droite des milieux]Dans un triangle ABC :
\par 1. la droite qui passe par les milieux de [AB] et de [AC] est parallèle à (BC) ;
\par 2. la droite qui passe par le milieu de [AB] parallèlement à (BC) coupe [AC] en son milieu ;
\par 3. le segment qui joint les milieux de [AB] et de [AC] mesure la moitié de BC.\end{propriete}
\begin{exemple}Si I et J sont les milieux de [AB] et [AC] et si $BC=9$ cm, alors $(IJ)\parallel(BC)$ et $IJ=4{,}5$ cm.\end{exemple}
\begin{savoirfaire}[Rédiger]Citer les hypothèses (angle droit, milieux, points sur un cercle), nommer la propriété utilisée, puis conclure en nommant précisément les objets. Pour montrer qu’un triangle n’est pas rectangle, le mot « contraposée » est attendu.\end{savoirfaire}
\clearpage\rubrique{Découverte}
\begin{blocLivre}{ex-12-1}\begin{exercice}[1]{Reconnaître l’hypoténuse}
ABC est rectangle en A. Nommer l’hypoténuse et écrire l’égalité de Pythagore.
\label{fig-ex-12-1}\figureGeo{pyth-01}
\end{exercice}\end{blocLivre}
\begin{blocLivre}{ex-12-2}\begin{exercice}[1]{Un triangle classique}
ABC est rectangle en A, avec AB = 6 cm et AC = 8 cm. Calculer BC.
\label{fig-ex-12-2}\figureGeo{pyth-02}
\end{exercice}\end{blocLivre}
\begin{blocLivre}{ex-12-3}\begin{exercice}[1]{Le côté manquant}
ABC est rectangle en A. BC = 13 cm et AB = 5 cm. Calculer AC.
\label{fig-ex-12-3}\figureGeo{pyth-03}
\end{exercice}\end{blocLivre}
\begin{blocLivre}{ex-12-4}\begin{exercice}[1]{Garder une racine}
ABC est rectangle en A. AB = 4 cm, AC = 7 cm. Donner BC exactement puis au dixième.
\label{fig-ex-12-4}\figureGeo{pyth-04}
\end{exercice}\end{blocLivre}
\begin{blocLivre}{ex-12-5}\begin{exercice}[1]{Le centre du cercle circonscrit}
ABC est rectangle en A, avec BC = 10 cm ; O est le milieu de [BC]. Où se trouve le centre du cercle circonscrit à ABC ? Donner son rayon et la longueur OA.
\label{fig-ex-12-5}\figureGeo{cercle-circonscrit}
\end{exercice}\end{blocLivre}

\clearpage\rubrique{Entraînement}
\begin{blocLivre}{ex-12-6}\begin{exercice}[2]{Tester la réciproque}
Dans ABC, AB = 9 cm, AC = 12 cm et BC = 15 cm. Le triangle est-il rectangle ?
\label{fig-ex-12-6}\figureGeo{pyth-05}
\end{exercice}\end{blocLivre}
\begin{blocLivre}{ex-12-7}\begin{exercice}[2]{Utiliser la contraposée}
AB = 5 cm, AC = 7 cm et BC = 9 cm. ABC est-il rectangle ? Nommer la propriété utilisée.
\label{fig-ex-12-7}\figureGeo{pyth-06}
\end{exercice}\end{blocLivre}
\begin{blocLivre}{ex-12-8}\begin{exercice}[2]{La droite des milieux}
Dans le triangle ABC, I est le milieu de [AB] et J celui de [AC] ; BC = 9 cm. Que peut-on dire des droites (IJ) et (BC) ? Calculer IJ.
\label{fig-ex-12-8}\figureGeo{milieux-parallele}
\end{exercice}\end{blocLivre}
\begin{blocLivre}{ex-12-9}\begin{exercice}[2]{Un milieu retrouvé}
Dans le triangle ABC, I est le milieu de [AB]. La parallèle à (BC) passant par I coupe [AC] en J. On sait que AC = 7,4 cm. Calculer AJ en justifiant.
\label{fig-ex-12-9}\figureGeo{milieux-milieu}
\end{exercice}\end{blocLivre}
\begin{blocLivre}{ex-12-10}\begin{exercice}[2]{Un angle à identifier}
AB = 10 cm, AC = 6 cm et BC = 8 cm. Prouver que le triangle est rectangle et préciser en quel sommet.
\label{fig-ex-12-10}\figureGeo{pyth-10}
\end{exercice}\end{blocLivre}

\clearpage\rubrique{Problèmes}
\begin{blocLivre}{ex-12-11}\begin{exercice}[3]{Une échelle contre un mur}
Le mur AC est vertical et le sol AB horizontal. Une échelle BC de 5 m est posée à 1,4 m du mur. Quelle hauteur AC atteint-elle ?
\label{fig-ex-12-11}\figureGeo{pyth-07}
\end{exercice}\end{blocLivre}
\begin{blocLivre}{ex-12-12}\begin{exercice}[3]{Un triangle dans un demi-cercle}
[BC] est un diamètre d’un cercle de centre O et de rayon 6,5 cm. A est un point de ce cercle, distinct de B et de C, tel que AB = 5 cm. Démontrer que ABC est rectangle, puis calculer AC.
\label{fig-ex-12-12}\figureGeo{demi-cercle}
\end{exercice}\end{blocLivre}
\begin{blocLivre}{ex-12-13}\begin{exercice}[3]{Un rectangle sans équerre}
On trace un cercle de centre O, puis deux diamètres [AC] et [BD] qui ne sont pas perpendiculaires. Démontrer que ABCD est un rectangle. Pourquoi ce procédé permet-il de construire un rectangle sans équerre ?
\label{fig-ex-12-13}\figureGeo{rectangle-compas}
\end{exercice}\end{blocLivre}
\begin{blocLivre}{ex-12-14}\begin{exercice}[3]{Une hauteur de triangle}
ABC est isocèle en C : CA = CB = 13 cm et AB = 10 cm. H est le milieu de [AB]. Calculer CH puis l’aire de ABC.
\label{fig-ex-12-14}\figureGeo{pyth-iso}
\end{exercice}\end{blocLivre}
\begin{blocLivre}{ex-12-15}\begin{exercice}[3]{Le théorème de Varignon}
ABCD est un quadrilatère quelconque, non croisé. I, J, K et L sont les milieux respectifs de [AB], [BC], [CD] et [DA]. Démontrer que IJKL est un parallélogramme.
\label{fig-ex-12-15}\figureGeo{varignon}
\end{exercice}\end{blocLivre}

\clearpage

\begin{ouverture}[geometrie]{13}{Thalès, réciproque et contraposée}{Partie 4 · Espace et géométrie}
{Alignements · Parallèles · Rapports égaux · Réciproque · Contraposée}
\objectif{Écrire les rapports de longueurs correspondantes.}
\objectif{Calculer une longueur dans les deux configurations.}
\objectif{Démontrer que des droites sont, ou ne sont pas, parallèles.}
\prerequis{Proportionnalité, fractions, produits en croix et droite des milieux.}
\end{ouverture}
\pagestyle{otMain}\essentiel\fontsize{11}{15}\selectfont
\essSection{1}{Le théorème de Thalès}
\begin{propriete}[Théorème de Thalès]Si A, M, B sont alignés, A, N, C sont alignés et $(MN)\parallel(BC)$, dans la configuration des triangles emboîtés ou dans celle dite du papillon, alors
\[\frac{AM}{AB}=\frac{AN}{AC}=\frac{MN}{BC}.\]
Les longueurs des côtés correspondants des triangles AMN et ABC sont proportionnelles.\end{propriete}
\begin{savoirfaire}[Associer les côtés]Dans les triangles AMN et ABC, A correspond à A, M à B et N à C. Tous les numérateurs appartiennent au même triangle. Choisir ensuite l’égalité contenant l’inconnue et deux données.\end{savoirfaire}
\begin{exemple}Si $AM=3$, $AB=5$, $AC=7$ et les hypothèses sont vérifiées, $AN/7=3/5$, donc $AN=4{,}2$ dans la même unité.\end{exemple}
\begin{attention}Si M est entre A et B, $MB=AB-AM$. Dans le papillon, où A est entre M et B, $MB=AM+AB$. Lire l’ordre des points avant de calculer.\end{attention}
\begin{exemple}[Une preuve par les aires]Dans le cas des triangles emboîtés, les triangles MNB et MNC ont la même base [MN] et la même hauteur, car $(BC)\parallel(MN)$. Comparer les aires conduit à l’égalité des rapports : la démonstration complète figure en fin de livre.\end{exemple}
\clearpage
\essSection{2}{Réciproque, contraposée et proportionnalité}
\begin{propriete}[Réciproque]Si A, M, B et A, N, C sont alignés dans le même ordre et si $\dfrac{AM}{AB}=\dfrac{AN}{AC}$, alors $(MN)\parallel(BC)$.\end{propriete}
\begin{propriete}[Contraposée]Si A, M, B et A, N, C sont alignés et si $\dfrac{AM}{AB}\ne\dfrac{AN}{AC}$, alors les droites (MN) et (BC) ne sont pas parallèles : sinon le théorème imposerait l’égalité.\end{propriete}
\begin{savoirfaire}[Comparer deux rapports]Comparer exactement, par exemple avec des produits en croix : $3/8$ et $5/12$ diffèrent, car $3\times12=36$ et $5\times8=40$. Un arrondi ne permet jamais de conclure à l’égalité.\end{savoirfaire}
\begin{propriete}[Lien avec la proportionnalité]Dans un repère, des points alignés avec l’origine ont des coordonnées proportionnelles : Thalès s’applique aux triangles formés avec leurs projetés sur l’axe des abscisses.\end{propriete}
\begin{savoirfaire}[Rédiger]Énoncer les alignements et le parallélisme, nommer le théorème, écrire l’égalité des rapports puis calculer. Dans une réciproque ou une contraposée, conclure en nommant les deux droites.\end{savoirfaire}
\clearpage\rubrique{Découverte}
\begin{blocLivre}{ex-13-1}\begin{exercice}[1]{Écrire les rapports}
Dans la figure, M appartient à [AB], N à [AC] et $(MN)\parallel(BC)$. Écrire les rapports égaux.
\label{fig-ex-13-1}\figureGeo{thales-01}
\end{exercice}\end{blocLivre}
\begin{blocLivre}{ex-13-2}\begin{exercice}[1]{Une longueur réduite}
M est sur [AB], N sur [AC] et $(MN)\parallel(BC)$. AM = 3 cm, AB = 5 cm, AC = 10 cm. Calculer AN.
\label{fig-ex-13-2}\figureGeo{thales-02}
\end{exercice}\end{blocLivre}
\begin{blocLivre}{ex-13-3}\begin{exercice}[1]{La petite base}
Dans la même configuration parallèle, AM = 4 cm, AB = 10 cm et BC = 15 cm. Calculer MN.
\label{fig-ex-13-3}\figureGeo{thales-03}
\end{exercice}\end{blocLivre}
\begin{blocLivre}{ex-13-4}\begin{exercice}[1]{Retrouver le grand côté}
M est sur [AB], N sur [AC], $(MN)\parallel(BC)$. AM = 6 cm, AN = 4 cm et AC = 10 cm. Calculer AB.
\label{fig-ex-13-4}\figureGeo{thales-04}
\end{exercice}\end{blocLivre}
\begin{blocLivre}{ex-13-5}\begin{exercice}[1]{De la partie au tout}
M est sur [AB], N sur [AC], $(MN)\parallel(BC)$. AM = 3 cm, MB = 2 cm et AC = 12 cm. Calculer AN.
\label{fig-ex-13-5}\figureGeo{thales-05}
\end{exercice}\end{blocLivre}

\clearpage\rubrique{Entraînement}
\begin{blocLivre}{ex-13-6}\begin{exercice}[2]{Démontrer un parallélisme}
A, M, B sont alignés dans cet ordre ; A, N, C aussi. AM = 4, AB = 10, AN = 6 et AC = 15, en cm. Montrer que (MN) et (BC) sont parallèles.
\label{fig-ex-13-6}\figureGeo{thales-06}
\end{exercice}\end{blocLivre}
\begin{blocLivre}{ex-13-7}\begin{exercice}[2]{Des droites non parallèles}
A, M, B et A, N, C sont alignés dans cet ordre. AM = 3, AB = 8, AN = 5, AC = 12, en cm. Les droites (MN) et (BC) sont-elles parallèles ?
\label{fig-ex-13-7}\figureGeo{thales-07}
\end{exercice}\end{blocLivre}
\begin{blocLivre}{ex-13-8}\begin{exercice}[2]{Le papillon}
M, A, B sont alignés dans cet ordre ; N, A, C aussi. $(MN)\parallel(BC)$, AM = 3 cm, AB = 5 cm et AC = 8 cm. Calculer AN.
\label{fig-ex-13-8}\figureGeo{thales-08}
\end{exercice}\end{blocLivre}
\begin{blocLivre}{ex-13-9}\begin{exercice}[2]{La distance entre deux points}
Dans le papillon précédent, AM = 3 cm et AB = 5 cm. Calculer MB. Si AN = 4,8 cm et AC = 8 cm, calculer NC.
\label{fig-ex-13-9}\figureGeo{thales-09}
\end{exercice}\end{blocLivre}
\begin{blocLivre}{ex-13-10}\begin{exercice}[2]{Une réciproque en papillon}
M, A, B sont alignés dans cet ordre, ainsi que N, A, C. AM = 2,4 cm, AB = 4 cm, AN = 3 cm et AC = 5 cm. Les droites (MN) et (BC) sont-elles parallèles ?
\label{fig-ex-13-10}\figureGeo{thales-10}
\end{exercice}\end{blocLivre}

\clearpage\rubrique{Problèmes}
\begin{blocLivre}{ex-13-11}\begin{exercice}[3]{Deux triangles de même aire}
Dans le triangle ABC, M appartient à [AB], N à [AC] et $(MN)\parallel(BC)$. Expliquer pourquoi les triangles MNB et MNC ont la même aire. En déduire que les triangles AMC et ANB ont la même aire.
\label{fig-ex-13-11}\figureGeo{thales-aires}
\end{exercice}\end{blocLivre}
\begin{blocLivre}{ex-13-12}\begin{exercice}[3]{L’ombre d’un arbre}
Un bâton vertical de 1,5 m est planté de sorte que l’extrémité S de son ombre coïncide avec celle de l’ombre d’un arbre vertical. L’ombre du bâton mesure 2 m, celle de l’arbre 12 m, sur un sol horizontal. Les rayons du soleil sont rectilignes. Calculer la hauteur de l’arbre.
\label{fig-ex-13-12}\figureGeo{thales-ombre}
\end{exercice}\end{blocLivre}
\begin{blocLivre}{ex-13-13}\begin{exercice}[3]{Une équation dans Thalès}
M est sur [AB], N sur [AC], (MN) est parallèle à (BC). AM = 4 cm, MB = $x$ cm, AN = 6 cm et NC = 3 cm. Trouver $x$.
\label{fig-ex-13-13}\figureGeo{thales-13}
\end{exercice}\end{blocLivre}
\begin{blocLivre}{ex-13-14}\begin{exercice}[3]{Calculer une largeur inaccessible}
A, B, C forment un triangle. M est le milieu de [AB] et N le milieu de [AC] ; [BC] traverse une zone inaccessible. On mesure MN = 7,4 m. Justifier le parallélisme de (MN) et (BC), puis calculer BC.
\label{fig-ex-13-14}\figureGeo{thales-milieux}
\end{exercice}\end{blocLivre}
\begin{blocLivre}{ex-13-15}\begin{exercice}[3]{Un contrôle critique}
Dans une configuration emboîtée avec (MN) parallèle à (BC), un élève écrit $AM/AB=NC/AC$. AM = 3 cm, AB = 8 cm, AC = 16 cm. Corriger et calculer NC.
\label{fig-ex-13-15}\figureGeo{thales-15}
\end{exercice}\end{blocLivre}

\clearpage

\begin{ouverture}[geometrie]{14}{Trigonométrie du triangle rectangle}{Partie 4 · Espace et géométrie}
{Cosinus · Sinus · Tangente · Longueurs · Angles}
\objectif{Choisir le rapport adapté à l’angle et aux données.}
\objectif{Calculer une longueur ou un angle aigu.}
\objectif{Démontrer et utiliser $(\cos\widehat B)^2+(\sin\widehat B)^2=1$.}
\prerequis{Pythagore, Thalès et proportionnalité.}
\end{ouverture}
\pagestyle{otMain}\essentiel\fontsize{11}{15}\selectfont
\essSection{1}{Trois rapports pour un angle aigu}
\begin{definition}Dans ABC rectangle en A, pour l’angle $\widehat B$, BC est l’hypoténuse, AB le côté adjacent et AC le côté opposé. Les deux derniers rôles s’échangent si l’on choisit l’angle $\widehat C$.\end{definition}
\begin{propriete}$\cos\widehat B=\dfrac{AB}{BC}$, $\sin\widehat B=\dfrac{AC}{BC}$ et $\tan\widehat B=\dfrac{AC}{AB}$. Ces rapports ne dépendent que de l’angle : deux triangles rectangles ayant le même angle aigu peuvent être placés l’un dans l’autre, et le théorème de Thalès donne les mêmes rapports.\end{propriete}
\begin{savoirfaire}[Trouver une longueur]Écrire le rapport avec les noms des côtés avant de remplacer. Si BC = 10 cm et $\widehat B=35^\circ$, $AC=10\sin35^\circ\approx5{,}74$ cm. Si AB = 10 cm, on utiliserait $AC=10\tan35^\circ$.\end{savoirfaire}
\begin{propriete}Pour tout angle aigu, $(\cos\widehat B)^2+(\sin\widehat B)^2=1$. En effet, la somme vaut $\dfrac{AB^2+AC^2}{BC^2}$, et $AB^2+AC^2=BC^2$ d’après Pythagore.\end{propriete}
\begin{attention}La calculatrice doit être en mode degrés. Un sinus ou cosinus d’angle aigu est entre 0 et 1 ; une hypoténuse reste le plus long côté.\end{attention}
\clearpage
\essSection{2}{Angles et problèmes}
\begin{savoirfaire}[Trouver un angle]Si AB = 6 cm et BC = 10 cm, $\cos\widehat B=0{,}6$. Utiliser la fonction inverse du cosinus, souvent notée $\cos^{-1}$ ou acos : $\widehat B\approx53{,}1^\circ$. Cette touche ne calcule pas $1/\cos$.\end{savoirfaire}
\begin{propriete}Les deux angles aigus d’un triangle rectangle ont pour somme $90^\circ$. Cette relation fournit un contrôle ou un autre moyen de calcul.\end{propriete}
\begin{savoirfaire}[Une mesure indirecte]Pour une hauteur inaccessible, dessiner l’horizontale au niveau de l’observateur et sa ligne de visée. La hauteur calculée par tangente est celle au-dessus de cette horizontale ; ajouter ensuite la hauteur de l’instrument si nécessaire.\end{savoirfaire}
\begin{attention}Conserver les valeurs de calculatrice non arrondies jusqu’à la dernière étape. Préciser l’unité et la précision demandée. Un angle observé et une distance mesurés restent des données de modèle, avec une précision limitée.\end{attention}
\clearpage\rubrique{Découverte}
\begin{blocLivre}{ex-14-1}\begin{exercice}[1]{Nommer les côtés}
ABC est rectangle en A. Pour l’angle $\widehat B$, identifier le côté opposé, l’adjacent et l’hypoténuse.
\label{fig-ex-14-1}\figureGeo{trigo-01}
\end{exercice}\end{blocLivre}
\begin{blocLivre}{ex-14-2}\begin{exercice}[1]{Choisir le cosinus}
ABC est rectangle en A. AB = 6 cm et BC = 10 cm. Donner $\cos\widehat B$.
\label{fig-ex-14-2}\figureGeo{trigo-02}
\end{exercice}\end{blocLivre}
\begin{blocLivre}{ex-14-3}\begin{exercice}[1]{Choisir le sinus}
ABC est rectangle en A. AC = 8 cm et BC = 10 cm. Donner $\sin\widehat B$.
\label{fig-ex-14-3}\figureGeo{trigo-03}
\end{exercice}\end{blocLivre}
\begin{blocLivre}{ex-14-4}\begin{exercice}[1]{Choisir la tangente}
ABC est rectangle en A. AB = 6 cm et AC = 8 cm. Donner $\tan\widehat B$.
\label{fig-ex-14-4}\figureGeo{trigo-04}
\end{exercice}\end{blocLivre}
\begin{blocLivre}{ex-14-5}\begin{exercice}[1]{Une première longueur}
ABC est rectangle en A, $\widehat B=30^\circ$ et BC = 12 cm. Calculer AC.
\label{fig-ex-14-5}\figureGeo{trigo-05}
\end{exercice}\end{blocLivre}

\clearpage\rubrique{Entraînement}
\begin{blocLivre}{ex-14-6}\begin{exercice}[2]{Un côté adjacent}
ABC est rectangle en A. $\widehat B=40^\circ$ et BC = 9 cm. Calculer AB au millimètre.
\label{fig-ex-14-6}\figureGeo{trigo-06}
\end{exercice}\end{blocLivre}
\begin{blocLivre}{ex-14-7}\begin{exercice}[2]{Une hypoténuse}
ABC est rectangle en A. $\widehat B=50^\circ$ et AB = 7 cm. Calculer BC au dixième de centimètre.
\label{fig-ex-14-7}\figureGeo{trigo-07}
\end{exercice}\end{blocLivre}
\begin{blocLivre}{ex-14-8}\begin{exercice}[2]{Un angle par cosinus}
ABC est rectangle en A, AB = 6 cm et BC = 10 cm. Calculer $\widehat B$ au degré.
\label{fig-ex-14-8}\figureGeo{trigo-08}
\end{exercice}\end{blocLivre}
\begin{blocLivre}{ex-14-9}\begin{exercice}[2]{Un angle par tangente}
ABC est rectangle en A, AC = 5 cm et AB = 12 cm. Calculer $\widehat B$ au dixième de degré.
\label{fig-ex-14-9}\figureGeo{trigo-09}
\end{exercice}\end{blocLivre}
\begin{blocLivre}{ex-14-10}\begin{exercice}[2]{Cosinus et sinus liés}
ABC est rectangle en A. Démontrer que $(\cos\widehat B)^2+(\sin\widehat B)^2=1$. Sachant que $\cos\widehat B=0{,}6$, en déduire $\sin\widehat B$.
\label{fig-ex-14-10}\figureGeo{trigo-10}
\end{exercice}\end{blocLivre}

\clearpage\rubrique{Problèmes}
\begin{blocLivre}{ex-14-11}\begin{exercice}[3]{La hauteur d’un bâtiment}
Un appareil à 1,6 m du sol vise le sommet d’un bâtiment avec un angle d’élévation de 32 degrés. Il est à 25 m horizontalement du pied. Le sol est horizontal. Calculer la hauteur au décimètre.
\label{fig-ex-14-11}\figureGeo{trigo-batiment}
\end{exercice}\end{blocLivre}
\begin{blocLivre}{ex-14-12}\begin{exercice}[3]{Une rampe et son angle}
Une rampe monte de 0,8 m pour 10 m d’avancée horizontale. Dans ABC, A est l’angle droit, AB l’avancée et AC la hauteur. Calculer $\widehat B$ au dixième et la pente définie par $100\,AC/AB$.
\label{fig-ex-14-12}\figureGeo{trigo-12}
\end{exercice}\end{blocLivre}
\begin{blocLivre}{ex-14-13}\begin{exercice}[3]{Un câble}
Un câble BC de 18 m fait un angle $\widehat B$ de 65 degrés avec le sol horizontal AB. AC est vertical. Calculer la hauteur AC et la distance AB au dixième de mètre.
\label{fig-ex-14-13}\figureGeo{trigo-13}
\end{exercice}\end{blocLivre}
\begin{blocLivre}{ex-14-14}\begin{exercice}[3]{Un toit symétrique}
Le toit a pour section un triangle isocèle de base 8 m ; chaque versant fait 35 degrés avec l’horizontale. Calculer sa hauteur puis la longueur d’un versant au centimètre.
\label{fig-ex-14-14}\figureGeo{trigo-toit}
\end{exercice}\end{blocLivre}
\begin{blocLivre}{ex-14-15}\begin{exercice}[3]{Diagnostiquer une erreur}
ABC est rectangle en A. $\widehat B=38^\circ$, AB = 7 cm. Un élève annonce AC = $7\sin38^\circ$. Corriger puis donner AC au dixième.
\label{fig-ex-14-15}\figureGeo{trigo-15}
\end{exercice}\end{blocLivre}

\clearpage

\begin{ouverture}[vecteursplan]{15}{Translations et vecteurs}{Partie 4 · Espace et géométrie}
{Symétries · Translations · Vecteurs égaux · Somme · Relation de Chasles}
\objectif{Construire l’image d’une figure par symétrie, demi-tour ou translation.}
\objectif{Reconnaître des vecteurs égaux et un parallélogramme.}
\objectif{Construire une somme de vecteurs et utiliser la relation de Chasles.}
\prerequis{Symétrie axiale, demi-tour, translation et parallélogrammes.}
\end{ouverture}
\pagestyle{otMain}\essentiel\fontsize{11}{15}\selectfont
\essSection{1}{Translations et vecteurs}
\begin{definition}La translation qui transforme A en B transforme tout point C en l’unique point D tel que ABDC soit un parallélogramme, éventuellement aplati. On l’appelle la translation de vecteur $\vect{AB}$.\end{definition}
\begin{definition}Pour deux points distincts A et B, le vecteur $\vect{AB}$ est défini par la direction de la droite (AB), le sens de A vers B et la longueur AB, appelée norme. On le représente par une flèche ; deux vecteurs sont égaux s’ils ont même direction, même sens et même longueur.\end{definition}
\begin{propriete}$\vect{AB}=\vect{CD}$ si et seulement si ABDC est un parallélogramme, éventuellement aplati. Deux vecteurs égaux définissent la même translation.\end{propriete}
\begin{definition}Le vecteur nul est $\vect{AA}=\vect{BB}=\vec0$ : la translation correspondante ne déplace aucun point. L’opposé de $\vect{AB}$ est $-\vect{AB}=\vect{BA}$, de même direction et même longueur, mais de sens contraire.\end{definition}
\begin{attention}L’ordre des lettres compte : $\vect{AB}=\vect{DC}$ signifie que ABCD est un parallélogramme, alors que $\vect{AB}=\vect{CD}$ correspond à ABDC. Symétries, demi-tours et translations conservent longueurs, angles, alignements et aires.\end{attention}
\clearpage
\essSection{2}{Somme de vecteurs et relation de Chasles}
\begin{propriete}Enchaîner la translation de vecteur $\vect{AB}$ puis celle de vecteur $\vect{BC}$ revient à appliquer la translation de vecteur $\vect{AC}$. On définit ainsi la somme :
\[\vect{AB}+\vect{BC}=\vect{AC}\qquad\text{(relation de Chasles).}\]
En particulier $\vect{AB}+\vect{BA}=\vect{AA}=\vec0$.\end{propriete}
\begin{savoirfaire}[Construire $\vec u+\vec v$]À partir d’un point O, tracer un représentant $\vect{OP}$ de $\vec u$, puis, à partir de P, un représentant $\vect{PS}$ de $\vec v$. Alors $\vec u+\vec v=\vect{OS}$. On note $2\vect{AB}$ la somme $\vect{AB}+\vect{AB}$ : même direction, même sens, longueur double.\end{savoirfaire}
\begin{propriete}[Règle du parallélogramme]Si ABCD est un parallélogramme, alors $\vect{AB}+\vect{AD}=\vect{AC}$ : on remplace $\vect{AD}$ par le vecteur égal $\vect{BC}$, puis on applique Chasles.\end{propriete}
\begin{attention}La somme de deux vecteurs n’est pas la somme de leurs longueurs : $AB+BC$ n’est égal à $AC$ que si B appartient au segment [AC].\end{attention}
\clearpage\rubrique{Découverte}
\begin{blocLivre}{ex-15-1}\begin{exercice}[1]{Symétrie et demi-tour}
Sur le quadrillage, construire l’image A’B’C’ du triangle ABC par la symétrie d’axe (d), puis son image A’’B’’C’’ par le demi-tour de centre O. Ces transformations conservent-elles les longueurs ?
\label{fig-ex-15-1}\figureGeo{vect-01}
\end{exercice}\end{blocLivre}
\begin{blocLivre}{ex-15-2}\begin{exercice}[1]{Caractériser un vecteur}
Le côté d’un carreau représente 1 cm. Décrire le vecteur $\vect{AB}$ : direction, sens et longueur. Construire ensuite le point D tel que $\vect{CD}=\vect{AB}$.
\label{fig-ex-15-2}\figureGeo{vect-02}
\end{exercice}\end{blocLivre}
\begin{blocLivre}{ex-15-3}\begin{exercice}[1]{Des vecteurs égaux}
Parmi les vecteurs $\vect{CD}$, $\vect{EF}$, $\vect{GH}$ et $\vect{IJ}$ dessinés, lesquels sont égaux à $\vect{AB}$ ? Lequel est son opposé ?
\label{fig-ex-15-3}\figureGeo{vect-03}
\end{exercice}\end{blocLivre}
\begin{blocLivre}{ex-15-4}\begin{exercice}[1]{Image par une translation}
Construire l’image D du point C par la translation de vecteur $\vect{AB}$. Quelle est la nature du quadrilatère ABDC ?
\label{fig-ex-15-4}\figureGeo{vect-04}
\end{exercice}\end{blocLivre}
\begin{blocLivre}{ex-15-5}\begin{exercice}[1]{Vecteur nul et opposé}
Simplifier $\vect{AA}$. Donner un autre nom du vecteur $-\vect{MN}$, puis calculer $\vect{MN}+\vect{NM}$.
\label{fig-ex-15-5}\figureGeo{vect-05}
\end{exercice}\end{blocLivre}

\clearpage\rubrique{Entraînement}
\begin{blocLivre}{ex-15-6}\begin{exercice}[2]{Parallélogramme et égalités}
ABCD est un parallélogramme. Compléter avec des points de la figure : $\vect{AB}=\vect{D\,\ldots}$ ; $\vect{AD}=\vect{\ldots\,C}$ ; $\vect{BA}=\vect{C\,\ldots}$ ; $\vect{CB}=\vect{\ldots\,A}$.
\label{fig-ex-15-6}\figureGeo{vect-06}
\end{exercice}\end{blocLivre}
\begin{blocLivre}{ex-15-7}\begin{exercice}[2]{Démontrer un parallélogramme}
On sait que $\vect{EF}=\vect{HG}$. Démontrer que EFGH est un parallélogramme, puis vérifier cette égalité sur le quadrillage.
\label{fig-ex-15-7}\figureGeo{vect-07}
\end{exercice}\end{blocLivre}
\begin{blocLivre}{ex-15-8}\begin{exercice}[2]{La relation de Chasles}
Simplifier : $\vect{AB}+\vect{BC}$ ; $\vect{MN}+\vect{NP}+\vect{PQ}$ ; $\vect{AB}+\vect{CA}$ ; $\vect{AB}+\vect{BA}$.
\label{fig-ex-15-8}\figureGeo{vect-08}
\end{exercice}\end{blocLivre}
\begin{blocLivre}{ex-15-9}\begin{exercice}[2]{Construire une somme}
Les vecteurs $\vec u$ et $\vec v$ sont dessinés ; le côté d’un carreau représente 1 cm. Construire un représentant de $\vec u+\vec v$ d’origine O, puis calculer sa longueur.
\label{fig-ex-15-9}\figureGeo{vect-09}
\end{exercice}\end{blocLivre}
\begin{blocLivre}{ex-15-10}\begin{exercice}[2]{Deux translations successives}
On applique au point M la translation de vecteur $\vect{AB}$, ce qui donne M’, puis la translation de vecteur $\vect{BC}$, ce qui donne M’’. Construire M’ et M’’. Par quelle translation unique passe-t-on de M à M’’ ?
\label{fig-ex-15-10}\figureGeo{vect-10}
\end{exercice}\end{blocLivre}

\clearpage\rubrique{Problèmes}
\begin{blocLivre}{ex-15-11}\begin{exercice}[3]{La règle du parallélogramme}
ABCD est un parallélogramme. Démontrer que $\vect{AB}+\vect{AD}=\vect{AC}$.
\label{fig-ex-15-11}\figureGeo{vect-11}
\end{exercice}\end{blocLivre}
\begin{blocLivre}{ex-15-12}\begin{exercice}[3]{Milieu et vecteurs}
I est le milieu de [AB]. Démontrer que $\vect{AI}=\vect{IB}$ et que $\vect{AB}=2\vect{AI}$. Réciproquement, si un point J vérifie $\vect{AJ}=\vect{JB}$, que peut-on dire de J ?
\label{fig-ex-15-12}\figureGeo{vect-12}
\end{exercice}\end{blocLivre}
\begin{blocLivre}{ex-15-13}\begin{exercice}[3]{Un drone}
Un drone part de D, se déplace de 60 m vers l’est, puis de 80 m vers le nord. Représenter les deux déplacements par des vecteurs sur le quadrillage, où un carreau représente 20 m. Quel déplacement unique produit le même effet ? Quelle distance sépare le point d’arrivée F du point de départ ?
\label{fig-ex-15-13}\figureGeo{vect-13}
\end{exercice}\end{blocLivre}
\begin{blocLivre}{ex-15-14}\begin{exercice}[3]{Une frise}
Une frise est obtenue en répétant un motif par la translation de vecteur $\vec u$, de longueur 4 cm. Quelle translation transforme le motif 1 en motif 6 ? Le motif 6 en motif 2 ? Donner les longueurs des vecteurs correspondants.
\label{fig-ex-15-14}\figureGeo{vect-frise}
\end{exercice}\end{blocLivre}
\begin{blocLivre}{ex-15-15}\begin{exercice}[3]{Démontrer avec des vecteurs}
ABCD est un parallélogramme et E est le point tel que $\vect{BE}=\vect{AB}$. Démontrer que DCEB est un parallélogramme.
\label{fig-ex-15-15}\figureGeo{vect-15}
\end{exercice}\end{blocLivre}

\clearpage

\begin{ouverture}[espace]{16}{Espace : sphère, boule et sections}{Partie 4 · Espace et géométrie}
{Solides · Volumes · Sphère et boule · Grands cercles · Sections planes}
\objectif{Reconnaître les solides et calculer leurs volumes.}
\objectif{Définir la sphère, la boule et leurs grands cercles.}
\objectif{Visualiser et calculer des sections planes.}
\prerequis{Aires usuelles, Pythagore, Thalès et unités de volume.}
\end{ouverture}
\pagestyle{otMain}\essentiel\fontsize{11}{15}\selectfont
\essSection{1}{Solides, sphère et boule}
\begin{propriete}Pour un pavé : $V=L\ell h$. Pour un prisme droit ou un cylindre : $V=\mathcal A_{\rm base}\times h$. Pour une pyramide ou un cône : $V=\mathcal A_{\rm base}\times h/3$. La hauteur est perpendiculaire à la base.\end{propriete}
\begin{definition}La sphère de centre O et de rayon R est l’ensemble des points de l’espace situés à la distance R de O. La boule de centre O et de rayon R est l’ensemble des points situés à une distance de O inférieure ou égale à R.\end{definition}
\begin{definition}Un grand cercle de la sphère est un cercle de centre O et de rayon R tracé sur la sphère. Deux points A et B de la sphère sont diamétralement opposés lorsque O est le milieu de [AB] ; alors $AB=2R$ et [AB] est un diamètre.\end{definition}
\begin{propriete}Le volume d’une boule de rayon $r$ est $V=\dfrac43\pi r^3$. Cette formule est admise.\end{propriete}
\begin{savoirfaire}[Assembler]Décomposer en solides usuels, additionner des volumes sans recouvrement ou soustraire un creux. Conserver $\pi$ jusqu’au résultat final, puis convertir : $1\text{ dm}^3=1\text{ L}$.\end{savoirfaire}
\clearpage
\essSection{2}{Sections planes}
\begin{propriete}La section d’un pavé par un plan parallèle à une face est un rectangle de mêmes dimensions que cette face ; par un plan parallèle à une arête, c’est un rectangle.\end{propriete}
\begin{propriete}La section d’un cylindre par un plan perpendiculaire à son axe est un disque de même rayon que la base ; par un plan parallèle à son axe, c’est un rectangle.\end{propriete}
\begin{propriete}La section d’une pyramide ou d’un cône par un plan parallèle à la base est une réduction de la base : même nature, dimensions proportionnelles. Le coefficient est le quotient de la distance au sommet par la hauteur.\end{propriete}
\begin{propriete}La section d’une boule par un plan est un disque. Si le plan est à la distance $d$ du centre, avec $d<R$, le rayon $r$ du disque vérifie $r^2=R^2-d^2$ d’après Pythagore. Le plan passant par le centre donne un grand cercle.\end{propriete}
\begin{savoirfaire}[Dessiner en vraie grandeur]Reporter les longueurs utiles à partir de la face avant de la perspective, ou les calculer, puis tracer la section à part, sans déformation.\end{savoirfaire}
\clearpage\rubrique{Découverte}
\begin{blocLivre}{ex-16-1}\begin{exercice}[1]{Un pavé}
Un pavé a pour dimensions 8 cm, 5 cm et 3 cm. Calculer son volume en cm$^3$ puis en litres.
\label{fig-ex-16-1}\figureGeo{espace-01}
\end{exercice}\end{blocLivre}
\begin{blocLivre}{ex-16-2}\begin{exercice}[1]{Un cylindre}
Un cylindre a un rayon de 3 cm et une hauteur de 10 cm. Donner son volume exact et au cm$^3$.
\label{fig-ex-16-2}\figureGeo{espace-02}
\end{exercice}\end{blocLivre}
\begin{blocLivre}{ex-16-3}\begin{exercice}[1]{Une pyramide}
Une pyramide a une base carrée de côté 6 cm et une hauteur de 9 cm. Calculer son volume. Quelle est la nature de sa base ?
\label{fig-ex-16-3}\figureGeo{espace-03}
\end{exercice}\end{blocLivre}
\begin{blocLivre}{ex-16-4}\begin{exercice}[1]{Un cône}
Un cône a pour diamètre 8 cm et hauteur 9 cm. Calculer son volume exact.
\label{fig-ex-16-4}\figureGeo{espace-04}
\end{exercice}\end{blocLivre}
\begin{blocLivre}{ex-16-5}\begin{exercice}[1]{Une boule}
Une boule a un rayon de 3 cm. Donner son volume exact puis arrondi au cm$^3$. Quelle est la longueur d’un grand cercle de sa sphère ?
\label{fig-ex-16-5}\figureGeo{espace-05}
\end{exercice}\end{blocLivre}

\clearpage\rubrique{Entraînement}
\begin{blocLivre}{ex-16-6}\begin{exercice}[2]{Sphère, boule et grands cercles}
La sphère de centre O a pour rayon 4 cm. A et B sont diamétralement opposés, C est sur la sphère, OM = 2 cm et ON = 6 cm. Quels points appartiennent à la sphère ? À la boule ? Calculer AB et la longueur du grand cercle passant par A et B.
\label{fig-ex-16-6}\figureGeo{espace-sphere-vocab}
\end{exercice}\end{blocLivre}
\begin{blocLivre}{ex-16-7}\begin{exercice}[2]{Une pyramide dans un cube}
ABCDEFGH est un cube d’arête 6 cm ; E est le sommet situé au-dessus de A. On considère la pyramide EABCD de sommet E et de base ABCD. Donner la nature de ses cinq faces, puis son volume. Quelle fraction du cube représente-t-elle ?
\label{fig-ex-16-7}\figureGeo{espace-pyramide-cube}
\end{exercice}\end{blocLivre}
\begin{blocLivre}{ex-16-8}\begin{exercice}[2]{Une section de pavé}
Un pavé de base 8 cm sur 5 cm et de hauteur 6 cm est coupé par un plan parallèle à la base à 2 cm de celle-ci. Quelle forme et quelle aire a la section ?
\label{fig-ex-16-8}\figureGeo{espace-section-pave}
\end{exercice}\end{blocLivre}
\begin{blocLivre}{ex-16-9}\begin{exercice}[2]{Sections d’un cylindre}
Un cylindre a un rayon de 3 cm et une hauteur de 10 cm. Décrire et calculer l’aire de sa section par un plan perpendiculaire à son axe, puis par un plan contenant son axe.
\label{fig-ex-16-9}\figureGeo{espace-section-cylindre}
\end{exercice}\end{blocLivre}
\begin{blocLivre}{ex-16-10}\begin{exercice}[2]{Une section de sphère}
Une sphère de centre O a un rayon de 5 cm. Un plan la coupe à une distance OH = 3 cm du centre ; H est le centre du cercle de section. Calculer le rayon de ce cercle et l’aire du disque de section.
\label{fig-ex-16-10}\figureGeo{espace-section-sphere}
\end{exercice}\end{blocLivre}

\clearpage\rubrique{Problèmes}
\begin{blocLivre}{ex-16-11}\begin{exercice}[3]{Un réservoir assemblé}
Un réservoir fermé est formé d’un cylindre de rayon 1 m et de hauteur 2 m, surmonté d’une demi-boule de même rayon. Calculer sa capacité au litre.
\label{fig-ex-16-11}\figureGeo{espace-reservoir}
\end{exercice}\end{blocLivre}
\begin{blocLivre}{ex-16-12}\begin{exercice}[3]{Une boule creusée}
Une boule pleine de rayon 5 cm contient une cavité sphérique de rayon 2 cm, entièrement intérieure. Quel volume de matière reste-t-il ?
\label{fig-ex-16-12}\figureGeo{espace-cavite}
\end{exercice}\end{blocLivre}
\begin{blocLivre}{ex-16-13}\begin{exercice}[3]{Un cône coupé}
Un cône de rayon 6 cm et de hauteur 12 cm est coupé par un plan parallèle à sa base, à mi-hauteur. Quelle est la section ? Calculer le volume du petit cône de même sommet, puis celui du tronc restant.
\label{fig-ex-16-13}\figureGeo{espace-cone-coupe}
\end{exercice}\end{blocLivre}
\begin{blocLivre}{ex-16-14}\begin{exercice}[3]{Une section de pyramide}
SABCD est une pyramide de base carrée ABCD de côté 6 cm et de hauteur SO = 9 cm. Un plan parallèle à la base coupe la pyramide à 3 cm du sommet S. Déterminer la nature et l’aire de la section, puis le volume de la petite pyramide obtenue.
\label{fig-ex-16-14}\figureGeo{espace-section-pyramide}
\end{exercice}\end{blocLivre}
\begin{blocLivre}{ex-16-15}\begin{exercice}[3]{Le cylindre d’Archimède}
Une boule de rayon 3 cm est posée dans un cylindre de rayon 3 cm et de hauteur 6 cm, qu’elle touche en haut, en bas et sur le côté. Calculer les deux volumes et le rapport du volume de la boule à celui du cylindre. Montrer que ce rapport ne dépend pas du rayon.
\label{fig-ex-16-15}\figureGeo{espace-archimede}
\end{exercice}\end{blocLivre}

\clearpage

\begin{ouverture}[algorithmique]{17}{Pensée informatique et algorithmique}{Partie 5 · Pensée informatique et raisonnement}
{Variables · Conditions composées · Boucles conditionnelles · Blocs personnalisés}
\objectif{Lire, exécuter à la main et compléter un script par blocs.}
\objectif{Utiliser plusieurs variables, des conditions composées et des boucles conditionnelles.}
\objectif{Structurer un programme avec des blocs personnalisés et simuler une expérience.}
\prerequis{Calcul numérique, divisibilité et probabilités.}
\end{ouverture}
\pagestyle{otMain}\essentiel\fontsize{11}{15}\selectfont
\essSection{1}{Variables, conditions et boucles}
\begin{definition}Une variable associe un nom à une valeur, que le programme peut lire et modifier. « Mettre x à 3 » remplace son contenu ; « ajouter 3 à x » augmente sa valeur actuelle. Un programme peut utiliser plusieurs variables, par exemple un compteur et une somme.\end{definition}
\begin{definition}Une condition est un test vrai ou faux, comme $x\ge10$. Une condition composée combine deux conditions : « A et B » est vraie lorsque les deux le sont ; « A ou B » est vraie lorsque l’une au moins l’est.\end{definition}
\begin{definition}La boucle « répéter $n$ fois » exécute son contenu un nombre fixé de fois. La boucle conditionnelle « répéter jusqu’à » recommence tant que sa condition est fausse et s’arrête dès qu’elle devient vraie ; elle ne termine que si le contenu fait évoluer les variables vers cette condition.\end{definition}
\begin{savoirfaire}[Exécuter à la main]Préparer un tableau avec une colonne par variable et une ligne par étape ou par tour. Exécuter dans l’ordre, en utilisant les valeurs déjà modifiées. Tester les valeurs frontières d’une condition.\end{savoirfaire}
\begin{attention}Un test n’est pas une affectation : « x = 6 » dans une condition compare, il ne modifie pas x. Les lignes indentées appartiennent au bloc qui les contient ; une ligne « fin » le ferme.\end{attention}
\clearpage
\essSection{2}{Structurer et simuler}
\begin{definition}Un bloc personnalisé regroupe des instructions sous un nom ; ses paramètres lui fournissent des valeurs. Le bloc polygone(n, c) peut répéter $n$ fois « avancer de c ; tourner de $360/n$ degrés ». On l’appelle ensuite comme une seule instruction.\end{definition}
\begin{savoirfaire}[Décomposer un problème]Isoler une tâche répétée dans un bloc personnalisé, la tester seule, puis écrire le programme principal qui l’appelle. Un programme bien structuré est plus court, plus lisible et plus facile à corriger.\end{savoirfaire}
\begin{savoirfaire}[Simuler]Choisir un entier aléatoire de 1 à 6 modélise un dé équilibré. Pour deux dés indépendants, tirer deux nombres séparément. Initialiser le compteur à zéro avant la boucle, l’augmenter quand l’évènement est réalisé, puis diviser par le nombre d’essais.\end{savoirfaire}
\begin{exemple}Pour obtenir le quotient et le reste de 23 par 5, on retire 5 tant que c’est possible : $23\to18\to13\to8\to3$. Quatre soustractions ont été faites : quotient 4, reste 3. Une boucle conditionnelle automatise ce procédé.\end{exemple}
\begin{attention}Les blocs du livre représentent les instructions d’un langage de type Scratch, sans prétendre reproduire son interface. Les multiplications sont explicites.\end{attention}
\clearpage\rubrique{Découverte}
\begin{blocLivre}{ex-17-1}\begin{exercice}[1]{Suivre trois opérations}
Exécuter le script et donner la valeur affichée.
\label{fig-ex-17-1}\figureGeo{script-01}
\end{exercice}\end{blocLivre}
\begin{blocLivre}{ex-17-2}\begin{exercice}[1]{Un nombre demandé}
Le script lit un nombre x. Donner le résultat pour 2 et pour $-4$, puis l’expression obtenue à partir de x.
\label{fig-ex-17-2}\figureGeo{script-02}
\end{exercice}\end{blocLivre}
\begin{blocLivre}{ex-17-3}\begin{exercice}[1]{Une boucle bornée}
Quelle valeur affiche le script ? Combien de fois le bloc ajouter est-il exécuté ?
\label{fig-ex-17-3}\figureGeo{script-03}
\end{exercice}\end{blocLivre}
\begin{blocLivre}{ex-17-4}\begin{exercice}[1]{Un test et deux branches}
Exécuter le script pour les entrées 9, 10 et 14.
\label{fig-ex-17-4}\figureGeo{script-04}
\end{exercice}\end{blocLivre}
\begin{blocLivre}{ex-17-5}\begin{exercice}[1]{Doubler plusieurs fois}
Quelle est la valeur finale de x ? Un élève propose $2+3\times2$. Expliquer son erreur.
\label{fig-ex-17-5}\figureGeo{script-05}
\end{exercice}\end{blocLivre}

\clearpage\rubrique{Entraînement}
\begin{blocLivre}{ex-17-6}\begin{exercice}[2]{Deux variables}
Suivre les variables s et n à chaque tour, puis donner le résultat. Que calcule ce script ?
\label{fig-ex-17-6}\figureGeo{script-06}
\end{exercice}\end{blocLivre}
\begin{blocLivre}{ex-17-7}\begin{exercice}[2]{Répéter jusqu’à}
Donner les valeurs successives de x et le nombre de tours.
\label{fig-ex-17-7}\figureGeo{script-07}
\end{exercice}\end{blocLivre}
\begin{blocLivre}{ex-17-8}\begin{exercice}[2]{Une condition composée}
Le script teste si un nombre appartient à une plage. Donner l’affichage pour 9, 10, 20 et 21. Que se passerait-il si l’on remplaçait « et » par « ou » ?
\label{fig-ex-17-8}\figureGeo{script-08}
\end{exercice}\end{blocLivre}
\begin{blocLivre}{ex-17-9}\begin{exercice}[2]{Énumérer des diviseurs}
Quelles valeurs sont affichées ? Pourquoi commence-t-on d à 1 et non à 0 ?
\label{fig-ex-17-9}\figureGeo{script-09}
\end{exercice}\end{blocLivre}
\begin{blocLivre}{ex-17-10}\begin{exercice}[2]{Le reste par soustractions}
Le script demande deux entiers positifs a et b. Exécuter le script pour $a=23$ et $b=5$, puis pour $a=3$ et $b=5$. Que calculent les variables q et r ?
\label{fig-ex-17-10}\figureGeo{script-10}
\end{exercice}\end{blocLivre}

\clearpage\rubrique{Problèmes}
\begin{blocLivre}{ex-17-11}\begin{exercice}[3]{Une boucle qui dépasse}
Combien de tours effectue le script ? Remplacer « multiplier par 2 » par « ajouter 0 » permettrait-il encore de terminer ?
\label{fig-ex-17-11}\figureGeo{script-11}
\end{exercice}\end{blocLivre}
\begin{blocLivre}{ex-17-12}\begin{exercice}[3]{Un bloc personnalisé}
Le bloc carré(c) trace un carré et revient au même point avec la même orientation. Expliquer ce que dessine le programme principal. Quelle longueur est tracée au total ?
\label{fig-ex-17-12}\figureGeo{script-carres}
\end{exercice}\end{blocLivre}
\begin{blocLivre}{ex-17-13}\begin{exercice}[3]{Simuler un dé}
Compléter le script pour compter les six sur 600 lancers d’un dé. Où placer l’initialisation et le calcul de fréquence ? Quelle valeur la fréquence affichée devrait-elle approcher ?
\label{fig-ex-17-13}\figureGeo{script-simulation}
\end{exercice}\end{blocLivre}
\begin{blocLivre}{ex-17-14}\begin{exercice}[3]{Deux dés et un « ou »}
Le script simule 1000 lancers de deux dés. Quel évènement compte-t-il ? Calculer sa probabilité. Parmi les affichages 0,17 ; 0,31 et 0,50, lequel est plausible ?
\label{fig-ex-17-14}\figureGeo{script-deux-des}
\end{exercice}\end{blocLivre}
\begin{blocLivre}{ex-17-15}\begin{exercice}[3]{Des polygones réguliers}
Le bloc polygone(n, c) trace un polygone régulier à $n$ côtés de longueur c. Le programme principal l’appelle avec n allant de 3 à 6. Nommer les figures tracées, donner l’angle de rotation pour l’hexagone et la longueur totale tracée.
\label{fig-ex-17-15}\figureGeo{script-polygones}
\end{exercice}\end{blocLivre}

\clearpage

\begin{ouverture}[algorithmique]{18}{Raisonner et préparer le brevet}{Partie 5 · Pensée informatique et raisonnement}
{Conjecturer · Démontrer · Réfuter · Rédiger · Vérifier}
\objectif{Distinguer exemple, conjecture, preuve et contre-exemple.}
\objectif{Utiliser réciproque, contraposée et raisonnement par l’absurde.}
\objectif{Organiser le temps et la rédaction de l’épreuve du brevet.}
\prerequis{Ensemble des chapitres précédents.}
\end{ouverture}
\pagestyle{otMain}\essentiel\fontsize{11}{15}\selectfont
\essSection{1}{Les formes du raisonnement}
\begin{definition}Une conjecture est une affirmation que l’on pense vraie après des observations. Une preuve établit une conclusion à partir d’hypothèses, de définitions et de propriétés. Un seul contre-exemple suffit à réfuter une affirmation générale ; plusieurs exemples ne suffisent pas à la démontrer.\end{definition}
\begin{definition}Pour une propriété « si $P$, alors $Q$ », la réciproque est « si $Q$, alors $P$ » et la contraposée « si $Q$ est faux, alors $P$ est faux ». La contraposée d’une propriété vraie est vraie ; sa réciproque peut être fausse.\end{definition}
\begin{savoirfaire}[Raisonner par l’absurde]Supposer le contraire de ce que l’on veut prouver, en déduire une contradiction, puis conclure que la supposition était fausse. Exemple : un triangle ne peut pas avoir deux angles droits, sinon la somme de ses angles dépasserait $180^\circ$.\end{savoirfaire}
\begin{savoirfaire}[Démontrer]Pour une propriété des entiers, traduire « pair » par $2n$, « impair » par $2n+1$, « multiple de $d$ » par $dn$. En géométrie, distinguer une définition, une propriété et une propriété caractéristique, puis citer les hypothèses avant le théorème.\end{savoirfaire}
\begin{attention}Une mesure sur le dessin, une image de calculatrice ou un essai de programme donne une observation ; cela ne justifie pas à soi seul une propriété pour tous les cas.\end{attention}
\clearpage
\essSection{2}{Une copie claire dans le temps imparti}
\begin{savoirfaire}[L’épreuve retenue]D’après la définition en vigueur à compter de la session 2027, retenue pour 2029, l’épreuve dure 2 h. Partie 1 : 20 min, sans calculatrice, 6 points d’automatismes et d’applications directes, à réponses courtes. Partie 2 : 1 h 40, calculatrice autorisée, 14 points de raisonnement et de résolution de problèmes. La première copie est relevée après 20 minutes.\end{savoirfaire}
\begin{propriete}La seconde partie ne comporte pas de QCM. Toutes les réponses doivent être justifiées, sauf indication contraire. La maîtrise de la langue est évaluée sur 2 points ; les essais et démarches engagés, même non aboutis, sont pris en compte.\end{propriete}
\begin{savoirfaire}[Rédiger]Écrire le calcul ou la propriété utile, donner les étapes essentielles, indiquer unités et arrondis, puis répondre par une phrase correctement construite. Si un résultat intermédiaire manque, poursuivre avec une valeur annoncée ou une démarche explicitée.\end{savoirfaire}
\begin{savoirfaire}[Vérifier]Relire les nombres copiés, le signe, la précision, les unités, les domaines et les ordres de grandeur. Une probabilité ne dépasse pas 1 ; une aire reste positive ; un achat entier se majore pour couvrir un besoin.\end{savoirfaire}
\clearpage\rubrique{Découverte}
\begin{blocLivre}{ex-18-1}\begin{exercice}[1]{Un contre-exemple}
« Le carré d’un nombre est toujours plus grand que ce nombre. » Réfuter cette phrase.
\end{exercice}\end{blocLivre}
\begin{blocLivre}{ex-18-2}\begin{exercice}[1]{Deux entiers impairs}
Démontrer que la somme de deux entiers impairs est paire.
\end{exercice}\end{blocLivre}
\begin{blocLivre}{ex-18-3}\begin{exercice}[1]{Réciproque et contraposée}
On considère la propriété vraie : « Si un entier est multiple de 6, alors il est pair. » Énoncer sa réciproque et sa contraposée. Lesquelles sont vraies ?
\end{exercice}\end{blocLivre}
\begin{blocLivre}{ex-18-4}\begin{exercice}[1]{Un arrondi utile}
Il faut 7,2 L d’un produit vendu en bidons de 3 L. Combien de bidons au minimum ?
\end{exercice}\end{blocLivre}
\begin{blocLivre}{ex-18-5}\begin{exercice}[1]{Une hypothèse manquante}
Un triangle est dessiné avec deux côtés de 3 cm et 4 cm. Peut-on affirmer que le troisième vaut 5 cm ?
\end{exercice}\end{blocLivre}

\clearpage\rubrique{Entraînement}
\begin{blocLivre}{ex-18-6}\begin{exercice}[2]{Un programme constant}
Choisir x, ajouter 5, doubler, puis retirer deux fois le nombre de départ. Montrer que le résultat est toujours le même.
\end{exercice}\end{blocLivre}
\begin{blocLivre}{ex-18-7}\begin{exercice}[2]{Un raisonnement par l’absurde}
Démontrer par l’absurde qu’un triangle rectangle ne peut pas avoir un angle obtus.
\end{exercice}\end{blocLivre}
\begin{blocLivre}{ex-18-8}\begin{exercice}[2]{Une conjecture de divisibilité}
On observe que $n(n+1)$ est pair pour plusieurs entiers n. Le démontrer.
\end{exercice}\end{blocLivre}
\begin{blocLivre}{ex-18-9}\begin{exercice}[2]{Vérifier une équation}
Un élève affirme que $x=4$ résout $3(x-2)=x+4$. Vérifier, puis résoudre.
\end{exercice}\end{blocLivre}
\begin{blocLivre}{ex-18-10}\begin{exercice}[2]{Une moyenne mal interprétée}
Une classe a une moyenne de 12/20. Est-il certain qu’au moins la moitié des notes sont supérieures ou égales à 12 ?
\end{exercice}\end{blocLivre}

\clearpage\rubrique{Problèmes}
\begin{blocLivre}{ex-18-11}\begin{exercice}[3]{Une commande au meilleur prix}
Pour 24 carnets : A vend à 2,80 € l’unité avec 6 € de livraison ; B vend des lots de 8 à 21 €, livraison gratuite. Comparer les coûts.
\end{exercice}\end{blocLivre}
\begin{blocLivre}{ex-18-12}\begin{exercice}[3]{Un modèle de bassin}
Un bassin pavé de 4 m sur 2,5 m contient 6000 L. La hauteur d’eau maximale est 0,8 m. Calculer la hauteur actuelle puis la quantité encore possible.
\label{fig-ex-18-12}\figureGeo{synthese-bassin}
\end{exercice}\end{blocLivre}
\begin{blocLivre}{ex-18-13}\begin{exercice}[3]{Un hasard et un budget}
Une tombola propose 200 billets à 3 €. Vingt billets gagnent chacun un lot coûtant 15 € à l’organisateur. Tous les billets sont vendus et chaque billet a la même chance. Calculer la probabilité qu’un billet donné gagne et le solde de l’organisation, sans autres frais.
\end{exercice}\end{blocLivre}
\begin{blocLivre}{ex-18-14}\begin{exercice}[3]{Un maximum expliqué}
Un rectangle a un périmètre de 40 cm. Montrer que son aire vaut $100-(x-10)^2$ si un côté vaut x, puis déterminer la plus grande aire.
\label{fig-ex-18-14}\figureGeo{rectangle-fonction}
\end{exercice}\end{blocLivre}
\begin{blocLivre}{ex-18-15}\begin{exercice}[3]{Un plan de révision}
Une élève réussit 8 automatismes sur 12, puis obtient 9/14 aux problèmes, rédaction comprise. Selon le barème des brevets blancs du livre, quelle note obtient-elle ? Quel travail cibler si ses quatre erreurs concernent les fractions ?
\end{exercice}\end{blocLivre}

\clearpage

\editionPage[\pdfbookmark[0]{Défis d’approfondissement}{bm-defis-debut-parent}]{defis-debut}{Dix-huit défis}{Un défi par chapitre · Parcours facultatif}{Chercher avant de consulter les deux indications.}
\begin{blocEdition}{defis-debut}
Ces problèmes demandent davantage d’initiative. Ils utilisent les outils du livre ; les prolongements sont définis dans l’énoncé. Aucun défi n’ajoute une exigence officielle au brevet.\par\bigskip
\begin{savoirfaire}[Une recherche qui avance]Faire un schéma, tester un petit cas, comparer deux méthodes, conserver les essais et préciser ce qui reste à prouver. Ouvrir la première indication seulement après une recherche réelle ; revenir au problème avant de lire la seconde.\end{savoirfaire}
\begin{savoirfaire}[Après la solution]Identifier l’idée décisive et la refaire sur une feuille vierge. Expliquer pourquoi elle fonctionne compte autant que retrouver le résultat.\end{savoirfaire}
\end{blocEdition}
\editionPage{defi-ex-1}{D01 · Défi}{Une fraction entre deux autres}{Approfondissement · 30 à 60 minutes indicatives}
\begin{blocEdition}{defi-ex-1}
Trouver une fraction de dénominateur 60 strictement comprise entre $5/12$ et $7/15$. Combien y en a-t-il ? Puis trouver une fraction comprise entre ces deux nombres, mais de dénominateur 120.
\end{blocEdition}
\editionPage{defi-ex-2}{D02 · Défi}{Comparer sans grands calculs}{Approfondissement · 30 à 60 minutes indicatives}
\begin{blocEdition}{defi-ex-2}
Comparer $2^{30}$ et $10^9$ sans développer les deux nombres. On pourra utiliser $2^{10}=1024$. Puis dire si $2^{60}$ est plus grand que $10^{18}$.
\end{blocEdition}
\editionPage{defi-ex-3}{D03 · Défi}{Une recherche entière exhaustive}{Approfondissement · 30 à 60 minutes indicatives}
\begin{blocEdition}{defi-ex-3}
Trouver tous les entiers positifs n tels que n divise 180 et que n soit un carré parfait. Justifier que la liste est complète.
\end{blocEdition}
\editionPage{defi-ex-4}{D04 · Défi}{L’identité de Sophie Germain}{Approfondissement · 30 à 60 minutes indicatives}
\begin{blocEdition}{defi-ex-4}
Vérifier que pour tous nombres $a$ et $b$ : $a^4+4b^4=(a^2+2b^2)^2-(2ab)^2$. En déduire une factorisation de $n^4+4$, puis démontrer que $n^4+4$ n’est jamais un nombre premier lorsque $n$ est un entier supérieur ou égal à 2.
\end{blocEdition}
\editionPage{defi-ex-5}{D05 · Défi}{Deux solutions et un contexte}{Approfondissement · 30 à 60 minutes indicatives}
\begin{blocEdition}{defi-ex-5}
Un rectangle a pour côtés x et $x+3$ cm. Un carré a pour côté $x+1$ cm. Pour quelles valeurs positives de x leurs aires diffèrent-elles de 1 cm$^2$, sans préciser laquelle est la plus grande ?
\end{blocEdition}
\editionPage{defi-ex-6}{D06 · Défi}{Promotions dans quel ordre ?}{Approfondissement · 30 à 60 minutes indicatives}
\begin{blocEdition}{defi-ex-6}
Un article coûte P euros, avec $P>20$. Le magasin applique une réduction de 20 \% et un bon de 10 €. Comparer les deux ordres possibles et montrer que la différence de prix ne dépend pas de P.
\end{blocEdition}
\editionPage{defi-ex-7}{D07 · Défi}{Une rencontre de débits}{Approfondissement · 30 à 60 minutes indicatives}
\begin{blocEdition}{defi-ex-7}
Une cuve vide reçoit 18 L/min et perd simultanément 3 L/min par une fuite constante. Après 12 min, on coupe l’arrivée d’eau ; la fuite continue seule. Au bout de combien de temps depuis le début la cuve redevient-elle vide ?
\end{blocEdition}
\editionPage{defi-ex-8}{D08 · Défi}{Une fonction cachée derrière des essais}{Approfondissement · 30 à 60 minutes indicatives}
\begin{blocEdition}{defi-ex-8}
On sait seulement que $f(-1)=3$, $f(0)=2$, $f(1)=3$. Un élève conclut que $f(x)=x^2+2$. Montrer que cette conclusion n’est pas imposée en proposant une autre formule donnant les mêmes trois images mais une image différente pour 2.
\end{blocEdition}
\editionPage{defi-ex-9}{D09 · Défi}{Trois tarifs, un choix argumenté}{Approfondissement · 30 à 60 minutes indicatives}
\begin{blocEdition}{defi-ex-9}
Pour n séances entières, de 1 à 20, trois offres coûtent A = 8n, B = $30+5n$, C = $70+2n$, en euros. Déterminer l’offre la moins chère pour chaque nombre de séances ; indiquer les égalités utiles.
\end{blocEdition}
\editionPage{defi-ex-10}{D10 · Défi}{Une série à reconstruire}{Approfondissement · 30 à 60 minutes indicatives}
\begin{blocEdition}{defi-ex-10}
Proposer une série de huit nombres entiers dont le minimum est 2, le premier quartile 5, la médiane 8, le troisième quartile 12, le maximum 20 et la moyenne 9. Une telle série peut-elle avoir une moyenne égale à 15 ?
\end{blocEdition}
\editionPage{defi-ex-11}{D11 · Défi}{Un jeu de somme équitable ?}{Approfondissement · 30 à 60 minutes indicatives}
\begin{blocEdition}{defi-ex-11}
Deux dés équilibrés sont lancés indépendamment. Le joueur A gagne si la somme est 6 ou 8 ; B gagne si elle est 7. Sinon on recommence sans modifier les règles. Les deux joueurs ont-ils la même chance de gagner lors d’un lancer décisif ? Comparer en comptant les couples, sans nouveau formalisme.
\end{blocEdition}
\editionPage{defi-ex-12}{D12 · Défi}{Une corde dans un cercle}{Approfondissement · 30 à 60 minutes indicatives}
\begin{blocEdition}{defi-ex-12}
Un cercle a pour centre O et rayon 13 cm. Une corde AB est à 5 cm de O. H est le pied de la perpendiculaire issue de O sur AB. On admet que H est le milieu de la corde. Calculer AB et l’aire de OAB.\label{fig-defi-12}\figureGeo{defi-corde}

\end{blocEdition}
\editionPage{defi-ex-13}{D13 · Défi}{Les diagonales d’un trapèze}{Approfondissement · 30 à 60 minutes indicatives}
\begin{blocEdition}{defi-ex-13}
ABCD est un trapèze de bases [AB] et [CD], parallèles, avec AB = 6 cm et CD = 10 cm. Ses diagonales [AC] et [BD] se coupent en O. Calculer $OA/OC$. Si AC = 12 cm, calculer OA et OC.\label{fig-defi-13}\figureGeo{defi-trapeze}

\end{blocEdition}
\editionPage{defi-ex-14}{D14 · Défi}{Deux angles de visée}{Approfondissement · 30 à 60 minutes indicatives}
\begin{blocEdition}{defi-ex-14}
Deux points d’observation A et B sont à 10 m l’un de l’autre sur une même droite horizontale allant vers le pied d’une tour. B est le plus proche. Les angles d’élévation du sommet sont de 30 degrés en A et 45 degrés en B. Les visées se font au niveau du sol dans ce modèle. Calculer la hauteur au décimètre.\label{fig-defi-14}\figureGeo{defi-visees}

\end{blocEdition}
\editionPage{defi-ex-15}{D15 · Défi}{Une somme de deux vecteurs}{Approfondissement · 30 à 60 minutes indicatives}
\begin{blocEdition}{defi-ex-15}
ABC est un triangle. On construit le point D tel que $\vect{AD}=\vect{AB}+\vect{AC}$, puis le point E tel que $\vect{BE}=\vect{AC}$. Démontrer que D et E sont confondus. Quelle est la nature de ABDC ? Si I est le milieu de [BC], montrer que $\vect{AB}+\vect{AC}=2\vect{AI}$.\label{fig-defi-15}\figureGeo{defi-vecteurs}

\end{blocEdition}
\editionPage{defi-ex-16}{D16 · Défi}{Une section d’aire imposée}{Approfondissement · 30 à 60 minutes indicatives}
\begin{blocEdition}{defi-ex-16}
Une boule de rayon 10 cm est coupée par un plan. Le disque de section doit avoir pour aire $64\pi$ cm$^2$. À quelle distance du centre placer le plan ? Y a-t-il plusieurs positions parallèles possibles ?\label{fig-defi-16}\figureGeo{defi-section}

\end{blocEdition}
\editionPage{defi-ex-17}{D17 · Défi}{Un script de recherche et sa preuve}{Approfondissement · 30 à 60 minutes indicatives}
\begin{blocEdition}{defi-ex-17}
Un script initialise s à 0, n à 1, puis répète tant que s est inférieur à 100 : ajouter n à s ; ajouter 1 à n. Quelle valeur de n possède-t-il à la sortie ? Écrire un script par blocs ou en français structuré et justifier l’arrêt.
\end{blocEdition}
\editionPage{defi-ex-18}{D18 · Défi}{Le meilleur rectangle entier}{Approfondissement · 30 à 60 minutes indicatives}
\begin{blocEdition}{defi-ex-18}
Un rectangle a des côtés entiers positifs et un périmètre de 46 cm. Déterminer les dimensions donnant l’aire maximale et prouver qu’aucun autre rectangle entier ne fait mieux.
\end{blocEdition}

\editionPage[\pdfbookmark[0]{Indications des défis}{bm-indices-debut-parent}]{indices-debut}{Indications progressives}{Défis 1 à 3}{Ne lire que l’aide nécessaire.}
\begin{blocEdition}{indices-debut}
\editionEtape{D01 · Une fraction entre deux autres}\textbf{Indice 1.} Mettre les deux bornes au dénominateur 60.\par\medskip\textbf{Indice 2.} Chercher les numérateurs entiers strictement entre les deux bornes.\par
\editionEtape{D02 · Comparer sans grands calculs}\textbf{Indice 1.} Écrire $2^{30}$ comme un cube.\par\medskip\textbf{Indice 2.} Élever au carré deux nombres positifs conserve leur ordre.\par
\editionEtape{D03 · Une recherche entière exhaustive}\textbf{Indice 1.} Décomposer 180 en facteurs premiers.\par\medskip\textbf{Indice 2.} Dans un carré, les exposants premiers sont pairs.\par

\end{blocEdition}
\editionPage{indices-4}{Indications progressives}{Défis 4 à 6}{Ne lire que l’aide nécessaire.}
\begin{blocEdition}{indices-4}
\editionEtape{D04 · L’identité de Sophie Germain}\textbf{Indice 1.} Développer $(a^2+2b^2)^2$ avec une identité remarquable.\par\medskip\textbf{Indice 2.} Factoriser une différence de deux carrés, puis montrer que chaque facteur vaut au moins 2.\par
\editionEtape{D05 · Deux solutions et un contexte}\textbf{Indice 1.} Calculer aire du carré moins aire du rectangle.\par\medskip\textbf{Indice 2.} Une différence de valeur absolue 1 correspond à deux égalités possibles.\par
\editionEtape{D06 · Promotions dans quel ordre ?}\textbf{Indice 1.} Écrire les deux prix avec P.\par\medskip\textbf{Indice 2.} Développer $0{,}8(P-10)$.\par

\end{blocEdition}
\editionPage{indices-7}{Indications progressives}{Défis 7 à 9}{Ne lire que l’aide nécessaire.}
\begin{blocEdition}{indices-7}
\editionEtape{D07 · Une rencontre de débits}\textbf{Indice 1.} Calculer le débit net pendant le remplissage.\par\medskip\textbf{Indice 2.} Distinguer la durée de vidange de la durée totale.\par
\editionEtape{D08 · Une fonction cachée derrière des essais}\textbf{Indice 1.} Ajouter une expression nulle en -1, 0 et 1.\par\medskip\textbf{Indice 2.} Le produit $(x+1)x(x-1)$ est nul pour les trois entrées.\par
\editionEtape{D09 · Trois tarifs, un choix argumenté}\textbf{Indice 1.} Comparer A et B, B et C, puis A et C.\par\medskip\textbf{Indice 2.} Tester les entiers proches des points d’égalité.\par

\end{blocEdition}
\editionPage{indices-10}{Indications progressives}{Défis 10 à 12}{Ne lire que l’aide nécessaire.}
\begin{blocEdition}{indices-10}
\editionEtape{D10 · Une série à reconstruire}\textbf{Indice 1.} Ranger les huit valeurs $v_1\le\dots\le v_8$ : $Q_1$ est la 2\textsuperscript{e}, $Q_3$ la 6\textsuperscript{e}, et la médiane est la moyenne des 4\textsuperscript{e} et 5\textsuperscript{e}.\par\medskip\textbf{Indice 2.} La somme doit valoir 72. Pour la seconde question, chercher la plus grande somme possible.\par
\editionEtape{D11 · Un jeu de somme équitable ?}\textbf{Indice 1.} Compter les couples pour chacune des trois sommes.\par\medskip\textbf{Indice 2.} Se limiter aux couples qui désignent un gagnant.\par
\editionEtape{D12 · Une corde dans un cercle}\textbf{Indice 1.} Appliquer Pythagore à OHA.\par\medskip\textbf{Indice 2.} La base AB a pour hauteur OH.\par

\end{blocEdition}
\editionPage{indices-13}{Indications progressives}{Défis 13 à 15}{Ne lire que l’aide nécessaire.}
\begin{blocEdition}{indices-13}
\editionEtape{D13 · Les diagonales d’un trapèze}\textbf{Indice 1.} Repérer une configuration en papillon de sommet O.\par\medskip\textbf{Indice 2.} Les droites (AB) et (CD) sont parallèles : appliquer le théorème de Thalès dans les triangles OAB et OCD.\par
\editionEtape{D14 · Deux angles de visée}\textbf{Indice 1.} Si la distance B-pied vaut x, la visée à 45 degrés donne une hauteur x.\par\medskip\textbf{Indice 2.} L’autre visée donne $x=(x+10)\tan30^\circ$.\par
\editionEtape{D15 · Une somme de deux vecteurs}\textbf{Indice 1.} Écrire $\vect{AE}$ avec la relation de Chasles en passant par B.\par\medskip\textbf{Indice 2.} Les diagonales d’un parallélogramme ont le même milieu.\par

\end{blocEdition}
\editionPage{indices-16}{Indications progressives}{Défis 16 à 18}{Ne lire que l’aide nécessaire.}
\begin{blocEdition}{indices-16}
\editionEtape{D16 · Une section d’aire imposée}\textbf{Indice 1.} Retrouver le rayon de section grâce à son aire.\par\medskip\textbf{Indice 2.} Dans le triangle rectangle des rayons, utiliser $10^2=OH^2+8^2$.\par
\editionEtape{D17 · Un script de recherche et sa preuve}\textbf{Indice 1.} Calculer les sommes successives $1+2+\cdots+k$.\par\medskip\textbf{Indice 2.} La variable n est augmentée après son ajout à s.\par
\editionEtape{D18 · Le meilleur rectangle entier}\textbf{Indice 1.} La somme des deux côtés vaut 23.\par\medskip\textbf{Indice 2.} Comparer $x(23-x)$ à $132$ ou écrire $(x-11)(x-12)$.\par

\end{blocEdition}

\editionPage[\pdfbookmark[0]{Deux brevets blancs}{bm-bilans-guide-parent}]{bilans-guide}{Deux brevets blancs}{Session 2029 · Sujets originaux}{Durée totale : 2 heures · Note sur 20 · Coefficient 2}
\begin{blocEdition}{bilans-guide}

Définition retenue : note de service du 11 septembre 2026 (BO spécial n° 4 du
17 septembre 2026), en vigueur à compter de la session 2027. Au 23 octobre 2026,
aucun texte propre à la session 2029 n’était publié ; vérifier avant l’examen
qu’aucune modification n’est intervenue.
\editionEtape{Partie 1 : 20 minutes, sans calculatrice}
Automatismes et applications directes : 6 points. Questions à réponses courtes,
portant sur les automatismes du cycle 4 ou l’application directe d’une propriété,
d’une formule ou d’un théorème. Ici, douze questions à 0,5 point. La première
copie est relevée au bout de 20 minutes.
\editionEtape{Partie 2 : 1 h 40, calculatrice autorisée}
Raisonnement et résolution de problèmes : 14 points. Exercices indépendants,
sans QCM ; toutes les réponses sont justifiées, sauf indication contraire. La
maîtrise de la langue est évaluée sur 2 points. Utiliser une seconde copie ;
le passage de la première à la seconde partie se fait sans pause.
\editionEtape{Le barème de ce livre}
Chaque problème vaut 3,5 points : trois questions à 1 point et 0,5 point pour
la maîtrise de la langue et la clarté de la rédaction. Les 2 points de maîtrise
de la langue sont ainsi inclus dans les 14 points. Ce partage est un choix
d’entraînement ; il ne préjuge pas des barèmes des sujets officiels.
\editionEtape{Après la séance}
Comparer les démarches, classer les erreurs par notion et refaire les questions
utiles sans correction. Les essais et démarches engagés, même non aboutis, sont
valorisés : un résultat numérique faux n’efface pas une démarche correcte.

\end{blocEdition}
\editionPage{bilan-A-auto}{Brevet blanc A}{Relier calculs, modèles et figures}{Partie 1 · 20 min · Sans calculatrice · 6 points}
\begin{blocEdition}{bilan-A-auto}
\begin{questions}
\item Calculer $-7+3\times5$.
\item Calculer $\dfrac23+\dfrac16$.
\item Donner la notation scientifique de 450000.
\item Résoudre $x^2=81$.
\item Développer $(x+3)^2$.
\item Résoudre l’inéquation $-2x\ge8$.
\item Calculer 15 \% de 80.
\item Pour $f(x)=2x+1$, calculer $f(-2)$.
\item Partager 35 € selon le ratio $2:5$.
\item Donner la médiane de $2\,;4\,;5\,;8\,;11$.
\item $P(A)=0{,}5$, $P(B)=0{,}4$ et $P(A\cap B)=0{,}2$. Calculer $P(A\cup B)$.
\item ABCD est un parallélogramme. Compléter : $\vect{AB}=\vect{D\ldots}$.\end{questions}
\end{blocEdition}
\editionPage{bilan-A-ex-1}{Brevet A · Problème 1}{Un atelier de carnets}{Partie 2 · 1 h 40 pour les quatre problèmes · 3,5 points}
\begin{blocEdition}{bilan-A-ex-1}
Un atelier possède 168 feuilles bleues et 126 feuilles blanches. Il prépare le plus grand nombre possible de carnets identiques, sans reste. Chaque carnet contient le même nombre de feuilles de chaque couleur.\begin{questions}
\item Déterminer le nombre maximal de carnets et leur composition.
\item L’atelier vend chaque carnet 4,50 €. Les fournitures ont coûté 78 €. Calculer le solde si tous les carnets sont vendus.
\item Un autre atelier vend le carnet 5 €, puis applique 10 \% de réduction. Comparer son prix unitaire au précédent.
\end{questions}\par\medskip{\small\color{otMuted}Barème : question 1 : 1 pt ; question 2 : 1 pt ; question 3 : 1 pt ; maîtrise de la langue et clarté de la rédaction : 0,5 pt.}
\end{blocEdition}
\editionPage{bilan-A-ex-2}{Brevet A · Problème 2}{Une voile divisée}{Partie 2 · 1 h 40 pour les quatre problèmes · 3,5 points}
\begin{blocEdition}{bilan-A-ex-2}
ABC est rectangle en A, AB = 12 m, AC = 9 m. M appartient à [AB] avec AM = 4 m. N appartient à [AC] et (MN) est parallèle à (BC).\begin{questions}
\item Calculer BC.
\item Calculer AN et MN.
\item On construit le point D, image de N par la translation de vecteur $\vect{MB}$. Quelle est la nature du quadrilatère MBDN ? Démontrer que D appartient à la droite (BC).
\end{questions}\par\medskip{\small\color{otMuted}Barème : question 1 : 1 pt ; question 2 : 1 pt ; question 3 : 1 pt ; maîtrise de la langue et clarté de la rédaction : 0,5 pt.}\label{fig-bilan-A-ex-2}\figureGeo{brevet-A-geo}

\end{blocEdition}
\editionPage{bilan-A-ex-3}{Brevet A · Problème 3}{Deux locations de matériel}{Partie 2 · 1 h 40 pour les quatre problèmes · 3,5 points}
\begin{blocEdition}{bilan-A-ex-3}
Pour $n$ journées entières de location, A facture $18+7n$ euros et B facture $10n$ euros, pour $1\le n\le15$.\begin{questions}
\item Calculer les deux prix pour 4 jours. Quelle offre choisir ?
\item Pour combien de jours les prix sont-ils égaux ? Résoudre l’inéquation $18+7n<10n$ et interpréter le résultat.
\item Avec un budget de 90 €, déterminer la durée maximale possible chez chaque loueur, en résolvant une inéquation.
\end{questions}\par\medskip{\small\color{otMuted}Barème : question 1 : 1 pt ; question 2 : 1 pt ; question 3 : 1 pt ; maîtrise de la langue et clarté de la rédaction : 0,5 pt.}\label{fig-bilan-A-ex-3}\figureGeo{brevet-A-tarifs}

\end{blocEdition}
\editionPage{bilan-A-ex-4}{Brevet A · Problème 4}{Des scores et un tirage}{Partie 2 · 1 h 40 pour les quatre problèmes · 3,5 points}
\begin{blocEdition}{bilan-A-ex-4}
Une série de 20 essais donne les scores du tableau. On tire ensuite deux boules avec remise dans un sac de 3 boules rouges et 1 bleue, numérotées $R_1$, $R_2$, $R_3$ et $B$ ; les tirages sont faits au hasard.\begin{center}\renewcommand{\arraystretch}{1.35}\begin{tabular}{c|cccc}\toprule
Score & $0$ & $1$ & $2$ & $3$\\\midrule
Effectif & $2$ & $5$ & $8$ & $5$\\
\bottomrule\end{tabular}\end{center}\begin{questions}
\item Calculer la moyenne, la médiane et les quartiles des scores.
\item À l’aide du tableau des couples, calculer la probabilité d’obtenir deux boules rouges.
\item Calculer la probabilité d’obtenir exactement une boule bleue, puis celle d’obtenir au moins une boule rouge.
\end{questions}\par\medskip{\small\color{otMuted}Barème : question 1 : 1 pt ; question 2 : 1 pt ; question 3 : 1 pt ; maîtrise de la langue et clarté de la rédaction : 0,5 pt.}\label{fig-bilan-A-ex-4}\figureGeo{brevet-A-tableau}

\end{blocEdition}
\editionPage{bilan-B-auto}{Brevet blanc B}{Choisir une méthode et contrôler}{Partie 1 · 20 min · Sans calculatrice · 6 points}
\begin{blocEdition}{bilan-B-auto}
\begin{questions}
\item Calculer $18-4\times3$.
\item Calculer $\dfrac35\div\dfrac9{10}$.
\item Écrire $10^3\times10^{-5}$ comme une puissance de 10.
\item Décomposer 60 en produit de facteurs premiers.
\item Factoriser $x^2-49$.
\item Résoudre $(x-3)(2x+1)=0$.
\item Augmenter 50 de 12 \%.
\item Pour $g(x)=-3x$, donner l’antécédent de 12.
\item Donner le volume exact d’une boule de rayon 3 cm.
\item Donner le premier quartile de $1\,;3\,;4\,;6\,;8\,;9\,;12\,;15$.
\item Un évènement a une probabilité de 0,3. Donner celle de son contraire.
\item Simplifier $\vect{AB}+\vect{BC}$.\end{questions}
\end{blocEdition}
\editionPage{bilan-B-ex-1}{Brevet B · Problème 1}{Un abri et sa pente}{Partie 2 · 1 h 40 pour les quatre problèmes · 3,5 points}
\begin{blocEdition}{bilan-B-ex-1}
Un toit est représenté par ABC rectangle en A. AB = 4 m horizontalement, AC = 2,5 m verticalement et BC est le versant. L’abri a une longueur de 6 m perpendiculaire au dessin.\begin{questions}
\item Calculer BC au centimètre.
\item Calculer l’angle $\widehat B$ au degré.
\item Le versant rectangulaire doit être couvert. Ajouter 8 \% à son aire pour les découpes ; chaque plaque couvre 2 m$^2$. Combien de plaques au minimum ? Utiliser la longueur BC non arrondie.
\end{questions}\par\medskip{\small\color{otMuted}Barème : question 1 : 1 pt ; question 2 : 1 pt ; question 3 : 1 pt ; maîtrise de la langue et clarté de la rédaction : 0,5 pt.}\label{fig-bilan-B-ex-1}\figureGeo{brevet-B-geo}

\end{blocEdition}
\editionPage{bilan-B-ex-2}{Brevet B · Problème 2}{Une bouteille et une boule}{Partie 2 · 1 h 40 pour les quatre problèmes · 3,5 points}
\begin{blocEdition}{bilan-B-ex-2}
Une bouteille est modélisée par un cylindre de rayon intérieur 4 cm et de hauteur utile 20 cm. Une machine la remplit à 50 mL/s, à débit constant.\begin{questions}
\item Calculer la capacité exacte en cm$^3$, puis l’arrondir au mL.
\item Calculer la durée nécessaire pour remplir 90 \% de la capacité, au dixième de seconde.
\item On verse le contenu d’une bouteille pleine dans un récipient en forme de boule de rayon intérieur 7 cm. Le récipient est-il assez grand ? Quel pourcentage de son volume est rempli, à l’unité près ?
\end{questions}\par\medskip{\small\color{otMuted}Barème : question 1 : 1 pt ; question 2 : 1 pt ; question 3 : 1 pt ; maîtrise de la langue et clarté de la rédaction : 0,5 pt.}\label{fig-bilan-B-ex-2}\figureGeo{brevet-B-cylindre}

\end{blocEdition}
\editionPage{bilan-B-ex-3}{Brevet B · Problème 3}{Un script de figures}{Partie 2 · 1 h 40 pour les quatre problèmes · 3,5 points}
\begin{blocEdition}{bilan-B-ex-3}
Le projet démarre stylo baissé, orienté vers la droite. Le bloc carré(c) trace un carré de côté c et revient au départ avec l’orientation initiale. Le programme principal est donné ci-dessous.\begin{questions}
\item Donner les côtés successifs des carrés tracés et le nombre de passages dans la boucle.
\item Calculer la somme des périmètres tracés.
\item Donner la valeur affichée à la fin et la distance totale parcourue stylo relevé. Expliquer pourquoi cette valeur finale de c n’est pas le côté d’un carré tracé.
\end{questions}\par\medskip{\small\color{otMuted}Barème : question 1 : 1 pt ; question 2 : 1 pt ; question 3 : 1 pt ; maîtrise de la langue et clarté de la rédaction : 0,5 pt.}\label{fig-bilan-B-ex-3}\figureGeo{brevet-B-script}

\end{blocEdition}
\editionPage{bilan-B-ex-4}{Brevet B · Problème 4}{Un jeu à deux étapes}{Partie 2 · 1 h 40 pour les quatre problèmes · 3,5 points}
\begin{blocEdition}{bilan-B-ex-4}
Une roue possède quatre secteurs égaux : trois portent la lettre A, un porte la lettre B. Un sac contient une boule rouge et une bleue. On tourne la roue, puis on tire une boule au hasard, indépendamment de la roue. Le joueur gagne s’il obtient A et rouge, ou B et bleue. Les secteurs A sont numérotés A1, A2, A3 dans le tableau.\begin{questions}
\item Combien d’issues équiprobables compte l’expérience ? Calculer la probabilité de gagner.
\item Sur 100 parties, peut-on garantir exactement 50 gains ? Expliquer.
\item On note E : « la roue donne A » et F : « la boule est rouge ». Calculer $P(E)$, $P(F)$, $P(E\cap F)$, puis $P(E\cup F)$.
\end{questions}\par\medskip{\small\color{otMuted}Barème : question 1 : 1 pt ; question 2 : 1 pt ; question 3 : 1 pt ; maîtrise de la langue et clarté de la rédaction : 0,5 pt.}\label{fig-bilan-B-ex-4}\figureGeo{brevet-B-tableau}

\end{blocEdition}

\editionPage[\pdfbookmark[0]{Démonstrations}{bm-preuves-debut-parent}]{preuves-debut}{Comprendre les preuves}{Identités remarquables et parité}{Les arguments ci-dessous éclairent des propriétés utilisées dans le livre.}
\begin{blocEdition}{preuves-debut}

\begin{demonstration}[Les trois identités remarquables]
Par double distributivité, $(a+b)^2=(a+b)(a+b)=a^2+ab+ba+b^2=a^2+2ab+b^2$.
De même, $(a-b)^2=a^2-ab-ba+b^2=a^2-2ab+b^2$. Enfin
$(a-b)(a+b)=a^2+ab-ab-b^2=a^2-b^2$ : les deux termes opposés s’annulent.
Lues de droite à gauche, ces égalités permettent de factoriser.
\end{demonstration}
\begin{demonstration}[Le carré d’un entier impair est impair]
Un entier impair s’écrit $2n+1$ avec n entier. Son carré vaut
$(2n+1)^2=4n^2+4n+1=2(2n^2+2n)+1$. Le nombre entre parenthèses est entier :
on obtient bien un multiple de 2 augmenté de 1.
\end{demonstration}
\editionRetenir{Une écriture générale prouve ce que quelques essais ne font que suggérer.}

\end{blocEdition}
\editionPage{preuves-2}{Comprendre les preuves}{Pythagore par les aires}{Les arguments ci-dessous éclairent des propriétés utilisées dans le livre.}
\begin{blocEdition}{preuves-2}

Considérons quatre copies d’un triangle rectangle de côtés perpendiculaires
a et b et d’hypoténuse c, disposées dans un carré de côté $a+b$.
\begin{center}\begin{tikzpicture}[scale=.65,font=\small]
\draw[geoTrait](0,0)rectangle(5,5);
\fill[otGreenTint](2,0)--(5,2)--(3,5)--(0,3)--cycle;
\draw[geoTrait](2,0)--(5,2)--(3,5)--(0,3)--cycle;
\node[below]at(1,0){$a$};\node[below]at(3.5,0){$b$};
\node[right]at(5,1){$a$};\node[right]at(5,3.5){$b$};
\node at(2.5,2.5){aire $c^2$};
\end{tikzpicture}\end{center}
Chaque côté du quadrilatère central est une hypoténuse de longueur c.
Ses angles sont droits : les deux angles aigus voisins des triangles
ont une somme de 90 degrés, et complètent un angle plat avec l’angle central.
Le quadrilatère est donc un carré.\par\medskip
L’aire du grand carré vaut $(a+b)^2$. Elle est aussi la somme de l’aire
du carré central et de celles des quatre triangles :
\[
(a+b)^2=c^2+4\times\frac{ab}{2}=c^2+2ab.
\]
En développant puis en soustrayant $2ab$, on obtient
$\res{a^2+b^2=c^2}$.

\end{blocEdition}
\editionPage{preuves-3}{Comprendre les preuves}{Thalès par les aires}{Les arguments ci-dessous éclairent des propriétés utilisées dans le livre.}
\begin{blocEdition}{preuves-3}

Dans le triangle ABC, M appartient à [AB], N à [AC] et $(MN)\parallel(BC)$.
On note $\mathcal A$ l’aire d’un triangle.
\begin{center}\begin{tikzpicture}[scale=.5,font=\small]
\draw[geoTrait](0,0)node[below left]{$A$}--(10,0)node[below right]{$B$}--(4,6)node[above]{$C$}--cycle;
\draw[geoTrait](6,0)node[below]{$M$}--(2.4,3.6)node[left]{$N$};
\draw[otMuted,dashed](6,0)--(4,6) (2.4,3.6)--(10,0);
\end{tikzpicture}\end{center}
\begin{demonstration}[Deux triangles de même aire]
B et C sont sur une droite parallèle à (MN) : ils sont à la même distance
de (MN). Les triangles MNB et MNC ont donc la même base [MN] et la même
hauteur : $\mathcal A(MNB)=\mathcal A(MNC)$. En ajoutant l’aire de AMN,
on obtient $\mathcal A(ANB)=\mathcal A(AMC)$.
\end{demonstration}
\begin{demonstration}[L’égalité des rapports]
Les triangles AMN et ABN ont la même hauteur issue de N ; leurs bases
[AM] et [AB] sont sur la même droite. Donc
$\dfrac{\mathcal A(AMN)}{\mathcal A(ABN)}=\dfrac{AM}{AB}$. De même, avec la
hauteur issue de M, $\dfrac{\mathcal A(AMN)}{\mathcal A(AMC)}=\dfrac{AN}{AC}$.
Les dénominateurs sont égaux d’après la première étape :
$\res{\dfrac{AM}{AB}=\dfrac{AN}{AC}}$.
\end{demonstration}

\end{blocEdition}
\editionPage{preuves-4}{Comprendre les preuves}{Trigonométrie et vecteurs}{Les arguments ci-dessous éclairent des propriétés utilisées dans le livre.}
\begin{blocEdition}{preuves-4}

\begin{demonstration}[$(\cos\widehat B)^2+(\sin\widehat B)^2=1$]
Dans ABC rectangle en A, $\cos\widehat B=AB/BC$ et $\sin\widehat B=AC/BC$. Donc
$(\cos\widehat B)^2+(\sin\widehat B)^2=\dfrac{AB^2+AC^2}{BC^2}$. D’après le théorème
de Pythagore, $AB^2+AC^2=BC^2$ : la somme vaut 1. Connaissant le cosinus
d’un angle aigu, on en déduit son sinus, qui est positif.
\end{demonstration}
\begin{demonstration}[La règle du parallélogramme]
Si ABCD est un parallélogramme, ses côtés opposés [AD] et [BC] sont
parallèles, de même longueur, et parcourus dans le même sens :
$\vect{AD}=\vect{BC}$. Par la relation de Chasles,
$\vect{AB}+\vect{AD}=\vect{AB}+\vect{BC}=\vect{AC}$.
\end{demonstration}
\editionRetenir{Une égalité de vecteurs traduit à la fois un parallélisme, un sens et une égalité de longueurs.}

\end{blocEdition}

\partieCorriges
\clearpage\pagestyle{otMain}\fontsize{11}{15}\selectfont\corrigesChapitre{1}{Nombres rationnels et fractions}
\begin{blocLivre}{cor-1-1}\begin{corrige}{1}{Une température}\demarche{Traduire chaque changement avec son signe.}
On respecte l’ordre des changements : $-6+11=5$, puis $5-4=\res{1\,^\circ\mathrm C}$. Une augmentation se traduit par une addition, même si la température initiale est négative.
\end{corrige}\end{blocLivre}
\begin{blocLivre}{cor-1-2}\begin{corrige}{2}{Respecter les priorités}\demarche{Calculer les parenthèses avant les produits.}
Dans $A$, la parenthèse vaut 7, puis $3\times7=21$ : $\res{A=-3}$. Dans $B$, les deux parenthèses donnent 15 et 7 : $\res{B=105}$. Les parenthèses ont changé les facteurs du produit.
\end{corrige}\end{blocLivre}
\begin{blocLivre}{cor-1-3}\begin{corrige}{3}{Deux écritures, un nombre}\demarche{Appliquer le même facteur aux deux termes.}
On multiplie 8 par 5, donc aussi 5 par 5 : $\res{25/40}$. On divise 42 par 21, donc aussi 63 par 21 : $\res{2/3}$. La même opération sur les deux termes conserve le quotient.
\end{corrige}\end{blocLivre}
\begin{blocLivre}{cor-1-4}\begin{corrige}{4}{Une somme exacte}\demarche{Choisir un dénominateur commun adapté.}
Un dénominateur commun est 36 : $7/12=21/36$ et $5/18=10/36$. La somme est $\res{31/36}$. Le numérateur 31 est premier et ne divise pas 36 : la fraction est irréductible.
\end{corrige}\end{blocLivre}
\begin{blocLivre}{cor-1-5}\begin{corrige}{5}{Produit et quotient}\demarche{Simplifier les facteurs et inverser le diviseur.}
Dans le produit, on simplifie 7 avec 14 et 9 avec 15 : $\res{-3/10}$. Dans le quotient, on multiplie par l’inverse : $(3/5)\times(-10/9)=\res{-2/3}$. Les deux résultats sont négatifs.
\end{corrige}\end{blocLivre}
\begin{blocLivre}{cor-1-6}\begin{corrige}{6}{Simplifications successives}\demarche{Chercher un diviseur commun visible, puis recommencer.}
$150/225=30/45=6/9=\res{2/3}$, en simplifiant par 5, par 5, puis par 3. $126/198=63/99=7/11$, en simplifiant par 2 puis par 9 : $\res{7/11}$. Comme $2\times11=22$ et $7\times3=21$, les fractions $\res{\text{ne sont pas égales}}$.
\end{corrige}\end{blocLivre}
\begin{blocLivre}{cor-1-7}\begin{corrige}{7}{Comparer sans arrondir}\demarche{Ramener les fractions à un dénominateur positif commun.}
Avec le dénominateur 12, les numérateurs sont $-9,-10,-7$. Ainsi $\res{-5/6<-3/4<-7/12}$. Parmi des négatifs, le nombre le plus éloigné de zéro est le plus petit.
\end{corrige}\end{blocLivre}
\begin{blocLivre}{cor-1-8}\begin{corrige}{8}{Une expression à étages}\demarche{Séparer le numérateur et le dénominateur avant de diviser.}
Le numérateur vaut $4/6+1/6=5/6$ ; le dénominateur vaut $5/4-2/4=3/4$, non nul. Donc $C=(5/6)\times(4/3)=\res{10/9}$.
\end{corrige}\end{blocLivre}
\begin{blocLivre}{cor-1-9}\begin{corrige}{9}{Des ensembles de nombres}\demarche{Simplifier l’écriture avant de choisir l’ensemble.}
$-7\in\res{\mathbb Z}$ ; $0\in\res{\mathbb N}$ ; $2{,}5=25/10\in\res{\mathbb D}$ ; $12/4=3\in\res{\mathbb N}$ ; $-5/3\in\res{\mathbb Q}$. Ce dernier n’est pas décimal : si $5/3$ s’écrivait $a/10^k$, on aurait $3a=5\times10^k$, impossible car 3 ne divise pas $5\times10^k$.
\end{corrige}\end{blocLivre}
\begin{blocLivre}{cor-1-10}\begin{corrige}{10}{Une réduction trompeuse}\demarche{Revenir au sens de la barre de fraction.}
La division porte sur toute la somme : $(8+4)/4=12/4=\res3$. On peut aussi diviser chaque terme : $8/4+4/4=2+1$. Simplifier seulement le terme 4 a laissé le 8 sans division.
\end{corrige}\end{blocLivre}
\begin{blocLivre}{cor-1-11}\begin{corrige}{11}{Deux étapes de lecture}\demarche{Identifier le tout auquel se rapporte chaque fraction.}
Après lundi, il reste $3/5$. Mardi, il lit $(1/3)\times(3/5)=1/5$ du livre entier. Il reste donc $3/5-1/5=\res{2/5}$. Cela représente $270\times2/5=\res{108\text{ pages}}$. Additionner $2/5$ et $1/3$ aurait confondu les références.
\end{corrige}\end{blocLivre}
\begin{blocLivre}{cor-1-12}\begin{corrige}{12}{Des bouquets vendus}\demarche{Exprimer ce qui reste comme une fraction du stock initial.}
Après le matin, il reste $1/3$ du stock. L’après-midi, elle vend les $3/4$ de ce tiers ; il en reste $1/4$, soit $(1/4)\times(1/3)=1/12$ du stock initial. Ces $1/12$ valent 9 bouquets : elle en avait $\res{108}$. Contrôle : 72 le matin, puis 27 l’après-midi, et $108-72-27=9$.
\end{corrige}\end{blocLivre}
\begin{blocLivre}{cor-1-13}\begin{corrige}{13}{Une réserve d’eau}\demarche{Exprimer les volumes comme fractions d’une même capacité.}
Après consommation, la part remplie est $(7/8)(1-2/7)=5/8$. Les 30 L ajoutés représentent $3/4-5/8=1/8$ de la capacité. La capacité vaut donc $30\times8=\res{240\text{ L}}$. Contrôle : 210 L au départ, puis 150 L, puis 180 L, soit $3/4$ de 240 L.
\end{corrige}\end{blocLivre}
\begin{blocLivre}{cor-1-14}\begin{corrige}{14}{Des votes et des absents}\demarche{Multiplier les proportions lorsque les ensembles sont emboîtés.}
La part de toute la classe qui vote A est $(5/6)(3/5)=1/2$. Les 15 voix représentent donc la moitié de la classe : $\res{30\text{ élèves}}$. Il y a $30/6=\res{5\text{ absents}}$ et 25 présents ; $3/5$ de 25 vaut bien 15.
\end{corrige}\end{blocLivre}
\begin{blocLivre}{cor-1-15}\begin{corrige}{15}{Trouver toutes les possibilités}\demarche{Transformer les comparaisons en un encadrement entier.}
Les trois fractions se comparent avec le dénominateur 12 : $2/3=8/12$ et $5/4=15/12$. Il faut donc $8<n<15$. Les solutions sont $\res{9,10,11,12,13,14}$. Les deux bornes sont exclues à cause du mot « strictement ».
\end{corrige}\end{blocLivre}

\clearpage\pagestyle{otMain}\fontsize{11}{15}\selectfont\corrigesChapitre{2}{Puissances, racines et notation scientifique}
\begin{blocLivre}{cor-2-1}\begin{corrige}{1}{Comprendre un exposant}\demarche{Relier l’écriture aux produits et aux inverses.}
Le premier produit vaut $\res{7^4}$. L’exposant négatif donne un inverse : $2^{-3}=1/2^3=\res{1/8}$. La base 9 est non nulle, donc $\res{9^0=1}$.
\end{corrige}\end{blocLivre}
\begin{blocLivre}{cor-2-2}\begin{corrige}{2}{Parenthèses et signe}\demarche{Déterminer la base exacte avant de calculer.}
Le carré de $-5$ vaut $\res{25}$, tandis que $-5^2=-(5^2)=\res{-25}$. Une puissance impaire de $-2$ est négative : $\res{-8}$ ; une puissance paire est positive : $\res{16}$.
\end{corrige}\end{blocLivre}
\begin{blocLivre}{cor-2-3}\begin{corrige}{3}{Règles de puissances}\demarche{Identifier la règle correspondant à chaque opération.}
Pour un produit de même base, on additionne les exposants : $\res{3^6}$. Pour le quotient, on les soustrait : $\res{5^4}$. Deux puissances de même exposant donnent $(2\times5)^3=\res{10^3}$. Enfin $10^{-2-4}=\res{10^{-6}}$.
\end{corrige}\end{blocLivre}
\begin{blocLivre}{cor-2-4}\begin{corrige}{4}{Écriture scientifique}\demarche{Contrôler le coefficient et le nombre de déplacements.}
On déplace la virgule jusqu’à obtenir un coefficient entre 1 et 10 : $\res{7{,}4\times10^7}$ et $\res{9{,}2\times10^{-5}}$. Le signe négatif de l’exposant traduit ici un nombre inférieur à 1.
\end{corrige}\end{blocLivre}
\begin{blocLivre}{cor-2-5}\begin{corrige}{5}{Quelques racines}\demarche{Utiliser les carrés connus en conservant le signe positif.}
On cherche des nombres positifs dont les carrés sont les données : $\res{12}$, $\res{0{,}9}$ et $\res{7/8}$. Par exemple $(7/8)^2=49/64$.
\end{corrige}\end{blocLivre}
\begin{blocLivre}{cor-2-6}\begin{corrige}{6}{Produit scientifique}\demarche{Regrouper puis normaliser le coefficient.}
Les coefficients donnent 24 et les puissances $10^{8-5}=10^3$. On obtient $24\times10^3=\res{2{,}4\times10^4}$. Le coefficient 24 n’était pas encore compris entre 1 et 10.
\end{corrige}\end{blocLivre}
\begin{blocLivre}{cor-2-7}\begin{corrige}{7}{Quotient scientifique}\demarche{Soustraire l’exposant du dénominateur, même s’il est positif.}
Le quotient des coefficients vaut 4 et celui des puissances vaut $10^{-3-5}=10^{-8}$. Le résultat est $\res{4\times10^{-8}}$. On vérifie qu’il est petit, puisque l’on divise un petit nombre par un grand.
\end{corrige}\end{blocLivre}
\begin{blocLivre}{cor-2-8}\begin{corrige}{8}{Des préfixes}\demarche{Remplacer le préfixe par sa puissance de 10, puis normaliser.}
Giga vaut $10^9$ : $\res{4{,}7\times10^9}$ octets. Nano vaut $10^{-9}$ : $250\times10^{-9}=\res{2{,}5\times10^{-7}\text{ m}}$. Micro vaut $10^{-6}$ : $0{,}6\times10^{-6}=\res{6\times10^{-7}\text{ s}}$. Les coefficients sont ensuite ramenés entre 1 et 10.
\end{corrige}\end{blocLivre}
\begin{blocLivre}{cor-2-9}\begin{corrige}{9}{Une équation $x^2=a$}\demarche{Examiner d’abord le signe du second membre.}
Les ensembles de solutions sont $\res{\{-7\,;7\}}$, $\res{\{-\sqrt7\,;\sqrt7\}}$, $\res{\{0\}}$ et $\res{\varnothing}$. Un carré n’est jamais négatif, ce qui explique le dernier cas.
\end{corrige}\end{blocLivre}
\begin{blocLivre}{cor-2-10}\begin{corrige}{10}{Réfuter une règle}\demarche{Choisir un contre-exemple dont les deux calculs sont exacts.}
Prendre $a=9$ et $b=16$ : $\sqrt{9+16}=\sqrt{25}=5$, tandis que $\sqrt9+\sqrt{16}=3+4=7$. Comme $5\ne7$, cet exemple suffit à réfuter l’affirmation générale.
\end{corrige}\end{blocLivre}
\begin{blocLivre}{cor-2-11}\begin{corrige}{11}{Une population qui double}\demarche{Compter les périodes avant d’appliquer la puissance.}
Deux heures représentent 120 minutes, donc six doublements. Le nombre prévu est $800\times2^6=\res{51200}$, soit $\res{5{,}12\times10^4}$. On multiplie six fois par 2 ; on n’ajoute pas six fois 800.
\end{corrige}\end{blocLivre}
\begin{blocLivre}{cor-2-12}\begin{corrige}{12}{Une masse minuscule}\demarche{Diviser puis convertir sans perdre les puissances de 10.}
On divise la masse totale par le nombre de particules : $1{,}8/(4\times10^{12})=0{,}45\times10^{-12}=\res{4{,}5\times10^{-13}\text{ g}}$. Comme $1\text{ g}=10^{-3}\text{ kg}$, cela vaut $\res{4{,}5\times10^{-16}\text{ kg}}$.
\end{corrige}\end{blocLivre}
\begin{blocLivre}{cor-2-13}\begin{corrige}{13}{Une distance astronomique}\demarche{Utiliser des unités compatibles avant de multiplier.}
La durée vaut $8\times60+20=500$ s. La distance est $3\times10^5\times500=\res{1{,}5\times10^8\text{ km}}$. On multiplie par $10^3$ pour convertir en mètres : $\res{1{,}5\times10^{11}\text{ m}}$.
\end{corrige}\end{blocLivre}
\begin{blocLivre}{cor-2-14}\begin{corrige}{14}{Choisir le côté d’une dalle}\demarche{Distinguer la longueur exacte et la mesure arrondie.}
Le côté $c$ vérifie $c^2=2$ avec $c>0$, donc $\res{c=\sqrt2\text{ m}}$, soit environ 1,414213 m. Au millimètre, on obtient $\res{1{,}414\text{ m}}$. Remplacer par 1,4 m donnerait $1{,}4^2=1{,}96$ m$^2$, différent de 2 : l’approximation ne conserve pas l’aire exacte.
\end{corrige}\end{blocLivre}
\begin{blocLivre}{cor-2-15}\begin{corrige}{15}{Un cercle et trois droites}\demarche{Se ramener à une équation $x^2=a$ et discuter selon le signe de $a$.}
Pour $y=3$ : $x^2=25-9=16$, donc deux points d’abscisses $\res{-4\text{ et }4}$. Pour $y=5$ : $x^2=0$, un seul point, $\res{(0\,;5)}$ : la droite touche le cercle. Pour $y=6$ : $x^2=-11$, $\res{\text{aucun point}}$. Le nombre d’intersections dépend de la distance de la droite au centre, comparée au rayon.
\label{fig-cor-2-15}\figureGeo[corrige]{cercle-droites}
\end{corrige}\end{blocLivre}

\clearpage\pagestyle{otMain}\fontsize{11}{15}\selectfont\corrigesChapitre{3}{Multiples, diviseurs et nombres premiers}
\begin{blocLivre}{cor-3-1}\begin{corrige}{1}{Tester des diviseurs}\demarche{Utiliser les critères puis contrôler les quotients.}
2340 est pair et se termine par 0 : 2, 5 et 10 conviennent. La somme de ses chiffres vaut 9 : 3 et 9 conviennent aussi. Les cinq nombres proposés sont donc des diviseurs.
\end{corrige}\end{blocLivre}
\begin{blocLivre}{cor-3-2}\begin{corrige}{2}{Premier ou composé ?}\demarche{Appliquer la définition et fournir des facteurs quand c’est possible.}
$\res{2\text{ et }17}$ sont premiers. Le nombre 1 n’a qu’un diviseur positif ; $21=3\times7$ et $49=7^2$ sont composés. Pour 17, les seuls diviseurs premiers à tester sont 2 et 3, car $5^2=25>17$ ; ils ne divisent pas 17.
\end{corrige}\end{blocLivre}
\begin{blocLivre}{cor-3-3}\begin{corrige}{3}{Décomposer un entier}\demarche{Poursuivre la décomposition jusqu’aux facteurs premiers.}
$420=42\times10=(2\times3\times7)(2\times5)=\res{2^2\times3\times5\times7}$. Tous les facteurs sont premiers ; leur produit redonne 420.
\end{corrige}\end{blocLivre}
\begin{blocLivre}{cor-3-4}\begin{corrige}{4}{Une fraction irréductible}\demarche{Simplifier tous les facteurs premiers communs.}
Les décompositions sont $198=2\times3^2\times11$ et $252=2^2\times3^2\times7$. En simplifiant $2\times3^2$, on obtient $\res{11/14}$. Les diviseurs de 11 ne divisent pas 14.
\end{corrige}\end{blocLivre}
\begin{blocLivre}{cor-3-5}\begin{corrige}{5}{Écrire des multiples}\demarche{Encadrer les multiplicateurs entiers.}
On part de $14\times5=70$ puis on ajoute 14 : $\res{70,84,98,112,126}$. Le suivant vaut 140 et dépasse la borne supérieure.
\end{corrige}\end{blocLivre}
\begin{blocLivre}{cor-3-6}\begin{corrige}{6}{Tous les diviseurs}\demarche{Chercher des paires de facteurs sans doublons.}
On utilise les paires $1\times72$, $2\times36$, $3\times24$, $4\times18$, $6\times12$ et $8\times9$. On obtient $\res{1,2,3,4,6,8,9,12,18,24,36,72}$. Au-delà de 8, chaque diviseur a déjà son partenaire dans une paire, car $9\times9>72$.
\end{corrige}\end{blocLivre}
\begin{blocLivre}{cor-3-7}\begin{corrige}{7}{Un chiffre inconnu}\demarche{Combiner les critères, puis vérifier les candidats.}
La divisibilité par 2 impose $a\in\{0,2,4,6,8\}$. La somme des chiffres vaut $9+a$ ; elle est divisible par 3 pour $a=0,3,6,9$. Les chiffres communs aux deux listes donnent $\res{a=0\text{ ou }6}$.
\end{corrige}\end{blocLivre}
\begin{blocLivre}{cor-3-8}\begin{corrige}{8}{Une grande fraction}\demarche{Soustraire les exposants seulement pour les facteurs de même base.}
On simplifie $2^2$ et $3^2$. Il reste $(2^2\times7)/(3\times5)=\res{28/15}$. Aucun facteur 2 ou 7 ne peut encore se simplifier avec 3 ou 5.
\end{corrige}\end{blocLivre}
\begin{blocLivre}{cor-3-9}\begin{corrige}{9}{Un carré parfait ?}\demarche{Observer la parité de chaque exposant premier.}
Un carré a des exposants pairs dans sa décomposition. L’exposant 3 du facteur 3 est impair : le nombre n’est pas un carré. Multiplier par $\res3$ donne $2^4\times3^4=(2^2\times3^2)^2=\res{36^2}$. Un facteur 3 est nécessaire.
\end{corrige}\end{blocLivre}
\begin{blocLivre}{cor-3-10}\begin{corrige}{10}{Un calcul par tableur}\demarche{Relier le reste nul à la définition d’un diviseur.}
La fonction MOD calcule le reste de la division par la valeur de A2. Un reste nul signifie que cet entier divise 84. La liste cherchée est $\res{1,2,3,4,6,7,12,14}$. On ne trouve pas les diviseurs supérieurs à 20 avec cette seule colonne.
\end{corrige}\end{blocLivre}
\begin{blocLivre}{cor-3-11}\begin{corrige}{11}{Préparer des lots}\demarche{Traduire le partage en diviseurs communs.}
Le nombre de lots doit diviser 156 et 204. Or $156=2^2\times3\times13$ et $204=2^2\times3\times17$ : le plus grand diviseur commun vaut 12. On peut former $\res{12\text{ lots}}$, contenant chacun $\res{13\text{ rouges et }17\text{ bleues}}$. Aucun diviseur commun n’est plus grand.
\end{corrige}\end{blocLivre}
\begin{blocLivre}{cor-3-12}\begin{corrige}{12}{Des cigales périodiques}\demarche{Chercher le premier multiple commun aux deux périodes.}
Le nombre d’années écoulées doit être un multiple commun de 13 et de 17. Ces deux nombres sont premiers et distincts : le plus petit multiple commun positif est $13\times17=221$. La prochaine coïncidence aura lieu en $\res{2245}$. Des périodes premières entre elles rendent les rencontres rares.
\end{corrige}\end{blocLivre}
\begin{blocLivre}{cor-3-13}\begin{corrige}{13}{Daller un rectangle}\demarche{Utiliser un diviseur commun pour les deux dimensions.}
Le côté doit diviser 180 et 252. Leur plus grand diviseur commun vaut 36. On place $252/36=7$ carrés sur la longueur et $180/36=5$ sur la largeur : $\res{35\text{ carrés de côté }36\text{ cm}}$. Le contrôle par les aires donne $35\times36^2=180\times252$.
\label{fig-cor-3-13}\figureGeo[corrige]{dallage}
\end{corrige}\end{blocLivre}
\begin{blocLivre}{cor-3-14}\begin{corrige}{14}{Un stock à retrouver}\demarche{Tester la compatibilité des contraintes avant de chercher un nombre.}
Un tel nombre doit être multiple de $2^3$, de $3^2$ et de 5, donc de $8\times9\times5=360$. Les multiples voisins sont 360 et 720. Aucun n’est entre 500 et 600 : $\res{\text{aucune solution}}$. Un énoncé peut conduire légitimement à l’impossibilité.
\end{corrige}\end{blocLivre}
\begin{blocLivre}{cor-3-15}\begin{corrige}{15}{Des engrenages}\demarche{Compter les dents au contact plutôt que confondre les nombres de tours.}
Le nombre de dents passées au contact doit être multiple de 24 et de 40. Le premier multiple commun est 120. La petite roue fait $120/24=\res5$ tours et la grande $120/40=\res3$ tours, en sens opposés. Le contact initial revient après $\res{120\text{ passages}}$.
\label{fig-cor-3-15}\figureGeo[corrige]{engrenages}
\end{corrige}\end{blocLivre}

\clearpage\pagestyle{otMain}\fontsize{11}{15}\selectfont\corrigesChapitre{4}{Calcul littéral et identités remarquables}
\begin{blocLivre}{cor-4-1}\begin{corrige}{1}{Remplacer une lettre}\demarche{Écrire la substitution avant de simplifier.}
Pour $x=-2$, $A=2\times4+6+1=\res{15}$. Pour $x=3$, $A=18-9+1=\res{10}$. Les parenthèses autour de $-2$ évitent de calculer un carré avec le mauvais signe.
\end{corrige}\end{blocLivre}
\begin{blocLivre}{cor-4-2}\begin{corrige}{2}{Somme ou produit ?}\demarche{Repérer la dernière opération effectuée.}
$A$, $D$ et $E$ sont des $\res{\text{sommes}}$ (ou différences) : la dernière opération est une addition ou une soustraction. $B$ et $C$ sont des $\res{\text{produits}}$. $E$ contient un produit, mais on lui ajoute 3 en dernier.
\end{corrige}\end{blocLivre}
\begin{blocLivre}{cor-4-3}\begin{corrige}{3}{Prendre l’opposé}\demarche{Distribuer chaque facteur à toute sa parenthèse.}
Le signe moins change les deux signes : $-4x+7+6-2x=\res{-6x+13}$. Le produit par 2 porte sur les deux termes de la deuxième parenthèse.
\end{corrige}\end{blocLivre}
\begin{blocLivre}{cor-4-4}\begin{corrige}{4}{Double distributivité}\demarche{Écrire les quatre produits avant la réduction.}
On écrit les quatre produits : $3x^2-15x+2x-10$. Les deux termes en $x$ se regroupent, d’où $\res{3x^2-13x-10}$.
\end{corrige}\end{blocLivre}
\begin{blocLivre}{cor-4-5}\begin{corrige}{5}{Un facteur commun}\demarche{Rechercher un facteur numérique et éventuellement littéral.}
Le premier facteur commun est 7 : $\res{7(2x+3)}$. Pour le second, 4 et $x$ sont communs : $\res{4x(3x-2)}$. Redévelopper retrouve les expressions initiales.
\end{corrige}\end{blocLivre}
\begin{blocLivre}{cor-4-6}\begin{corrige}{6}{Simplifier des produits et des rapports}\demarche{Regrouper les facteurs de même nature, puis simplifier les facteurs communs.}
Le produit vaut $\res{6a^2b^2c}$. Dans le rapport, on simplifie $3ab$ : $\res{\dfrac{2b}{3a}}$. Dans le dernier, $\pi$ et $R^2$ sont communs : $\res{\dfrac43R}$. On simplifie des facteurs, jamais des termes d’une somme.
\end{corrige}\end{blocLivre}
\begin{blocLivre}{cor-4-7}\begin{corrige}{7}{Différences de carrés}\demarche{Identifier les deux carrés avant d’appliquer l’identité.}
On reconnaît $81=9^2$ et $25x^2=(5x)^2$. Les résultats sont $\res{(x-9)(x+9)}$ et $\res{(5x-2)(5x+2)}$.
\end{corrige}\end{blocLivre}
\begin{blocLivre}{cor-4-8}\begin{corrige}{8}{Deux carrés de sommes}\demarche{Conserver le double produit.}
Le premier développement est $\res{x^2+12x+36}$. Le second est $\res{4x^2-12x+9}$. Dans chaque cas, le terme du milieu est le double produit des deux termes de la parenthèse, avec son signe.
\end{corrige}\end{blocLivre}
\begin{blocLivre}{cor-4-9}\begin{corrige}{9}{Reconnaître un carré}\demarche{Identifier $a$ et $b$, puis contrôler le double produit.}
Dans la première, $x^2$ et $25=5^2$ sont des carrés et $10x=2\times x\times5$ : $\res{(x+5)^2}$. Dans la seconde, $9x^2=(3x)^2$, $4=2^2$ et $12x=2\times3x\times2$ : $\res{(3x-2)^2}$. Le signe du double produit choisit l’identité.
\end{corrige}\end{blocLivre}
\begin{blocLivre}{cor-4-10}\begin{corrige}{10}{Factoriser en deux temps}\demarche{Repérer un facteur commun avant une identité éventuelle.}
Dans la première expression, la parenthèse entière est commune : $\res{(x-2)(3x+5)}$. Dans la seconde, sortir 9 puis utiliser une différence de carrés : $\res{9(x-2)(x+2)}$.
\end{corrige}\end{blocLivre}
\begin{blocLivre}{cor-4-11}\begin{corrige}{11}{Deux programmes équivalents}\demarche{Traduire les étapes avec un nombre quelconque.}
Pour un nombre $x$, A donne $4(x+5)-8=4x+20-8=4x+12$. B donne directement $4x+12$. Les deux expressions sont identiques pour tout $x$. Tester seulement quelques nombres n’aurait pas suffi.
\end{corrige}\end{blocLivre}
\begin{blocLivre}{cor-4-12}\begin{corrige}{12}{Une bordure}\demarche{Soustraire les aires en calculant d’abord le bon côté extérieur.}
Le grand carré a pour côté $x+4$, puisque la bordure est des deux côtés. L’aire cherchée vaut $(x+4)^2-x^2=\res{8x+16}$ cm$^2$. Pour $x=7$, elle vaut $\res{72\text{ cm}^2}$. On contrôle $11^2-7^2=72$.
\label{fig-cor-4-12}\figureGeo[corrige]{bordure}
\end{corrige}\end{blocLivre}
\begin{blocLivre}{cor-4-13}\begin{corrige}{13}{Des entiers et leur parité}\demarche{Écrire la forme générale de chaque parité.}
Un impair s’écrit $2n+1$ avec $n$ entier. Son carré est $4n^2+4n+1=2(2n^2+2n)+1$, donc impair. Un pair s’écrit $2n$ ; son carré vaut $4n^2$, donc est bien multiple de 4. Les preuves portent sur un entier quelconque.
\end{corrige}\end{blocLivre}
\begin{blocLivre}{cor-4-14}\begin{corrige}{14}{Une égalité impossible ?}\demarche{Confronter la solution algébrique au domaine géométrique.}
La différence vaut $(x+3)^2-x^2=6x+9$. Pour obtenir 9, il faudrait $x=0$, exclu par la positivité du côté. Pour 39, $6x+9=39$ donne $x=5$, qui convient. Les côtés sont alors 8 et 5 ; $64-25=39$.
\end{corrige}\end{blocLivre}
\begin{blocLivre}{cor-4-15}\begin{corrige}{15}{Une différence toujours positive}\demarche{Comparer les deux quantités par leur différence.}
On calcule la différence : $A-B=x^2+2x+1-(x^2+2x)=1$. Donc $\res{A=B+1}$ pour tout $x$ ; A est toujours strictement supérieur à B. Cette méthode évite de distinguer les signes des facteurs de B.
\end{corrige}\end{blocLivre}

\clearpage\pagestyle{otMain}\fontsize{11}{15}\selectfont\corrigesChapitre{5}{Équations, inéquations et produits nuls}
\begin{blocLivre}{cor-5-1}\begin{corrige}{1}{Une opération inverse}\demarche{Appliquer l’opération inverse aux deux membres.}
On soustrait 13 à la première égalité : $\res{x=-9}$. On divise la seconde par $-4$ : $\res{x=-9/2}$. Les deux équations sont indépendantes.
\end{corrige}\end{blocLivre}
\begin{blocLivre}{cor-5-2}\begin{corrige}{2}{Deux étapes}\demarche{Isoler le produit avant l’inconnue.}
On ajoute 5 : $7x=35$. Puis on divise par 7 : $\res{x=5}$. Contrôle : $7\times5-5=30$.
\end{corrige}\end{blocLivre}
\begin{blocLivre}{cor-5-3}\begin{corrige}{3}{L’inconnue des deux côtés}\demarche{Rassembler les termes inconnus avant les constantes.}
On soustrait $5x$ : $3x+3=-12$. On retire 3 : $3x=-15$. Ainsi $\res{x=-5}$. Les deux membres initiaux valent $-37$.
\end{corrige}\end{blocLivre}
\begin{blocLivre}{cor-5-4}\begin{corrige}{4}{Avec une parenthèse}\demarche{Développer avant de regrouper.}
On développe : $6x-15=4x+9$, donc $2x=24$ et $\res{x=12}$. Vérification : $3(24-5)=57=48+9$.
\end{corrige}\end{blocLivre}
\begin{blocLivre}{cor-5-5}\begin{corrige}{5}{Une première inéquation}\demarche{Regarder le signe du nombre par lequel on divise.}
En divisant par 3, positif, on obtient $\res{x\ge4}$. En divisant par $-2$, négatif, le sens change : $\res{x\le-4}$. Les crochets fermés indiquent que 4 et $-4$ sont solutions.
\label{fig-cor-5-5}\figureGeo[corrige]{inequation-droites}
\end{corrige}\end{blocLivre}
\begin{blocLivre}{cor-5-6}\begin{corrige}{6}{Une équation fractionnaire}\demarche{Multiplier chaque terme par le même dénominateur commun.}
On multiplie toute l’égalité par 12 : $4x+3x=168$. Donc $7x=168$ et $\res{x=24}$. Contrôle : $24/3+24/4=8+6=14$.
\end{corrige}\end{blocLivre}
\begin{blocLivre}{cor-5-7}\begin{corrige}{7}{Tout ou rien}\demarche{Interpréter l’égalité restante lorsque l’inconnue disparaît.}
La première égalité se réduit à $4x-8=4x-8$ : $\res{\text{tous les nombres sont solutions}}$. La seconde imposerait $-8=1$ après soustraction de $4x$ : elle n’a $\res{\text{aucune solution}}$.
\end{corrige}\end{blocLivre}
\begin{blocLivre}{cor-5-8}\begin{corrige}{8}{Un produit nul}\demarche{Séparer les deux facteurs, sans développer.}
Au moins un facteur est nul : $2x-7=0$ ou $x+4=0$. On obtient $\res{x=7/2\text{ ou }x=-4}$. Les deux valeurs annulent bien le produit.
\end{corrige}\end{blocLivre}
\begin{blocLivre}{cor-5-9}\begin{corrige}{9}{Factoriser pour résoudre}\demarche{Factoriser, par un facteur commun ou une identité, avant d’appliquer le produit nul.}
$x^2-6x=x(x-6)$, donc $\res{x=0\text{ ou }x=6}$ ; diviser par $x$ aurait perdu 0. Ensuite $(x+3)^2-4^2=(x+3-4)(x+3+4)=(x-1)(x+7)$ : $\res{x=1\text{ ou }x=-7}$.
\end{corrige}\end{blocLivre}
\begin{blocLivre}{cor-5-10}\begin{corrige}{10}{Lire une inéquation sur un graphique}\demarche{Lire les abscisses des points situés sur ou au-dessus de la droite horizontale.}
La droite horizontale $y=6$ coupe la droite au point d’abscisse 4 ; la droite est au-dessus pour les abscisses plus grandes. Les solutions sont $\res{x\ge4}$. Par le calcul, on divise par 1,5, positif : $x\ge6/1{,}5=4$.
\label{fig-cor-5-10}\figureGeo[corrige]{inequation-graphique}
\end{corrige}\end{blocLivre}
\begin{blocLivre}{cor-5-11}\begin{corrige}{11}{Des billets de spectacle}\demarche{Exprimer les deux effectifs avec une seule inconnue.}
Notons $x$ le nombre de jeunes ; il y a $28-x$ autres billets. L’équation est $6x+10(28-x)=232$, soit $280-4x=232$. Ainsi $\res{x=12}$. Il reste 16 autres billets ; $72+160=232$. La solution est entière et comprise entre 0 et 28.
\end{corrige}\end{blocLivre}
\begin{blocLivre}{cor-5-12}\begin{corrige}{12}{Un rectangle à retrouver}\demarche{Traduire le périmètre avec les deux côtés.}
Notons $x$ sa largeur, sa longueur vaut $x+7$. Alors $2x+2(x+7)=58$, donc $4x=44$ et $x=11$. Les dimensions sont $\res{11\text{ cm et }18\text{ cm}}$, l’aire $\res{198\text{ cm}^2}$.
\label{fig-cor-5-12}\figureGeo[corrige]{rectangle-equation}
\end{corrige}\end{blocLivre}
\begin{blocLivre}{cor-5-13}\begin{corrige}{13}{Un achat rentable}\demarche{Résoudre l’inéquation, puis revenir aux valeurs entières du contexte.}
L’inéquation est $38+12n\le30n$, soit $38\le18n$, donc $n\ge38/18\approx2{,}1$. Le nombre d’années étant entier, la gourde devient rentable à partir de la $\res{3^{\text{e}}\text{ année}}$. Contrôle : 74 € contre 90 € à trois ans, 62 € contre 60 € à deux ans.
\end{corrige}\end{blocLivre}
\begin{blocLivre}{cor-5-14}\begin{corrige}{14}{Deux aires égales}\demarche{Développer pour faire disparaître les termes carrés communs.}
On pose $(x+2)^2=x(x+5)$. Après développement : $x^2+4x+4=x^2+5x$, donc $\res{x=4}$. Le carré a pour côté 6 et le rectangle pour côtés 4 et 9. L’aire commune est $\res{36\text{ cm}^2}$. Toutes les dimensions sont positives.
\label{fig-cor-5-14}\figureGeo[corrige]{aires-egales}
\end{corrige}\end{blocLivre}
\begin{blocLivre}{cor-5-15}\begin{corrige}{15}{Analyse et synthèse}\demarche{Déduire les valeurs possibles, puis vérifier chacune dans le problème.}
On développe : $(x+1)^2-25=x^2+2x-24=x(x+2)-24$. Analyse : si $x$ convient, $(x+1)^2-5^2=0$, donc $(x-4)(x+6)=0$, soit $x=4$ ou $x=-6$. Synthèse : une largeur est positive, donc seul 4 peut convenir ; un rectangle de $\res{4\text{ cm sur }6\text{ cm}}$ a bien une aire de 24 cm$^2$.
\end{corrige}\end{blocLivre}

\clearpage\pagestyle{otMain}\fontsize{11}{15}\selectfont\corrigesChapitre{6}{Proportionnalité, ratios et pourcentages}
\begin{blocLivre}{cor-6-1}\begin{corrige}{1}{Passer par l’unité}\demarche{Calculer le prix d’une unité avant celui du lot.}
Un cahier coûte $10{,}80/4=2{,}70$ €. Sept coûtent $7\times2{,}70=\res{18{,}90\text{ €}}$.
\end{corrige}\end{blocLivre}
\begin{blocLivre}{cor-6-2}\begin{corrige}{2}{Un tableau proportionnel ?}\demarche{Comparer tous les rapports correspondant à la même direction.}
Les trois rapports prix sur quantité valent $7/2=17{,}5/5=28/8=3{,}5$. Le tableau est proportionnel, de coefficient $\res{3{,}5}$, et modélisé par $\res{x\mapsto3{,}5x}$.
\end{corrige}\end{blocLivre}
\begin{blocLivre}{cor-6-3}\begin{corrige}{3}{Appliquer une réduction}\demarche{Multiplier par le coefficient de baisse.}
Le coefficient est $1-25/100=0{,}75$. Le nouveau prix vaut $64\times0{,}75=\res{48\text{ €}}$. La réduction est de 16 €.
\end{corrige}\end{blocLivre}
\begin{blocLivre}{cor-6-4}\begin{corrige}{4}{Calculer un taux}\demarche{Diviser la variation par la valeur de départ.}
La hausse vaut $184-160=24$. Rapportée à l’effectif initial, elle représente $24/160=0{,}15$, soit $\res{15\,\%}$.
\end{corrige}\end{blocLivre}
\begin{blocLivre}{cor-6-5}\begin{corrige}{5}{Partager selon un ratio}\demarche{Compter toutes les parts du ratio.}
Il y a $3+5=8$ parts, chacune de $56/8=7$ billes. Les deux groupes reçoivent $\res{21\text{ et }35\text{ billes}}$. Le total et le ratio sont vérifiés.
\end{corrige}\end{blocLivre}
\begin{blocLivre}{cor-6-6}\begin{corrige}{6}{Retrouver un prix initial}\demarche{Diviser par le coefficient d’évolution.}
Si le prix initial est $P$, $1{,}08P=81$. Ainsi $P=81/1{,}08=\res{75\text{ €}}$. Soustraire 8 \% de 81 utiliserait une autre référence.
\end{corrige}\end{blocLivre}
\begin{blocLivre}{cor-6-7}\begin{corrige}{7}{Deux baisses}\demarche{Appliquer chaque taux à la valeur actualisée.}
Les coefficients se multiplient : $250\times0{,}9\times0{,}8=\res{180\text{ €}}$. Le coefficient global est 0,72, donc la baisse totale est $\res{28\,\%}$, et non 30 \%.
\end{corrige}\end{blocLivre}
\begin{blocLivre}{cor-6-8}\begin{corrige}{8}{Une échelle}\demarche{Mettre les longueurs dans la même unité.}
On convertit 7,2 km en 720000 cm. Le rapport plan sur réel vaut $9/720000=1/80000$ : l’échelle est $\res{1/80000}$. Ainsi 4 cm représentent $4\times80000=320000$ cm, soit $\res{3{,}2\text{ km}}$.
\end{corrige}\end{blocLivre}
\begin{blocLivre}{cor-6-9}\begin{corrige}{9}{Une masse en trois parts}\demarche{Additionner les termes du ratio avant de partager.}
Les neuf parts totalisent 4,5 kg ; une part vaut 0,5 kg. Les masses sont donc $\res{2{,}5\text{ kg},\ 1{,}5\text{ kg},\ 0{,}5\text{ kg}}$. Leur somme vaut bien 4,5 kg.
\end{corrige}\end{blocLivre}
\begin{blocLivre}{cor-6-10}\begin{corrige}{10}{Proportionnellement à l’âge}\demarche{Simplifier le ratio, puis partager le total.}
Le partage suit le ratio $12:18$, soit $2:3$. Il y a 5 parts de $450/5=90$ €. Les sœurs reçoivent $\res{180\text{ €}}$ et $\res{270\text{ €}}$. Contrôle : $180/12=270/18=15$ € par année d’âge.
\end{corrige}\end{blocLivre}
\begin{blocLivre}{cor-6-11}\begin{corrige}{11}{Comparer deux promotions}\demarche{Traduire exactement l’ordre des avantages.}
A conduit à $120-30=\res{90\text{ €}}$. B conduit à $120\times0{,}8-10=\res{86\text{ €}}$. Avec les conditions indiquées, B fait payer 4 € de moins. Le calcul dépend de l’utilisation immédiate du bon.
\end{corrige}\end{blocLivre}
\begin{blocLivre}{cor-6-12}\begin{corrige}{12}{Un mélange à corriger}\demarche{Fixer la quantité conservée et chercher le nouveau total.}
Les 2 L de sirop doivent représenter 0,16 du volume final $V$ : $V=2/0{,}16=12{,}5$ L. Le mélange initial fait 10 L. Il faut ajouter $\res{2{,}5\text{ L d’eau}}$. Le volume de sirop reste 2 L, mais le total augmente.
\end{corrige}\end{blocLivre}
\begin{blocLivre}{cor-6-13}\begin{corrige}{13}{Deux pourcentages emboîtés}\demarche{Nommer le groupe de référence de chaque proportion.}
L’association compte $720\times0{,}45=324$ élèves. Le badminton en compte $324\times2/9=\res{72}$. Leur part dans le collège est $72/720=\res{10\,\%}$. On multiplie 45 \% par $2/9$, pas par 2.
\end{corrige}\end{blocLivre}
\begin{blocLivre}{cor-6-14}\begin{corrige}{14}{Lire une annonce}\demarche{Distinguer points de pourcentage, taux relatif et effectifs.}
La part augmente de $\res{3\text{ points}}$. Relativement à 12 \%, la hausse vaut $(15-12)/12=\res{25\,\%}$. Les effectifs passent de 96 à 120, soit $\res{24}$ cyclistes supplémentaires. Le maintien du total de 800 est indispensable à cette dernière comparaison.
\end{corrige}\end{blocLivre}
\begin{blocLivre}{cor-6-15}\begin{corrige}{15}{Proportionnalité et alignement}\demarche{Repérer deux triangles emboîtés de sommet $O$.}
Les droites $(AH)$ et $(BK)$ sont toutes deux perpendiculaires à l’axe des abscisses : elles sont parallèles. Les points $O,H,K$ et $O,A,B$ sont alignés. Thalès donne $\dfrac{AH}{BK}=\dfrac{OH}{OK}$, soit $\dfrac{3}{7{,}5}=\dfrac{4}{10}$. Donc $\dfrac34=\dfrac{7{,}5}{10}=\res{0{,}75}$ : les ordonnées s’obtiennent en multipliant les abscisses par 0,75.
\label{fig-cor-6-15}\figureGeo[corrige]{thales-repere}
\end{corrige}\end{blocLivre}

\clearpage\pagestyle{otMain}\fontsize{11}{15}\selectfont\corrigesChapitre{7}{Grandeurs quotients et conversions}
\begin{blocLivre}{cor-7-1}\begin{corrige}{1}{Des longueurs et des aires}\demarche{Distinguer le facteur des longueurs de celui des aires.}
On obtient $\res{3800\text{ m}}$, $\res{4200\text{ cm}^2}$ et $\res{0{,}65\text{ m}^2}$. Les deux dernières conversions utilisent le facteur 10000.
\end{corrige}\end{blocLivre}
\begin{blocLivre}{cor-7-2}\begin{corrige}{2}{Volumes et contenances}\demarche{Utiliser les équivalences entre volumes et contenances.}
$2{,}7\text{ m}^3=\res{2700\text{ L}}$. Comme un mL est un cm$^3$, 450 mL valent $\res{450\text{ cm}^3}$, soit $\res{0{,}45\text{ L}}$.
\end{corrige}\end{blocLivre}
\begin{blocLivre}{cor-7-3}\begin{corrige}{3}{Une durée décimale}\demarche{Convertir la partie d’heure avec 60.}
$0{,}75\times60=45$, donc $\res{1\text{ h }45\text{ min}}$. Inversement, $42/60=0{,}7$, donc $\res{2{,}7\text{ h}}$.
\end{corrige}\end{blocLivre}
\begin{blocLivre}{cor-7-4}\begin{corrige}{4}{Une vitesse moyenne}\demarche{Mettre toute la durée dans l’unité choisie.}
La durée vaut 2,5 h. La vitesse moyenne est $150/2{,}5=\res{60\text{ km/h}}$. Diviser par 2,30 aurait mal traduit les minutes.
\end{corrige}\end{blocLivre}
\begin{blocLivre}{cor-7-5}\begin{corrige}{5}{Un débit}\demarche{Calculer le volume d’une minute.}
Le débit est $84/7=\res{12\text{ L/min}}$. Le volume en 25 minutes vaut $12\times25=\res{300\text{ L}}$.
\end{corrige}\end{blocLivre}
\begin{blocLivre}{cor-7-6}\begin{corrige}{6}{Convertir une vitesse}\demarche{Revenir aux unités de base pour justifier le facteur 3,6.}
$90/3{,}6=\res{25\text{ m/s}}$. Dans l’autre sens, $12\times3{,}6=\res{43{,}2\text{ km/h}}$. Un contrôle direct emploie 1000 m et 3600 s.
\end{corrige}\end{blocLivre}
\begin{blocLivre}{cor-7-7}\begin{corrige}{7}{Le temps de réaction}\demarche{Séparer la réaction du freinage.}
La vitesse vaut $54/3{,}6=15$ m/s. La distance pendant 0,8 s est $15\times0{,}8=\res{12\text{ m}}$. Ce calcul ne comprend pas la distance de freinage.
\end{corrige}\end{blocLivre}
\begin{blocLivre}{cor-7-8}\begin{corrige}{8}{Une masse volumique}\demarche{Multiplier le volume par la masse d’une unité de volume.}
La masse vaut $\rho V=2{,}5\times360=\res{900\text{ g}}$, soit $\res{0{,}9\text{ kg}}$. Les cm$^3$ du volume sont cohérents avec ceux du dénominateur de la masse volumique.
\end{corrige}\end{blocLivre}
\begin{blocLivre}{cor-7-9}\begin{corrige}{9}{Une consommation électrique}\demarche{Convertir la durée avant de calculer l’énergie.}
45 minutes valent 0,75 h. L’énergie vaut $1{,}2\times0{,}75=\res{0{,}9\text{ kWh}}$. Le coût vaut $0{,}9\times0{,}24=\res{0{,}216\text{ €}}$, soit 0,22 € au centime. Le tarif est une donnée de l’exercice.
\end{corrige}\end{blocLivre}
\begin{blocLivre}{cor-7-10}\begin{corrige}{10}{Un téléchargement}\demarche{Remplacer les préfixes par des puissances de 10 avant de diviser.}
$2{,}4\text{ Go}=2{,}4\times10^9$ octets et $40\text{ Mo/s}=4\times10^7$ octets par seconde. La durée vaut $2{,}4\times10^9/(4\times10^7)=0{,}6\times10^2=\res{60\text{ s}}$, soit une minute.
\end{corrige}\end{blocLivre}
\begin{blocLivre}{cor-7-11}\begin{corrige}{11}{Aller et retour}\demarche{Additionner distances et durées avant de diviser.}
L’aller dure $12/18=2/3$ h ; le retour dure 1 h. Le total est 24 km en $5/3$ h. La vitesse moyenne vaut $24/(5/3)=\res{14{,}4\text{ km/h}}$. La moyenne 15 km/h des deux vitesses ne convient pas, car les durées diffèrent.
\end{corrige}\end{blocLivre}
\begin{blocLivre}{cor-7-12}\begin{corrige}{12}{Deux robinets}\demarche{Additionner les fractions de cuve remplies par minute.}
En une minute, ils remplissent respectivement $1/30$ et $1/45$ de la cuve. Ensemble : $1/30+1/45=1/18$. Le remplissage prend donc $\res{18\text{ min}}$. Additionner les deux durées aurait ignoré l’addition des débits.
\end{corrige}\end{blocLivre}
\begin{blocLivre}{cor-7-13}\begin{corrige}{13}{Choisir assez de peinture}\demarche{Calculer l’aire utile, les couches, puis l’achat entier.}
L’aire à peindre est $5\times2{,}6-2=11$ m$^2$. Les deux couches représentent 22 m$^2$ à couvrir, donc $22/8=2{,}75$ L. Un pot ne suffit pas ; il faut $\res{2\text{ pots}}$. On n’arrondit pas le besoin au-dessous.
\end{corrige}\end{blocLivre}
\begin{blocLivre}{cor-7-14}\begin{corrige}{14}{Une fuite discrète}\demarche{Construire le débit puis multiplier par la durée totale.}
Le débit est 0,1 mL/s. Une journée compte 86400 s, d’où $8640$ mL, soit $\res{8{,}64\text{ L}}$. En 30 jours : $\res{259{,}2\text{ L}}$. Le résultat dépend du volume de goutte choisi dans le modèle.
\end{corrige}\end{blocLivre}
\begin{blocLivre}{cor-7-15}\begin{corrige}{15}{Une livraison de gravier}\demarche{Employer les mètres pour les trois dimensions du volume.}
On convertit l’épaisseur en 0,06 m. Le volume vaut $12\times1{,}5\times0{,}06=1{,}08$ m$^3$. La masse est $1{,}08\times1{,}6=1{,}728$ tonne. Elle dépasse 1,5 tonne : $\res{2\text{ voyages}}$ suffisent et un seul est impossible.
\end{corrige}\end{blocLivre}

\clearpage\pagestyle{otMain}\fontsize{11}{15}\selectfont\corrigesChapitre{8}{Fonctions, images et fonction carré}
\begin{blocLivre}{cor-8-1}\begin{corrige}{1}{Calculer des images}\demarche{Substituer en gardant les parenthèses pour les négatifs.}
On remplace $x$ : $\res{f(0)=-4}$, $\res{f(2)=2}$ et $\res{f(-3)=-13}$. Chaque entrée donne une seule image.
\end{corrige}\end{blocLivre}
\begin{blocLivre}{cor-8-2}\begin{corrige}{2}{Lire un tableau}\demarche{Distinguer ce qui est connu dans le tableau et ce qui ne l’est pas.}
À la colonne de $-2$, l’image est $\res5$. La valeur 1 apparaît sous les entrées 0 et 3 : les antécédents indiqués sont $\res{0\text{ et }3}$. Le tableau ne dit pas s’il en existe ailleurs.
\end{corrige}\end{blocLivre}
\begin{blocLivre}{cor-8-3}\begin{corrige}{3}{Traduire une phrase}\demarche{Identifier le rôle de l’entrée et de la sortie.}
La première phrase s’écrit $\res{g(-1)=7}$. La deuxième signifie que 7 est l’image de 4, ou que 4 est un antécédent de 7. Elles sont compatibles : un nombre peut avoir plusieurs antécédents, mais une entrée n’a qu’une image.
\end{corrige}\end{blocLivre}
\begin{blocLivre}{cor-8-4}\begin{corrige}{4}{Un point sur une courbe}\demarche{Calculer l’image de chaque abscisse.}
$f(3)=9-2=7$ : $\res{A\text{ appartient à la courbe}}$. $f(-2)=4-2=2\ne1$ : $\res{B\text{ n’y appartient pas}}$. Une lecture approximative ne remplace pas ces calculs.
\end{corrige}\end{blocLivre}
\begin{blocLivre}{cor-8-5}\begin{corrige}{5}{Un antécédent algébrique}\demarche{Écrire l’égalité correspondant à l’image demandée.}
On résout $5x+2=17$, d’où $5x=15$ et $\res{x=3}$. Vérification : $h(3)=17$. Le coefficient 5 non nul assure une seule solution.
\end{corrige}\end{blocLivre}
\begin{blocLivre}{cor-8-6}\begin{corrige}{6}{Une formule et un programme}\demarche{Traduire chaque étape avant de transformer l’expression.}
La fonction est $\res{f:x\mapsto(x+3)^2}$. Ainsi $f(1)=16$ et $f(-3)=0$. Écrire $x^2+3$ n’aurait pas respecté l’ordre des instructions.
\end{corrige}\end{blocLivre}
\begin{blocLivre}{cor-8-7}\begin{corrige}{7}{Lire une ligne brisée}\demarche{Parcourir toute la courbe pour ne pas oublier une intersection.}
Le sommet d’abscisse $-1$ a pour ordonnée $\res3$. La droite horizontale d’ordonnée 1 rencontre la ligne aux abscisses $\res{-3,0,3}$. Sur le segment central, $f(x)=1-2x$, ce qui confirme $f(0)=1$.
\label{fig-cor-8-7}\figureGeo[corrige]{fonction-brisee}
\end{corrige}\end{blocLivre}
\begin{blocLivre}{cor-8-8}\begin{corrige}{8}{Plusieurs ou aucun antécédent}\demarche{Résoudre une équation de carré en examinant le signe.}
Les antécédents de 16 sont $\res{-4\text{ et }4}$. Celui de 0 est seulement $\res0$. Le nombre $-1$ n’a $\res{\text{aucun antécédent}}$, car un carré est positif ou nul.
\label{fig-cor-8-8}\figureGeo[corrige]{fonction-carre}
\end{corrige}\end{blocLivre}
\begin{blocLivre}{cor-8-9}\begin{corrige}{9}{Comparer deux représentations}\demarche{Relier une intersection à l’égalité des deux images.}
Le graphique suggère l’intersection $(3\,;5)$. L’équation $x+2=2x-1$ donne $\res{x=3}$ ; les deux images valent $\res5$. Le calcul valide la lecture.
\label{fig-cor-8-9}\figureGeo[corrige]{fonction-intersection}
\end{corrige}\end{blocLivre}
\begin{blocLivre}{cor-8-10}\begin{corrige}{10}{Représenter la fonction carré}\demarche{Calculer les images, placer les points, puis tracer une courbe lisse.}
Les images sont $\res{9\,;4\,;1\,;0{,}25\,;0\,;0{,}25\,;1\,;4\,;9}$. Deux nombres opposés ont la même image : la courbe est symétrique par rapport à l’axe des ordonnées. On relie les points par une courbe régulière, sans segments, le corrigé la montre en vert.
\label{fig-cor-8-10}\figureGeo[corrige]{carre-trace}
\end{corrige}\end{blocLivre}
\begin{blocLivre}{cor-8-11}\begin{corrige}{11}{Un réservoir}\demarche{Choisir les unités avant d’écrire la fonction.}
Le volume en m$^3$ vaut $\res{V(h)=0{,}8h}$, avec $0\le h\le1{,}5$. On convertit 600 L en 0,6 m$^3$. L’équation $0{,}8h=0{,}6$ donne $\res{h=0{,}75\text{ m}}$, qui respecte la hauteur maximale.
\label{fig-cor-8-11}\figureGeo[corrige]{cuve-hauteur}
\end{corrige}\end{blocLivre}
\begin{blocLivre}{cor-8-12}\begin{corrige}{12}{Résoudre graphiquement $x^2=a$}\demarche{Tracer la droite horizontale, puis confirmer par le calcul.}
La droite d’équation $y=5$ coupe la parabole en deux points d’abscisses comprises entre 2 et 2,5 en valeur absolue. Les valeurs exactes sont $\res{-\sqrt5\text{ et }\sqrt5}$, soit environ $\res{-2{,}2\text{ et }2{,}2}$ au dixième. L’équation $x^2=0$ a une seule solution, 0 ; $x^2=-2$ n’en a aucune, car la droite $y=-2$ ne rencontre pas la courbe.
\label{fig-cor-8-12}\figureGeo[corrige]{carre-equation}
\end{corrige}\end{blocLivre}
\begin{blocLivre}{cor-8-13}\begin{corrige}{13}{Des aires possibles}\demarche{Exprimer l’autre dimension avec la contrainte de périmètre.}
L’autre côté vaut $20-x$. L’aire est $\res{A(x)=x(20-x)}$ cm$^2$, avec $0<x<20$ pour un rectangle non aplati. Les images de 3 et de 17 valent toutes deux $\res{51\text{ cm}^2}$. Une même aire peut donc avoir plusieurs antécédents.
\label{fig-cor-8-13}\figureGeo[corrige]{rectangle-fonction}
\end{corrige}\end{blocLivre}
\begin{blocLivre}{cor-8-14}\begin{corrige}{14}{Une lecture de déplacement}\demarche{Distinguer distance au départ et longueur totale du parcours.}
On s’éloigne jusqu’à 600 m pendant les dix premières minutes, on reste à cette distance cinq minutes, puis on revient en dix minutes. Dans le modèle rectiligne annoncé, la distance totale est $\res{1200\text{ m}}$. La dernière ordonnée 0 indique le retour au départ, pas un trajet total nul.
\label{fig-cor-8-14}\figureGeo[corrige]{promenade}
\end{corrige}\end{blocLivre}
\begin{blocLivre}{cor-8-15}\begin{corrige}{15}{Résoudre par un produit nul}\demarche{Choisir une forme adaptée à l’image recherchée.}
Le produit est nul pour $x-2=0$ ou $3x+6=0$, d’où $\res{2\text{ et }-2}$. Par exemple $f(0)=(-2)\times6=\res{-12}$, strictement négatif. Un exemple suffit pour cette dernière demande d’existence.
\end{corrige}\end{blocLivre}

\clearpage\pagestyle{otMain}\fontsize{11}{15}\selectfont\corrigesChapitre{9}{Fonctions linéaires et affines}
\begin{blocLivre}{cor-9-1}\begin{corrige}{1}{Reconnaître les formes}\demarche{Comparer les expressions aux définitions.}
Les fonctions $\res{f,g,k}$ sont affines, car $k(x)=0x+7$. Seule $\res f$ est linéaire parmi elles. La fonction carré $h$ n’est pas affine.
\end{corrige}\end{blocLivre}
\begin{blocLivre}{cor-9-2}\begin{corrige}{2}{Calculer avec une linéaire}\demarche{Appliquer puis inverser la multiplication.}
$f(4)=-2{,}5\times4=\res{-10}$. L’égalité $-2{,}5x=15$ donne $\res{x=-6}$. On divise par un coefficient non nul.
\end{corrige}\end{blocLivre}
\begin{blocLivre}{cor-9-3}\begin{corrige}{3}{Lire les coefficients}\demarche{Relier les coefficients à leur rôle graphique.}
Le coefficient directeur est $\res{-3}$, l’ordonnée à l’origine $\res8$ et $g(2)=-6+8=\res2$. La droite descend de gauche à droite.
\end{corrige}\end{blocLivre}
\begin{blocLivre}{cor-9-4}\begin{corrige}{4}{Tracer une linéaire}\demarche{Choisir une entrée donnant une image simple.}
La droite passe par l’origine $O(0\,;0)$. Une autre image facile est $f(2)=3$, qui donne $A(2\,;3)$. On trace la droite passant par ces deux points ; le corrigé la montre en vert.
\label{fig-cor-9-4}\figureGeo[corrige]{affine-trace-lineaire}
\end{corrige}\end{blocLivre}
\begin{blocLivre}{cor-9-5}\begin{corrige}{5}{Tracer une affine}\demarche{Prendre les intersections avec les axes quand elles sont simples.}
On peut utiliser $(0\,;4)$ et $(4\,;0)$. L’antécédent de 0 est $\res4$, puisque $-4+4=0$. Le graphique et l’équation concordent.
\label{fig-cor-9-5}\figureGeo[corrige]{affine-trace}
\end{corrige}\end{blocLivre}
\begin{blocLivre}{cor-9-6}\begin{corrige}{6}{Une linéaire à retrouver}\demarche{Utiliser la forme annoncée avant les données numériques.}
On écrit $f(x)=ax$. Comme $6a=21$, $a=7/2$. Donc $\res{f(x)=3{,}5x}$ et $\res{f(-2)=-7}$. L’information « linéaire » est essentielle pour utiliser un seul point.
\end{corrige}\end{blocLivre}
\begin{blocLivre}{cor-9-7}\begin{corrige}{7}{Deux points, une formule}\demarche{Calculer le coefficient directeur puis utiliser un point pour la constante.}
Le coefficient directeur est $(10-1)/(5-2)=3$. On remplace les coordonnées de A dans $y=3x+b$ : $1=6+b$, donc $b=-5$. Ainsi $\res{f(x)=3x-5}$. On contrôle $f(5)=10$.
\label{fig-cor-9-7}\figureGeo[corrige]{affine-deux-points}
\end{corrige}\end{blocLivre}
\begin{blocLivre}{cor-9-8}\begin{corrige}{8}{Retrouver un tarif}\demarche{Interpréter la variation de prix avant de retrouver le forfait.}
L’augmentation de 12 € correspond à trois heures : le prix horaire est $\res{4\text{ €}}$. Pour deux heures, le forfait vaut $14-2\times4=\res{6\text{ €}}$. Le modèle est $P(h)=4h+6$.
\end{corrige}\end{blocLivre}
\begin{blocLivre}{cor-9-9}\begin{corrige}{9}{Coefficients lus sur un graphique}\demarche{Lire l’ordonnée à l’origine, puis la variation pour une unité d’abscisse.}
La droite coupe l’axe des ordonnées en 3 : $\res{b=3}$. Quand l’abscisse augmente de 1, par exemple de 0 à 1, l’ordonnée passe de 3 à 1 : elle diminue de 2, donc $\res{a=-2}$. La fonction est $\res{f(x)=-2x+3}$ ; la droite descend, ce qui confirme $a<0$.
\label{fig-cor-9-9}\figureGeo[corrige]{affine-lecture}
\end{corrige}\end{blocLivre}
\begin{blocLivre}{cor-9-10}\begin{corrige}{10}{Comparer par une différence}\demarche{Calculer une différence indépendante de l’entrée.}
La différence $g(x)-f(x)=7$ est toujours positive. Donc $\res{g(x)=f(x)+7}$ pour tout $x$. Les droites ont le même coefficient directeur et des ordonnées à l’origine différentes : elles sont parallèles.
\label{fig-cor-9-10}\figureGeo[corrige]{affine-paralleles}
\end{corrige}\end{blocLivre}
\begin{blocLivre}{cor-9-11}\begin{corrige}{11}{Deux abonnements}\demarche{Résoudre le seuil, puis l’inéquation, et revenir aux nombres entiers.}
L’égalité $20+3n=8+5n$ donne $\res{n=6}$ et un coût de 38 €. L’inéquation équivaut à $12<2n$, soit $n>6$ : A est moins chère $\res{\text{à partir de 7 séances}}$, B l’est pour moins de 6 séances. Les droites se croisent au point $(6\,;38)$.
\label{fig-cor-9-11}\figureGeo[corrige]{affine-tarifs}
\end{corrige}\end{blocLivre}
\begin{blocLivre}{cor-9-12}\begin{corrige}{12}{Une vidange à débit constant}\demarche{Associer la formule à son domaine physique.}
Le modèle est $\res{V(t)=900-35t}$. L’égalité $900-35t=200$ donne $\res{t=20\text{ min}}$. À 30 minutes, la formule donne $-150$ L, impossible : le modèle s’arrête quand la cuve est vide, vers 25,7 minutes.
\end{corrige}\end{blocLivre}
\begin{blocLivre}{cor-9-13}\begin{corrige}{13}{Une vitesse et une inéquation}\demarche{Modéliser par une fonction linéaire, puis résoudre l’inéquation.}
La distance est proportionnelle à la durée : $\res{d(t)=5t}$, fonction linéaire. L’inéquation $5t>120$ équivaut à $\res{t>24}$, car 5 est positif. Graphiquement, la droite $y=5t$ passe au-dessus de la droite horizontale $y=120$ après l’abscisse 24.
\end{corrige}\end{blocLivre}
\begin{blocLivre}{cor-9-14}\begin{corrige}{14}{Un point suffit-il ?}\demarche{Compter les informations et construire deux exemples.}
Non : $f(x)=x+5$ et $g(x)=3x+1$ donnent toutes deux 7 en 2, mais diffèrent pour d’autres entrées. Une condition unique ne détermine pas les deux coefficients. Pour une fonction linéaire, $b$ serait déjà imposé à zéro.
\end{corrige}\end{blocLivre}
\begin{blocLivre}{cor-9-15}\begin{corrige}{15}{Un seuil avec contrainte}\demarche{Résoudre chaque inéquation, puis prendre le plus grand entier solution.}
À 6 jours, A vaut 117 € et B 123 € : seule A respecte le budget. Pour A, $45+12d\le120$ donne $d\le6{,}25$ : le maximum est $\res6$ jours. Pour B, $15+18d\le120$ donne $d\le5{,}83\ldots$ : le maximum est $\res5$ jours. On ne garde que les durées entières.
\end{corrige}\end{blocLivre}

\clearpage\pagestyle{otMain}\fontsize{11}{15}\selectfont\corrigesChapitre{10}{Statistiques et boîtes à moustaches}
\begin{blocLivre}{cor-10-1}\begin{corrige}{1}{Trois indicateurs}\demarche{Additionner, repérer le rang central, puis soustraire les extrêmes.}
La somme vaut 55 : la moyenne vaut $55/5=\res{11}$. La valeur centrale, de rang 3, donne la médiane $\res{10}$. L’étendue vaut $15-7=\res8$. Moyenne et médiane répondent à des questions différentes.
\end{corrige}\end{blocLivre}
\begin{blocLivre}{cor-10-2}\begin{corrige}{2}{Une fréquence}\demarche{Identifier le total de référence.}
La fréquence est $12/30=\res{0{,}4}$, soit $\res{40\,\%}$. Le dénominateur est l’effectif de tout le groupe.
\end{corrige}\end{blocLivre}
\begin{blocLivre}{cor-10-3}\begin{corrige}{3}{Des effectifs cumulés}\demarche{Ajouter chaque effectif au cumul précédent.}
On additionne les effectifs de proche en proche : $\res{2\,;5\,;11\,;16\,;20}$. L’effectif total est le dernier effectif cumulé, $\res{20}$. Par exemple, 11 élèves ont une note inférieure ou égale à 10.
\end{corrige}\end{blocLivre}
\begin{blocLivre}{cor-10-4}\begin{corrige}{4}{Ordonner avant la médiane}\demarche{Ranger avant de choisir le rang central.}
La première liste ordonnée est $4\,;7\,;9\,;11\,;13$ : la médiane est la 3\textsuperscript{e} valeur, $\res9$. La seconde a six valeurs : on retient $(6+8)/2=\res7$. La valeur 20 n’influence pas cette médiane.
\end{corrige}\end{blocLivre}
\begin{blocLivre}{cor-10-5}\begin{corrige}{5}{Lire une boîte à moustaches}\demarche{Associer chaque trait de la boîte à un indicateur.}
On lit $\res{4\,;8\,;11\,;14\,;19}$ pour le minimum, $Q_1$, la médiane, $Q_3$ et le maximum. L’étendue vaut $19-4=\res{15\text{ min}}$. Comme $Q_3=14$, au moins $\res{75\,\%}$ des durées sont inférieures ou égales à 14 min.
\label{fig-cor-10-5}\figureGeo[corrige]{boite-lecture}
\end{corrige}\end{blocLivre}
\begin{blocLivre}{cor-10-6}\begin{corrige}{6}{Moyenne et quartiles d’un tableau}\demarche{Utiliser les effectifs cumulés pour trouver les rangs.}
La somme pondérée vaut $12+24+60+60+64=220$ : moyenne $\res{11}$. Avec $N=20$, les rangs 10 et 11 portent 10 : médiane $\res{10}$. $Q_1$ est la 5\textsuperscript{e} valeur, d’effectif cumulé atteint à 8 : $\res{Q_1=8}$. $Q_3$ est la 15\textsuperscript{e}, atteinte à 12 : $\res{Q_3=12}$.
\end{corrige}\end{blocLivre}
\begin{blocLivre}{cor-10-7}\begin{corrige}{7}{Quartiles d’une série brute}\demarche{Calculer les rangs, arrondis à l’entier supérieur.}
$N=11$. La médiane est la 6\textsuperscript{e} valeur : $\res{10}$. $11\times0{,}25=2{,}75$, donc $Q_1$ est la 3\textsuperscript{e} valeur : $\res6$. $11\times0{,}75=8{,}25$, donc $Q_3$ est la 9\textsuperscript{e} valeur : $\res{15}$.
\end{corrige}\end{blocLivre}
\begin{blocLivre}{cor-10-8}\begin{corrige}{8}{Construire une boîte}\demarche{Placer les cinq valeurs avant de dessiner la boîte.}
On place le minimum 3, $Q_1=6$, la médiane 10, $Q_3=15$ et le maximum 20. La boîte va de 6 à 15, avec un trait à 10 ; les moustaches vont de 3 à 6 et de 15 à 20. Le corrigé la montre en vert, avec l’échelle de l’axe.
\label{fig-cor-10-8}\figureGeo[corrige]{boite-construire}
\end{corrige}\end{blocLivre}
\begin{blocLivre}{cor-10-9}\begin{corrige}{9}{Quartiles et diagramme en barres}\demarche{Passer des hauteurs aux effectifs cumulés.}
Les effectifs 3, 7, 9, 6, 5 donnent les cumuls 3, 10, 19, 25, 30 : $\res{N=30}$. Les rangs 15 et 16 valent 3 livres : médiane $\res3$. $30\times0{,}25=7{,}5$ : rang 8, $\res{Q_1=2}$. $30\times0{,}75=22{,}5$ : rang 23, $\res{Q_3=4}$.
\label{fig-cor-10-9}\figureGeo[corrige]{stat-barres}
\end{corrige}\end{blocLivre}
\begin{blocLivre}{cor-10-10}\begin{corrige}{10}{Des formules de tableur}\demarche{Désigner une plage de cellules, puis contrôler à la main.}
On peut saisir \code{=MOYENNE(A2:A9)}, \code{=MEDIANE(A2:A9)} et \code{=MAX(A2:A9)-MIN(A2:A9)}. La moyenne vaut $114/8=\res{14{,}25}$. La série ordonnée est 9, 11, 12, 14, 15, 15, 18, 20 : médiane $(14+15)/2=\res{14{,}5}$. L’étendue vaut $\res{11}$.
\end{corrige}\end{blocLivre}
\begin{blocLivre}{cor-10-11}\begin{corrige}{11}{Réunir deux groupes}\demarche{Reconstituer les sommes de chaque groupe.}
Les sommes valent $12\times13=156$ et $18\times11=198$. La moyenne commune vaut $(156+198)/30=\res{11{,}8}$. Faire $(13+11)/2$ donnerait un poids égal à des groupes de tailles différentes.
\end{corrige}\end{blocLivre}
\begin{blocLivre}{cor-10-12}\begin{corrige}{12}{Comparer deux collèges}\demarche{Comparer d’abord les médianes, puis les longueurs des boîtes et les étendues.}
Les médianes sont 15 min pour A et 14 min pour B : positions proches. A est plus dispersé : étendue $40-5=35$ contre $25-8=17$, boîte de longueur $22-10=12$ contre $17-12=5$. Dans $\res{A}$, $Q_3=22$ : moins de 75 \% des durées sont inférieures à 22 min, donc au moins 25 \% des élèves mettent 22 min ou plus. Dans B, $Q_3=17$ : on sait seulement qu’au plus un quart des élèves mettent plus de 17 min, on ne peut pas l’affirmer.
\label{fig-cor-10-12}\figureGeo[corrige]{boites-comparer}
\end{corrige}\end{blocLivre}
\begin{blocLivre}{cor-10-13}\begin{corrige}{13}{Ajouter une observation}\demarche{Séparer ce que détermine la somme de ce qui demande l’ordre.}
L’ancienne somme est 126 ; la nouvelle vaut 150 pour 10 valeurs. La moyenne est $\res{15}$. La médiane reste inconnue : la somme ne renseigne pas sur les valeurs centrales. Par exemple neuf fois 14, ou huit fois 13 et une fois 22, ont la même moyenne initiale mais des médianes différentes après l’ajout.
\end{corrige}\end{blocLivre}
\begin{blocLivre}{cor-10-14}\begin{corrige}{14}{Même moyenne, mêmes résultats ?}\demarche{Comparer un indicateur de position et un de dispersion.}
Les deux sommes valent 40 : les moyennes sont 10. Les deux médianes sont 10. Les étendues sont $\res2$ et $\res{16}$. B est beaucoup plus dispersée ; les seuls indicateurs de position masquent cet écart.
\end{corrige}\end{blocLivre}
\begin{blocLivre}{cor-10-15}\begin{corrige}{15}{Un graphique trompeur}\demarche{Ramener le pourcentage à la valeur initiale.}
La hausse réelle est $(84-80)/80=\res{5\,\%}$. Les hauteurs visibles sont 2 et 6, dans un rapport 3. Ce rapport visuel ne traduit pas un triplement des inscriptions. Lire l’origine et les graduations avant de comparer.
\label{fig-cor-10-15}\figureGeo[corrige]{stat-axe-tronque}
\end{corrige}\end{blocLivre}

\clearpage\pagestyle{otMain}\fontsize{11}{15}\selectfont\corrigesChapitre{11}{Probabilités, évènements et simulations}
\begin{blocLivre}{cor-11-1}\begin{corrige}{1}{Un dé}\demarche{Énumérer les issues favorables.}
Les issues favorables sont 2, 4 et 6, parmi six issues équiprobables. La probabilité vaut $3/6=\res{1/2}$.
\end{corrige}\end{blocLivre}
\begin{blocLivre}{cor-11-2}\begin{corrige}{2}{Des boules colorées}\demarche{Compter les boules plutôt que les couleurs.}
Le total vaut 10. $P(R)=5/10=\res{1/2}$ et $P(\overline R)=1-1/2=\res{1/2}$. Trois couleurs ne signifient pas trois issues équiprobables.
\end{corrige}\end{blocLivre}
\begin{blocLivre}{cor-11-3}\begin{corrige}{3}{Écrire des évènements}\demarche{Lister les issues de chaque évènement avant de les combiner.}
$A=\{2;4;6\}$, $B=\{3;6\}$. $\res{A\cap B=\{6\}}$, $\res{A\cup B=\{2;3;4;6\}}$ et $\res{\overline A=\{1;3;5\}}$. L’issue 6 n’est comptée qu’une fois dans la réunion.
\end{corrige}\end{blocLivre}
\begin{blocLivre}{cor-11-4}\begin{corrige}{4}{Impossible ou certain}\demarche{Comparer l’évènement à toutes les issues.}
Obtenir 9 est impossible : probabilité $\res0$ ; cet évènement est l’ensemble vide $\varnothing$. Toutes les faces sont inférieures à 7 : l’évènement est certain, de probabilité $\res1$.
\end{corrige}\end{blocLivre}
\begin{blocLivre}{cor-11-5}\begin{corrige}{5}{Une fréquence expérimentale}\demarche{Ne pas confondre résultat observé et modèle.}
La fréquence est $46/100=\res{0{,}46}$. Cet écart à 0,5 peut se produire avec une pièce équilibrée ; ce seul résultat n’est pas une preuve. La fréquence dépend de l’expérience réalisée.
\end{corrige}\end{blocLivre}
\begin{blocLivre}{cor-11-6}\begin{corrige}{6}{La relation de la réunion}\demarche{Isoler $P(A\cup B)$ dans la relation.}
$P(A)=3/6=1/2$, $P(B)=2/6=1/3$ et $P(A\cap B)=1/6$. La relation donne $P(A\cup B)=1/2+1/3-1/6=\res{2/3}$. Contrôle : $A\cup B$ contient 4 issues sur 6, et $4/6=2/3$.
\end{corrige}\end{blocLivre}
\begin{blocLivre}{cor-11-7}\begin{corrige}{7}{Deux pièces}\demarche{Suivre toutes les branches de l’arbre.}
Les quatre issues équiprobables sont PP, PF, FP, FF. Exactement une face : PF et FP, d’où $\res{1/2}$. Au moins une face : PF, FP, FF, d’où $\res{3/4}$, ou $1-1/4$ par le contraire PP.
\label{fig-cor-11-7}\figureGeo[corrige]{proba-pieces}
\end{corrige}\end{blocLivre}
\begin{blocLivre}{cor-11-8}\begin{corrige}{8}{La somme de deux dés}\demarche{Compter des couples et non des sommes.}
Les 36 couples ordonnés sont équiprobables. Les six couples $(1;6),(2;5),(3;4),(4;3),(5;2),(6;1)$ conviennent : $\res{1/6}$. Les sommes elles-mêmes ne sont pas équiprobables.
\label{fig-cor-11-8}\figureGeo[corrige]{proba-des}
\end{corrige}\end{blocLivre}
\begin{blocLivre}{cor-11-9}\begin{corrige}{9}{Un « ou » exclusif}\demarche{Retirer de la réunion les issues communes.}
$A\cup B=\{2;4;5;6\}$ et $A\cap B=\{4;6\}$. Les issues cherchées sont 2 et 5 : probabilité $2/6=\res{1/3}$. Par le calcul : $P(A\cup B)-P(A\cap B)=4/6-2/6$.
\end{corrige}\end{blocLivre}
\begin{blocLivre}{cor-11-10}\begin{corrige}{10}{Au moins un six}\demarche{Utiliser la relation de la réunion pour éviter un double comptage.}
« Au moins un six » est $A\cup B$. $P(A)=P(B)=6/36$ et $A\cap B$ ne contient que $(6;6)$. Donc $P(A\cup B)=6/36+6/36-1/36=\res{11/36}$. Additionner $1/6+1/6$ aurait compté deux fois le couple $(6;6)$.
\label{fig-cor-11-10}\figureGeo[corrige]{proba-des}
\end{corrige}\end{blocLivre}
\begin{blocLivre}{cor-11-11}\begin{corrige}{11}{Deux tirages avec remise}\demarche{Numéroter les boules pour obtenir des issues équiprobables.}
Le tableau compte $5\times5=25$ couples équiprobables. Deux rouges : $3\times3=9$ couples, d’où $\res{9/25}$. Deux bleues : 4 couples. Même couleur : $9+4=13$ couples, d’où $\res{13/25}$, car les deux évènements sont incompatibles.
\label{fig-cor-11-11}\figureGeo[corrige]{proba-tableau-boules}
\end{corrige}\end{blocLivre}
\begin{blocLivre}{cor-11-12}\begin{corrige}{12}{Un jeu à points}\demarche{Distinguer gain brut, mise et bilan net.}
Le gain brut moyen vaut $8/6+4/6=\res2$ points. Il égale la mise : le bilan net moyen théorique est 0. Cela ne garantit ni un remboursement ni un gain à chaque partie ; un résultat individuel peut être $-2$, 2 ou 6 points nets.
\end{corrige}\end{blocLivre}
\begin{blocLivre}{cor-11-13}\begin{corrige}{13}{Un effectif inconnu}\demarche{Traduire le total avant le produit en croix.}
On écrit $6/(6+b)=2/5$. Les produits en croix donnent $30=12+2b$, donc $\res{b=9}$. Avec 15 boules, $6/15=2/5$, ce qui confirme le résultat.
\end{corrige}\end{blocLivre}
\begin{blocLivre}{cor-11-14}\begin{corrige}{14}{Un code à deux chiffres}\demarche{Préciser si la répétition est autorisée, puis utiliser la relation de la réunion.}
Il y a $4\times4=16$ codes équiprobables, dont $4\times3=12$ à chiffres distincts : $\res{3/4}$. Ensuite $P(\text{pair})=8/16$, $P(\text{début }0)=4/16$ et l’intersection contient 00 et 02 : $P=8/16+4/16-2/16=\res{5/8}$.
\end{corrige}\end{blocLivre}
\begin{blocLivre}{cor-11-15}\begin{corrige}{15}{Estimer sans promettre}\demarche{Comparer les écarts absolus sans en faire une règle monotone.}
Les fréquences sont 0,327 et 0,3354. Leurs écarts à $1/3$ valent environ 0,00633 et 0,00207. La deuxième est plus proche, ce qui illustre la stabilisation des fréquences. Aucune amélioration à chaque étape n’est garantie : un prochain succès ou échec peut augmenter l’écart.
\end{corrige}\end{blocLivre}

\clearpage\pagestyle{otMain}\fontsize{11}{15}\selectfont\corrigesChapitre{12}{Pythagore, cercle et droite des milieux}
\begin{blocLivre}{cor-12-1}\begin{corrige}{1}{Reconnaître l’hypoténuse}\demarche{Repérer d’abord le sommet de l’angle droit.}
Le côté opposé à A est $\res{BC}$. L’égalité est $\res{BC^2=AB^2+AC^2}$. Le nom du sommet de l’angle droit détermine le bon côté.
\label{fig-cor-12-1}\figureGeo[corrige]{pyth-01}
\end{corrige}\end{blocLivre}
\begin{blocLivre}{cor-12-2}\begin{corrige}{2}{Un triangle classique}\demarche{Additionner les carrés pour trouver l’hypoténuse.}
Pythagore donne $BC^2=6^2+8^2=100$. Comme BC est positive, $\res{BC=10\text{ cm}}$.
\label{fig-cor-12-2}\figureGeo[corrige]{pyth-02}
\end{corrige}\end{blocLivre}
\begin{blocLivre}{cor-12-3}\begin{corrige}{3}{Le côté manquant}\demarche{Isoler le carré inconnu avant d’extraire la racine.}
$AC^2=13^2-5^2=144$, donc $\res{AC=12\text{ cm}}$. AC est inférieure à l’hypoténuse 13 cm.
\label{fig-cor-12-3}\figureGeo[corrige]{pyth-03}
\end{corrige}\end{blocLivre}
\begin{blocLivre}{cor-12-4}\begin{corrige}{4}{Garder une racine}\demarche{Arrondir uniquement le résultat final.}
$BC^2=16+49=65$, donc $\res{BC=\sqrt{65}\text{ cm}}\approx\res{8{,}1\text{ cm}}$. L’arrondi n’est pas une égalité exacte.
\label{fig-cor-12-4}\figureGeo[corrige]{pyth-04}
\end{corrige}\end{blocLivre}
\begin{blocLivre}{cor-12-5}\begin{corrige}{5}{Le centre du cercle circonscrit}\demarche{Relier l’angle droit au cercle de diamètre l’hypoténuse.}
Comme ABC est rectangle en A, A appartient au cercle de diamètre [BC]. Le centre est donc $\res{O\text{, milieu de }[BC]}$, le rayon vaut $\res{5\text{ cm}}$ et $\res{OA=5\text{ cm}}$ : la médiane issue de l’angle droit mesure la moitié de l’hypoténuse.
\label{fig-cor-12-5}\figureGeo[corrige]{cercle-circonscrit}
\end{corrige}\end{blocLivre}
\begin{blocLivre}{cor-12-6}\begin{corrige}{6}{Tester la réciproque}\demarche{Comparer les deux membres séparément.}
BC est le plus grand côté. $15^2=225$ et $9^2+12^2=225$. Par la réciproque du théorème de Pythagore, ABC est $\res{\text{rectangle en A}}$.
\label{fig-cor-12-6}\figureGeo[corrige]{pyth-05}
\end{corrige}\end{blocLivre}
\begin{blocLivre}{cor-12-7}\begin{corrige}{7}{Utiliser la contraposée}\demarche{Comparer exactement, sans se fier au dessin.}
$BC^2=81$ alors que $AB^2+AC^2=25+49=74$. Ces nombres diffèrent ; d’après la $\res{\text{contraposée du théorème de Pythagore}}$, ABC $\res{\text{n’est pas rectangle}}$.
\label{fig-cor-12-7}\figureGeo[corrige]{pyth-06}
\end{corrige}\end{blocLivre}
\begin{blocLivre}{cor-12-8}\begin{corrige}{8}{La droite des milieux}\demarche{Citer les deux milieux avant d’appliquer le théorème.}
I et J sont les milieux de deux côtés : d’après le théorème de la droite des milieux, $\res{(IJ)\parallel(BC)}$ et $IJ=BC/2=\res{4{,}5\text{ cm}}$.
\label{fig-cor-12-8}\figureGeo[corrige]{milieux-parallele}
\end{corrige}\end{blocLivre}
\begin{blocLivre}{cor-12-9}\begin{corrige}{9}{Un milieu retrouvé}\demarche{Identifier le théorème dont la conclusion est un milieu.}
La droite passe par le milieu de [AB] et elle est parallèle à (BC) : d’après le deuxième théorème des milieux, J est le milieu de [AC]. Donc $\res{AJ=3{,}7\text{ cm}}$.
\label{fig-cor-12-9}\figureGeo[corrige]{milieux-milieu}
\end{corrige}\end{blocLivre}
\begin{blocLivre}{cor-12-10}\begin{corrige}{10}{Un angle à identifier}\demarche{Nommer le sommet opposé au plus long côté.}
Le plus grand côté est AB. $AB^2=100=36+64=AC^2+BC^2$. La réciproque du théorème de Pythagore donne un angle droit $\res{\text{en C}}$, opposé à AB. La position du dessin ne détermine pas l’angle droit.
\label{fig-cor-12-10}\figureGeo[corrige]{pyth-10}
\end{corrige}\end{blocLivre}
\begin{blocLivre}{cor-12-11}\begin{corrige}{11}{Une échelle contre un mur}\demarche{Justifier l’angle droit par le mur et le sol.}
ABC est rectangle en A, car le mur est perpendiculaire au sol. $AC^2=5^2-1{,}4^2=23{,}04$, donc $\res{AC=4{,}8\text{ m}}$. La hauteur reste inférieure à la longueur de l’échelle.
\label{fig-cor-12-11}\figureGeo[corrige]{pyth-07}
\end{corrige}\end{blocLivre}
\begin{blocLivre}{cor-12-12}\begin{corrige}{12}{Un triangle dans un demi-cercle}\demarche{Utiliser le cercle pour obtenir l’angle droit, puis Pythagore.}
A appartient au cercle de diamètre [BC] et diffère de B et C : ABC est $\res{\text{rectangle en A}}$. Le diamètre vaut $BC=13$ cm. Pythagore donne $AC^2=13^2-5^2=144$, donc $\res{AC=12\text{ cm}}$.
\label{fig-cor-12-12}\figureGeo[corrige]{demi-cercle}
\end{corrige}\end{blocLivre}
\begin{blocLivre}{cor-12-13}\begin{corrige}{13}{Un rectangle sans équerre}\demarche{Utiliser une propriété caractéristique du rectangle par ses diagonales.}
Les diagonales [AC] et [BD] ont le même milieu O : ABCD est un parallélogramme. Elles ont la même longueur, celle d’un diamètre : un parallélogramme dont les diagonales ont même longueur est $\res{\text{un rectangle}}$. On n’a utilisé que le compas et la règle : $\res{\text{aucun angle droit n’est tracé}}$.
\label{fig-cor-12-13}\figureGeo[corrige]{rectangle-compas}
\end{corrige}\end{blocLivre}
\begin{blocLivre}{cor-12-14}\begin{corrige}{14}{Une hauteur de triangle}\demarche{Justifier la perpendicularité avant Pythagore.}
Dans un triangle isocèle, la médiane [CH] est aussi une hauteur : ACH est rectangle en H. AH = 5 cm et $CH^2=13^2-5^2=144$, donc CH = 12 cm. L’aire vaut $10\times12/2=\res{60\text{ cm}^2}$.
\label{fig-cor-12-14}\figureGeo[corrige]{pyth-iso}
\end{corrige}\end{blocLivre}
\begin{blocLivre}{cor-12-15}\begin{corrige}{15}{Le théorème de Varignon}\demarche{Utiliser la diagonale [AC] comme intermédiaire dans deux triangles.}
Dans le triangle ABC, I et J sont des milieux : $(IJ)\parallel(AC)$ et $IJ=AC/2$. Dans le triangle ACD, L et K sont des milieux : $(LK)\parallel(AC)$ et $LK=AC/2$. Les côtés opposés [IJ] et [LK] sont donc parallèles et de même longueur : IJKL est $\res{\text{un parallélogramme}}$, quelle que soit la forme de ABCD.
\label{fig-cor-12-15}\figureGeo[corrige]{varignon}
\end{corrige}\end{blocLivre}

\clearpage\pagestyle{otMain}\fontsize{11}{15}\selectfont\corrigesChapitre{13}{Thalès, réciproque et contraposée}
\begin{blocLivre}{cor-13-1}\begin{corrige}{1}{Écrire les rapports}\demarche{Garder la même correspondance des sommets.}
Les alignements et le parallélisme permettent d’appliquer le théorème de Thalès : $\res{AM/AB=AN/AC=MN/BC}$. Les côtés du petit triangle sont tous au numérateur.
\label{fig-cor-13-1}\figureGeo[corrige]{thales-01}
\end{corrige}\end{blocLivre}
\begin{blocLivre}{cor-13-2}\begin{corrige}{2}{Une longueur réduite}\demarche{Multiplier la grande longueur par le rapport de réduction.}
Thalès donne $AN/10=3/5$. Ainsi $\res{AN=6\text{ cm}}$. Cette longueur est inférieure à AC, conformément à l’emboîtement.
\label{fig-cor-13-2}\figureGeo[corrige]{thales-02}
\end{corrige}\end{blocLivre}
\begin{blocLivre}{cor-13-3}\begin{corrige}{3}{La petite base}\demarche{Choisir le rapport portant l’inconnue.}
$MN/15=4/10$, donc $\res{MN=6\text{ cm}}$. Le rapport des longueurs est 0,4.
\label{fig-cor-13-3}\figureGeo[corrige]{thales-03}
\end{corrige}\end{blocLivre}
\begin{blocLivre}{cor-13-4}\begin{corrige}{4}{Retrouver le grand côté}\demarche{Une longueur agrandie se trouve en divisant par le rapport.}
$6/AB=4/10$. Le produit en croix donne $4AB=60$, donc $\res{AB=15\text{ cm}}$.
\label{fig-cor-13-4}\figureGeo[corrige]{thales-04}
\end{corrige}\end{blocLivre}
\begin{blocLivre}{cor-13-5}\begin{corrige}{5}{De la partie au tout}\demarche{Reconstituer le côté entier avant Thalès.}
Comme M est entre A et B, AB = 3 + 2 = 5 cm. Puis $AN=12\times3/5=\res{7{,}2\text{ cm}}$. Employer MB au dénominateur aurait associé les mauvais côtés.
\label{fig-cor-13-5}\figureGeo[corrige]{thales-05}
\end{corrige}\end{blocLivre}
\begin{blocLivre}{cor-13-6}\begin{corrige}{6}{Démontrer un parallélisme}\demarche{Vérifier les conditions géométriques avant les rapports.}
$AM/AB=4/10=2/5$ et $AN/AC=6/15=2/5$. Les alignements dans le même ordre sont donnés. La réciproque du théorème de Thalès donne $\res{(MN)\parallel(BC)}$.
\label{fig-cor-13-6}\figureGeo[corrige]{thales-06}
\end{corrige}\end{blocLivre}
\begin{blocLivre}{cor-13-7}\begin{corrige}{7}{Des droites non parallèles}\demarche{Comparer exactement, sans arrondir les quotients.}
$3/8\ne5/12$, car $3\times12=36$ et $5\times8=40$. D’après la contraposée du théorème de Thalès, les droites (MN) et (BC) $\res{\text{ne sont pas parallèles}}$.
\label{fig-cor-13-7}\figureGeo[corrige]{thales-07}
\end{corrige}\end{blocLivre}
\begin{blocLivre}{cor-13-8}\begin{corrige}{8}{Le papillon}\demarche{Conserver les longueurs positives dans les rapports.}
Thalès s’applique aussi au papillon : $AN/8=3/5$. Donc $\res{AN=4{,}8\text{ cm}}$. Les longueurs restent positives, même si les points sont de part et d’autre de A.
\label{fig-cor-13-8}\figureGeo[corrige]{thales-08}
\end{corrige}\end{blocLivre}
\begin{blocLivre}{cor-13-9}\begin{corrige}{9}{La distance entre deux points}\demarche{Lire l’ordre des points.}
A est entre les points de chaque paire : $\res{MB=3+5=8\text{ cm}}$ et $\res{NC=4{,}8+8=12{,}8\text{ cm}}$. Une soustraction correspondrait à une autre disposition.
\label{fig-cor-13-9}\figureGeo[corrige]{thales-09}
\end{corrige}\end{blocLivre}
\begin{blocLivre}{cor-13-10}\begin{corrige}{10}{Une réciproque en papillon}\demarche{Contrôler l’ordre des points, puis comparer les rapports.}
$AM/AB=2{,}4/4=0{,}6$ et $AN/AC=3/5=0{,}6$. Les points sont alignés dans le même ordre, de part et d’autre de A. La réciproque du théorème de Thalès donne $\res{(MN)\parallel(BC)}$.
\label{fig-cor-13-10}\figureGeo[corrige]{thales-10}
\end{corrige}\end{blocLivre}
\begin{blocLivre}{cor-13-11}\begin{corrige}{11}{Deux triangles de même aire}\demarche{Repérer une base commune et deux hauteurs égales.}
B et C sont sur une droite parallèle à (MN) : ils sont à la même distance de (MN). Les triangles MNB et MNC ont la même base [MN] et la même hauteur : $\res{\text{même aire}}$. Ensuite $\mathcal A(AMC)=\mathcal A(AMN)+\mathcal A(MNC)=\mathcal A(AMN)+\mathcal A(MNB)=\mathcal A(ANB)$. C’est l’idée de la preuve de Thalès par les aires.
\label{fig-cor-13-11}\figureGeo[corrige]{thales-aires}
\end{corrige}\end{blocLivre}
\begin{blocLivre}{cor-13-12}\begin{corrige}{12}{L’ombre d’un arbre}\demarche{Identifier le sommet commun et les deux parallèles.}
Le bâton et l’arbre sont verticaux, donc parallèles. Les pieds sont alignés avec S, ainsi que les sommets. Le théorème de Thalès dans les triangles emboîtés de sommet S donne $h/1{,}5=12/2$, d’où $\res{h=9\text{ m}}$. Les hypothèses de verticalité et de sol horizontal justifient le modèle.
\label{fig-cor-13-12}\figureGeo[corrige]{thales-ombre}
\end{corrige}\end{blocLivre}
\begin{blocLivre}{cor-13-13}\begin{corrige}{13}{Une équation dans Thalès}\demarche{Exprimer les côtés entiers avant de résoudre.}
AB = $4+x$ et AC = 9. Thalès donne $4/(4+x)=6/9=2/3$. Ainsi $12=8+2x$, donc $\res{x=2\text{ cm}}$. La valeur est positive et redonne $4/6=6/9$.
\label{fig-cor-13-13}\figureGeo[corrige]{thales-13}
\end{corrige}\end{blocLivre}
\begin{blocLivre}{cor-13-14}\begin{corrige}{14}{Calculer une largeur inaccessible}\demarche{Séparer la preuve du parallélisme et le calcul.}
M et N sont les milieux de deux côtés : par le théorème de la droite des milieux, ou par la réciproque de Thalès avec $AM/AB=AN/AC=1/2$, on a $(MN)\parallel(BC)$. Puis $MN/BC=1/2$, donc $\res{BC=14{,}8\text{ m}}$.
\label{fig-cor-13-14}\figureGeo[corrige]{thales-milieux}
\end{corrige}\end{blocLivre}
\begin{blocLivre}{cor-13-15}\begin{corrige}{15}{Un contrôle critique}\demarche{Vérifier que chaque terme est un côté correspondant.}
AM correspond à AN, donc $AN/16=3/8$, d’où AN = 6 cm. Puis $NC=16-6=\res{10\text{ cm}}$. NC n’est pas un côté du petit triangle AMN.
\label{fig-cor-13-15}\figureGeo[corrige]{thales-15}
\end{corrige}\end{blocLivre}

\clearpage\pagestyle{otMain}\fontsize{11}{15}\selectfont\corrigesChapitre{14}{Trigonométrie du triangle rectangle}
\begin{blocLivre}{cor-14-1}\begin{corrige}{1}{Nommer les côtés}\demarche{Commencer par l’hypoténuse.}
Le côté opposé est $\res{AC}$, l’adjacent est $\res{AB}$ et l’hypoténuse est $\res{BC}$. AB touche B mais n’est pas l’hypoténuse.
\label{fig-cor-14-1}\figureGeo[corrige]{trigo-01}
\end{corrige}\end{blocLivre}
\begin{blocLivre}{cor-14-2}\begin{corrige}{2}{Choisir le cosinus}\demarche{Écrire les noms des côtés avant les nombres.}
$\cos\widehat B=AB/BC=6/10=\res{0{,}6}$. Le rapport utilise l’adjacent et l’hypoténuse.
\label{fig-cor-14-2}\figureGeo[corrige]{trigo-02}
\end{corrige}\end{blocLivre}
\begin{blocLivre}{cor-14-3}\begin{corrige}{3}{Choisir le sinus}\demarche{Repérer le côté qui ne touche pas l’angle B.}
$\sin\widehat B=AC/BC=8/10=\res{0{,}8}$. AC est opposé à $\widehat B$.
\label{fig-cor-14-3}\figureGeo[corrige]{trigo-03}
\end{corrige}\end{blocLivre}
\begin{blocLivre}{cor-14-4}\begin{corrige}{4}{Choisir la tangente}\demarche{La tangente ne fait pas intervenir l’hypoténuse.}
$\tan\widehat B=AC/AB=8/6=\res{4/3}$. Une tangente peut dépasser 1, contrairement au sinus ou au cosinus d’un angle aigu.
\label{fig-cor-14-4}\figureGeo[corrige]{trigo-04}
\end{corrige}\end{blocLivre}
\begin{blocLivre}{cor-14-5}\begin{corrige}{5}{Une première longueur}\demarche{Isoler la longueur dans le rapport.}
$AC=BC\sin30^\circ=12\times0{,}5=\res{6\text{ cm}}$. Le côté opposé est la moitié de l’hypoténuse pour cet angle.
\label{fig-cor-14-5}\figureGeo[corrige]{trigo-05}
\end{corrige}\end{blocLivre}
\begin{blocLivre}{cor-14-6}\begin{corrige}{6}{Un côté adjacent}\demarche{Relier la précision à l’unité de travail.}
$AB=9\cos40^\circ\approx6{,}8944$ cm, donc $\res{AB\approx6{,}9\text{ cm}}$ au millimètre. Le millimètre représente 0,1 cm.
\label{fig-cor-14-6}\figureGeo[corrige]{trigo-06}
\end{corrige}\end{blocLivre}
\begin{blocLivre}{cor-14-7}\begin{corrige}{7}{Une hypoténuse}\demarche{Résoudre l’égalité avant de calculer.}
$\cos50^\circ=7/BC$, donc $BC=7/\cos50^\circ\approx\res{10{,}9\text{ cm}}$. Multiplier par le cosinus aurait donné un côté trop court.
\label{fig-cor-14-7}\figureGeo[corrige]{trigo-07}
\end{corrige}\end{blocLivre}
\begin{blocLivre}{cor-14-8}\begin{corrige}{8}{Un angle par cosinus}\demarche{Utiliser la fonction inverse en mode degrés.}
$\cos\widehat B=0{,}6$. La fonction inverse donne $\widehat B\approx53{,}13^\circ$, donc $\res{\widehat B\approx53^\circ}$. L’angle $\widehat C$ vaut alors environ 37 degrés.
\label{fig-cor-14-8}\figureGeo[corrige]{trigo-08}
\end{corrige}\end{blocLivre}
\begin{blocLivre}{cor-14-9}\begin{corrige}{9}{Un angle par tangente}\demarche{Choisir le rapport contenant les deux données.}
$\tan\widehat B=5/12$, donc $\widehat B\approx\res{22{,}6^\circ}$. L’opposé est inférieur à l’adjacent : un angle inférieur à 45 degrés est cohérent.
\label{fig-cor-14-9}\figureGeo[corrige]{trigo-09}
\end{corrige}\end{blocLivre}
\begin{blocLivre}{cor-14-10}\begin{corrige}{10}{Cosinus et sinus liés}\demarche{Écrire les rapports, puis utiliser l’égalité de Pythagore.}
$(\cos\widehat B)^2+(\sin\widehat B)^2=\dfrac{AB^2}{BC^2}+\dfrac{AC^2}{BC^2}=\dfrac{AB^2+AC^2}{BC^2}=1$ d’après Pythagore. Ensuite $(\sin\widehat B)^2=1-0{,}36=0{,}64$ ; un sinus d’angle aigu est positif, donc $\res{\sin\widehat B=0{,}8}$.
\label{fig-cor-14-10}\figureGeo[corrige]{trigo-10}
\end{corrige}\end{blocLivre}
\begin{blocLivre}{cor-14-11}\begin{corrige}{11}{La hauteur d’un bâtiment}\demarche{Séparer hauteur au-dessus de la visée et hauteur totale.}
Au-dessus de l’appareil, la hauteur vaut $25\tan32^\circ\approx15{,}6217$ m. En ajoutant 1,6 m, la hauteur totale est $\res{17{,}2\text{ m}}$ au décimètre. Oublier la hauteur de visée sous-estimerait le résultat.
\label{fig-cor-14-11}\figureGeo[corrige]{trigo-batiment}
\end{corrige}\end{blocLivre}
\begin{blocLivre}{cor-14-12}\begin{corrige}{12}{Une rampe et son angle}\demarche{Distinguer pourcentage de pente et mesure angulaire.}
$\widehat B=\tan^{-1}(0{,}8/10)\approx\res{4{,}6^\circ}$. La pente est $\res{8\,\%}$. 8 \% ne signifie donc pas un angle de 8 degrés.
\label{fig-cor-14-12}\figureGeo[corrige]{trigo-12}
\end{corrige}\end{blocLivre}
\begin{blocLivre}{cor-14-13}\begin{corrige}{13}{Un câble}\demarche{Calculer les deux projections de l’hypoténuse.}
$AC=18\sin65^\circ\approx\res{16{,}3\text{ m}}$ et $AB=18\cos65^\circ\approx\res{7{,}6\text{ m}}$. Les calculs non arrondis vérifient Pythagore.
\label{fig-cor-14-13}\figureGeo[corrige]{trigo-13}
\end{corrige}\end{blocLivre}
\begin{blocLivre}{cor-14-14}\begin{corrige}{14}{Un toit symétrique}\demarche{Exploiter la symétrie pour obtenir un triangle rectangle.}
La hauteur partage la base en deux segments de 4 m. Elle vaut $4\tan35^\circ\approx\res{2{,}80\text{ m}}$. Un versant mesure $4/\cos35^\circ\approx\res{4{,}88\text{ m}}$. On utilise 4 m, pas la base entière.
\label{fig-cor-14-14}\figureGeo[corrige]{trigo-toit}
\end{corrige}\end{blocLivre}
\begin{blocLivre}{cor-14-15}\begin{corrige}{15}{Diagnostiquer une erreur}\demarche{Contrôler les côtés du rapport, pas seulement le résultat.}
AB est adjacent et AC opposé à $\widehat B$ ; leur rapport est la tangente. $\tan38^\circ=AC/7$, donc $\res{AC=7\tan38^\circ\approx5{,}5\text{ cm}}$. Le sinus exigerait l’hypoténuse au dénominateur.
\label{fig-cor-14-15}\figureGeo[corrige]{trigo-15}
\end{corrige}\end{blocLivre}

\clearpage\pagestyle{otMain}\fontsize{11}{15}\selectfont\corrigesChapitre{15}{Translations et vecteurs}
\begin{blocLivre}{cor-15-1}\begin{corrige}{1}{Symétrie et demi-tour}\demarche{Reporter les distances sur le quadrillage, point par point.}
Par la symétrie, chaque point est reporté perpendiculairement à (d), de l’autre côté et à la même distance : A et C sont à 6 carreaux de (d), B à 3 carreaux ; on obtient $\res{A'B'C'}$, tracé en vert. Par le demi-tour, O est le milieu de chaque segment point-image : on reproduit au-delà de O le déplacement de A vers O, soit 4 carreaux à droite et 3 en haut, pour placer A’’ ; de même pour $\res{B''\text{ et }C''}$. Les deux transformations $\res{\text{conservent les longueurs}}$.
\label{fig-cor-15-1}\figureGeo[corrige]{vect-01}
\end{corrige}\end{blocLivre}
\begin{blocLivre}{cor-15-2}\begin{corrige}{2}{Caractériser un vecteur}\demarche{Lire le déplacement horizontal et vertical, puis le reproduire.}
La direction est celle de la droite (AB), le sens va de A vers B. Le déplacement est de 4 carreaux vers la droite et 3 vers le haut : $AB=\sqrt{4^2+3^2}=\res{5\text{ cm}}$. On applique le même déplacement à partir de C : D est 4 carreaux à droite et 3 au-dessus de C.
\label{fig-cor-15-2}\figureGeo[corrige]{vect-02}
\end{corrige}\end{blocLivre}
\begin{blocLivre}{cor-15-3}\begin{corrige}{3}{Des vecteurs égaux}\demarche{Comparer direction, sens et longueur, ou le déplacement sur le quadrillage.}
$\vect{CD}$ et $\vect{IJ}$ correspondent au même déplacement que $\vect{AB}$ (2 à droite, 1 en haut) : $\res{\vect{CD}=\vect{IJ}=\vect{AB}}$. $\vect{EF}$ a même direction et même longueur, mais un sens contraire : $\res{\vect{EF}=-\vect{AB}=\vect{BA}}$. $\vect{GH}$ n’a pas la même direction.
\label{fig-cor-15-3}\figureGeo[corrige]{vect-03}
\end{corrige}\end{blocLivre}
\begin{blocLivre}{cor-15-4}\begin{corrige}{4}{Image par une translation}\demarche{Reproduire le déplacement, puis citer la définition.}
On applique à C le déplacement de A vers B : 3 carreaux à droite, 1 en haut. D’après la définition de la translation, $\res{ABDC\text{ est un parallélogramme}}$ et $\vect{CD}=\vect{AB}$. Les côtés [AB] et [CD] sont parallèles et de même longueur.
\label{fig-cor-15-4}\figureGeo[corrige]{vect-04}
\end{corrige}\end{blocLivre}
\begin{blocLivre}{cor-15-5}\begin{corrige}{5}{Vecteur nul et opposé}\demarche{Utiliser les définitions, puis la relation de Chasles.}
$\res{\vect{AA}=\vec0}$, le vecteur nul. L’opposé de $\vect{MN}$ est $\res{\vect{NM}}$ : même direction, même longueur, sens contraire. Par Chasles, $\vect{MN}+\vect{NM}=\vect{MM}=\res{\vec0}$.
\label{fig-cor-15-5}\figureGeo[corrige]{vect-05}
\end{corrige}\end{blocLivre}
\begin{blocLivre}{cor-15-6}\begin{corrige}{6}{Parallélogramme et égalités}\demarche{Suivre les côtés opposés dans le même sens.}
Dans ABCD, les côtés opposés donnent $\res{\vect{AB}=\vect{DC}}$ et $\res{\vect{AD}=\vect{BC}}$. En changeant le sens : $\res{\vect{BA}=\vect{CD}}$ et $\res{\vect{CB}=\vect{DA}}$. L’ordre des lettres suit le tour du parallélogramme.
\label{fig-cor-15-6}\figureGeo[corrige]{vect-06}
\end{corrige}\end{blocLivre}
\begin{blocLivre}{cor-15-7}\begin{corrige}{7}{Démontrer un parallélogramme}\demarche{Remplacer les lettres dans la propriété sans changer leur ordre.}
D’après la propriété du cours, $\vect{AB}=\vect{CD}$ équivaut à « ABDC est un parallélogramme ». Avec A = E, B = F, C = H et D = G, on obtient que $\res{EFGH\text{ est un parallélogramme}}$. Sur le quadrillage, les deux flèches correspondent au déplacement 4 à droite, 1 en haut.
\label{fig-cor-15-7}\figureGeo[corrige]{vect-07}
\end{corrige}\end{blocLivre}
\begin{blocLivre}{cor-15-8}\begin{corrige}{8}{La relation de Chasles}\demarche{Enchaîner les vecteurs bout à bout : l’extrémité de l’un est l’origine du suivant.}
$\res{\vect{AB}+\vect{BC}=\vect{AC}}$, comme sur la figure. $\vect{MN}+\vect{NP}+\vect{PQ}=\res{\vect{MQ}}$. En changeant l’ordre des termes, $\vect{CA}+\vect{AB}=\res{\vect{CB}}$. Enfin $\vect{AB}+\vect{BA}=\res{\vec0}$.
\label{fig-cor-15-8}\figureGeo[corrige]{vect-08}
\end{corrige}\end{blocLivre}
\begin{blocLivre}{cor-15-9}\begin{corrige}{9}{Construire une somme}\demarche{Mettre les représentants bout à bout.}
On trace $\vect{OP}=\vec u$ (3 à droite, 1 en haut), puis $\vect{PS}=\vec v$ (1 à droite, 2 en haut). Alors $\vec u+\vec v=\vect{OS}$, soit 4 à droite et 3 en haut : $OS=\sqrt{16+9}=\res{5\text{ cm}}$. Les longueurs de $\vec u$ et $\vec v$, $\sqrt{10}$ et $\sqrt5$, ont une somme différente.
\label{fig-cor-15-9}\figureGeo[corrige]{vect-09}
\end{corrige}\end{blocLivre}
\begin{blocLivre}{cor-15-10}\begin{corrige}{10}{Deux translations successives}\demarche{Remplacer chaque déplacement par un vecteur, puis additionner.}
$\vect{MM'}=\vect{AB}$ et $\vect{M'M''}=\vect{BC}$. Par Chasles, $\vect{MM''}=\vect{MM'}+\vect{M'M''}=\vect{AB}+\vect{BC}=\vect{AC}$ : c’est la $\res{\text{translation de vecteur }\vect{AC}}$. On le vérifie sur le quadrillage : 4 à droite, 3 en haut.
\label{fig-cor-15-10}\figureGeo[corrige]{vect-10}
\end{corrige}\end{blocLivre}
\begin{blocLivre}{cor-15-11}\begin{corrige}{11}{La règle du parallélogramme}\demarche{Remplacer un vecteur par un vecteur égal, puis appliquer Chasles.}
ABCD est un parallélogramme, donc $\vect{AD}=\vect{BC}$. Ainsi $\vect{AB}+\vect{AD}=\vect{AB}+\vect{BC}$, et la relation de Chasles donne $\res{\vect{AB}+\vect{AD}=\vect{AC}}$. La diagonale issue de A représente la somme.
\label{fig-cor-15-11}\figureGeo[corrige]{vect-11}
\end{corrige}\end{blocLivre}
\begin{blocLivre}{cor-15-12}\begin{corrige}{12}{Milieu et vecteurs}\demarche{Traduire « milieu » en direction, sens et longueur.}
I est sur [AB] à égale distance de A et B : $\vect{AI}$ et $\vect{IB}$ ont même direction, même sens et même longueur, donc $\res{\vect{AI}=\vect{IB}}$. Par Chasles, $\vect{AB}=\vect{AI}+\vect{IB}=\vect{AI}+\vect{AI}=\res{2\vect{AI}}$. Si $\vect{AJ}=\vect{JB}$, alors AJBJ est un parallélogramme aplati : ses diagonales [AB] et [JJ] ont même milieu, donc $\res{J\text{ est le milieu de }[AB]}$.
\label{fig-cor-15-12}\figureGeo[corrige]{vect-12}
\end{corrige}\end{blocLivre}
\begin{blocLivre}{cor-15-13}\begin{corrige}{13}{Un drone}\demarche{Additionner les déplacements, puis calculer la longueur avec Pythagore.}
Le premier déplacement est $\vect{DE}$, le second $\vect{EF}$. Par Chasles, leur somme est $\res{\vect{DF}}$. Le triangle DEF est rectangle en E : $DF=\sqrt{60^2+80^2}=\res{100\text{ m}}$, alors que le drone a parcouru 140 m.
\label{fig-cor-15-13}\figureGeo[corrige]{vect-13}
\end{corrige}\end{blocLivre}
\begin{blocLivre}{cor-15-14}\begin{corrige}{14}{Une frise}\demarche{Compter les intervalles entre les motifs.}
Du motif 1 au motif 6, il y a cinq déplacements : $\vec u+\vec u+\vec u+\vec u+\vec u$, noté $5\vec u$, de longueur $\res{20\text{ cm}}$. Du motif 6 au motif 2, on revient de quatre motifs : $\res{-4\vec u}$, même direction, sens contraire, longueur $\res{16\text{ cm}}$.
\label{fig-cor-15-14}\figureGeo[corrige]{vect-frise}
\end{corrige}\end{blocLivre}
\begin{blocLivre}{cor-15-15}\begin{corrige}{15}{Démontrer avec des vecteurs}\demarche{Chaîner deux égalités de vecteurs, puis citer la propriété.}
ABCD est un parallélogramme, donc $\vect{DC}=\vect{AB}$. Par hypothèse $\vect{BE}=\vect{AB}$. Donc $\vect{DC}=\vect{BE}$, ce qui signifie que $\res{DCEB\text{ est un parallélogramme}}$, d’après la propriété du cours appliquée aux points D, C, B, E.
\label{fig-cor-15-15}\figureGeo[corrige]{vect-15}
\end{corrige}\end{blocLivre}

\clearpage\pagestyle{otMain}\fontsize{11}{15}\selectfont\corrigesChapitre{16}{Espace : sphère, boule et sections}
\begin{blocLivre}{cor-16-1}\begin{corrige}{1}{Un pavé}\demarche{Multiplier les trois dimensions, puis convertir.}
$V=8\times5\times3=\res{120\text{ cm}^3}$. Comme $1\text{ L}=1000\text{ cm}^3$, cela fait $\res{0{,}12\text{ L}}$.
\label{fig-cor-16-1}\figureGeo[corrige]{espace-01}
\end{corrige}\end{blocLivre}
\begin{blocLivre}{cor-16-2}\begin{corrige}{2}{Un cylindre}\demarche{Calculer l’aire du disque de base.}
$V=\pi\times3^2\times10=\res{90\pi\text{ cm}^3}\approx\res{283\text{ cm}^3}$. Le rayon est déjà fourni, il ne faut pas le diviser par deux.
\label{fig-cor-16-2}\figureGeo[corrige]{espace-02}
\end{corrige}\end{blocLivre}
\begin{blocLivre}{cor-16-3}\begin{corrige}{3}{Une pyramide}\demarche{Utiliser la hauteur perpendiculaire à la base.}
La base est un carré d’aire 36 cm$^2$. Le volume vaut $36\times9/3=\res{108\text{ cm}^3}$. La division par trois distingue la pyramide du prisme de même base et de même hauteur.
\label{fig-cor-16-3}\figureGeo[corrige]{espace-03}
\end{corrige}\end{blocLivre}
\begin{blocLivre}{cor-16-4}\begin{corrige}{4}{Un cône}\demarche{Passer du diamètre au rayon.}
Le rayon est 4 cm. $V=\pi\times4^2\times9/3=\res{48\pi\text{ cm}^3}$. Employer 8 comme rayon multiplierait le résultat par quatre.
\label{fig-cor-16-4}\figureGeo[corrige]{espace-04}
\end{corrige}\end{blocLivre}
\begin{blocLivre}{cor-16-5}\begin{corrige}{5}{Une boule}\demarche{Élever le rayon au cube, puis multiplier par $4\pi/3$.}
$V=\dfrac43\pi\times27=\res{36\pi\text{ cm}^3}\approx\res{113\text{ cm}^3}$. Un grand cercle a pour rayon 3 cm : sa longueur vaut $2\pi\times3=\res{6\pi\text{ cm}}$, environ 18,8 cm.
\label{fig-cor-16-5}\figureGeo[corrige]{espace-05}
\end{corrige}\end{blocLivre}
\begin{blocLivre}{cor-16-6}\begin{corrige}{6}{Sphère, boule et grands cercles}\demarche{Comparer chaque distance au centre avec le rayon.}
A, B et C sont à 4 cm de O : ils appartiennent à la $\res{\text{sphère}}$, donc aussi à la boule. M est à 2 cm de O : il est $\res{\text{dans la boule, pas sur la sphère}}$. N est à 6 cm : $\res{\text{à l’extérieur}}$. [AB] est un diamètre : $\res{AB=8\text{ cm}}$. Le grand cercle passant par A et B a pour longueur $\res{8\pi\text{ cm}}$.
\label{fig-cor-16-6}\figureGeo[corrige]{espace-sphere-vocab}
\end{corrige}\end{blocLivre}
\begin{blocLivre}{cor-16-7}\begin{corrige}{7}{Une pyramide dans un cube}\demarche{Repérer la base, la hauteur [EA] et les angles droits des faces.}
La base ABCD est un carré. EAB et EAD sont des triangles $\res{\text{rectangles isocèles en A}}$, moitiés de faces du cube. La droite (BC) est perpendiculaire à la face ABFE, donc à (BE) : EBC est $\res{\text{rectangle en B}}$ ; de même EDC est rectangle en D. Le volume vaut $36\times6/3=\res{72\text{ cm}^3}$, soit $\res{1/3}$ du cube de $216\text{ cm}^3$.
\label{fig-cor-16-7}\figureGeo[corrige]{espace-pyramide-cube}
\end{corrige}\end{blocLivre}
\begin{blocLivre}{cor-16-8}\begin{corrige}{8}{Une section de pavé}\demarche{Reconnaître une section parallèle à une face.}
La section est un $\res{\text{rectangle de 8 cm sur 5 cm}}$, comme la base. Son aire vaut $\res{40\text{ cm}^2}$. La hauteur de coupe ne modifie pas cette section parallèle dans un pavé.
\label{fig-cor-16-8}\figureGeo[corrige]{espace-section-pave}
\end{corrige}\end{blocLivre}
\begin{blocLivre}{cor-16-9}\begin{corrige}{9}{Sections d’un cylindre}\demarche{Identifier la direction du plan par rapport à l’axe.}
Perpendiculairement à l’axe, la section est un $\res{\text{disque de rayon 3 cm}}$, d’aire $\res{9\pi\text{ cm}^2}$. Un plan contenant l’axe donne un $\res{\text{rectangle de 6 cm sur 10 cm}}$ : sa largeur est un diamètre de la base. Son aire vaut $\res{60\text{ cm}^2}$.
\label{fig-cor-16-9}\figureGeo[corrige]{espace-section-cylindre}
\end{corrige}\end{blocLivre}
\begin{blocLivre}{cor-16-10}\begin{corrige}{10}{Une section de sphère}\demarche{Passer par le triangle rectangle de la section.}
Pour M sur le cercle, OHM est rectangle en H. $HM^2=OM^2-OH^2=25-9=16$, donc $\res{HM=4\text{ cm}}$. Le disque a pour aire $\res{16\pi\text{ cm}^2}$.
\label{fig-cor-16-10}\figureGeo[corrige]{espace-section-sphere}
\end{corrige}\end{blocLivre}
\begin{blocLivre}{cor-16-11}\begin{corrige}{11}{Un réservoir assemblé}\demarche{Additionner les volumes puis convertir.}
Le cylindre vaut $2\pi$ m$^3$, la demi-boule $\frac12\times\frac43\pi=\frac{2\pi}3$ m$^3$. Le total est $8\pi/3$ m$^3$, soit $8000\pi/3$ L, environ $\res{8378\text{ L}}$. Les deux parties ne se recouvrent qu’à leur frontière.
\label{fig-cor-16-11}\figureGeo[corrige]{espace-reservoir}
\end{corrige}\end{blocLivre}
\begin{blocLivre}{cor-16-12}\begin{corrige}{12}{Une boule creusée}\demarche{Soustraire les volumes du plein et du creux.}
On soustrait les volumes : $V=\frac43\pi(5^3-2^3)=\frac43\pi\times117=\res{156\pi\text{ cm}^3}$. Il faut soustraire des cubes, pas calculer le volume d’une boule de rayon $5-2$.
\label{fig-cor-16-12}\figureGeo[corrige]{espace-cavite}
\end{corrige}\end{blocLivre}
\begin{blocLivre}{cor-16-13}\begin{corrige}{13}{Un cône coupé}\demarche{Utiliser les dimensions proportionnelles de la section, puis soustraire.}
La section est une réduction de la base de coefficient $6/12=1/2$ : un $\res{\text{disque de rayon 3 cm}}$. Le petit cône a pour rayon 3 cm et hauteur 6 cm : $V=\pi\times9\times6/3=\res{18\pi\text{ cm}^3}$. Le grand cône vaut $\pi\times36\times12/3=144\pi$ cm$^3$ ; le tronc vaut $\res{126\pi\text{ cm}^3}$.
\label{fig-cor-16-13}\figureGeo[corrige]{espace-cone-coupe}
\end{corrige}\end{blocLivre}
\begin{blocLivre}{cor-16-14}\begin{corrige}{14}{Une section de pyramide}\demarche{Calculer le coefficient de réduction à partir des hauteurs.}
La section est une réduction de la base de coefficient $3/9=1/3$ : un $\res{\text{carré de côté 2 cm}}$, d’aire $\res{4\text{ cm}^2}$. La petite pyramide a pour hauteur 3 cm : $V=4\times3/3=\res{4\text{ cm}^3}$.
\label{fig-cor-16-14}\figureGeo[corrige]{espace-section-pyramide}
\end{corrige}\end{blocLivre}
\begin{blocLivre}{cor-16-15}\begin{corrige}{15}{Le cylindre d’Archimède}\demarche{Écrire les deux volumes avec le même rayon, puis simplifier.}
Le cylindre vaut $\pi\times9\times6=54\pi$ cm$^3$ et la boule $36\pi$ cm$^3$. Le rapport vaut $36/54=\res{2/3}$. Pour un rayon $r$ : le cylindre vaut $\pi r^2\times2r=2\pi r^3$ et la boule $\frac43\pi r^3$, d’où $\res{\frac{4/3}{2}=\frac23}$ pour tout $r$. Archimède avait découvert ce résultat.
\label{fig-cor-16-15}\figureGeo[corrige]{espace-archimede}
\end{corrige}\end{blocLivre}

\clearpage\pagestyle{otMain}\fontsize{11}{15}\selectfont\corrigesChapitre{17}{Pensée informatique et algorithmique}
\begin{blocLivre}{cor-17-1}\begin{corrige}{1}{Suivre trois opérations}\demarche{Écrire chaque valeur dans l’ordre.}
x prend successivement 3, 8 et 16. Le bloc dire affiche $\res{16}$. Ajouter puis multiplier ne donne pas le même résultat que multiplier puis ajouter.
\label{fig-cor-17-1}\figureGeo[corrige]{script-01}
\end{corrige}\end{blocLivre}
\begin{blocLivre}{cor-17-2}\begin{corrige}{2}{Un nombre demandé}\demarche{Traduire les modifications successives.}
Après ajout, on a $x+4$ ; après multiplication, $3(x+4)$. Les résultats sont $\res{18}$ et $\res0$. Les parenthèses traduisent l’ordre des blocs.
\label{fig-cor-17-2}\figureGeo[corrige]{script-02}
\end{corrige}\end{blocLivre}
\begin{blocLivre}{cor-17-3}\begin{corrige}{3}{Une boucle bornée}\demarche{Repérer les limites de la boucle.}
Le bloc ajouter s’exécute quatre fois. Les valeurs sont 3, 6, 9, 12. La sortie est $\res{12}$. Le bloc dire placé après la boucle s’exécute une seule fois.
\label{fig-cor-17-3}\figureGeo[corrige]{script-03}
\end{corrige}\end{blocLivre}
\begin{blocLivre}{cor-17-4}\begin{corrige}{4}{Un test et deux branches}\demarche{Tester la valeur située à la frontière.}
Le test est $x\ge10$, donc il est faux pour 9 et vrai pour 10 et 14. Les sorties sont respectivement $\res{\text{petit, grand, grand}}$. La borne 10 conduit bien à la branche du test vrai.
\label{fig-cor-17-4}\figureGeo[corrige]{script-04}
\end{corrige}\end{blocLivre}
\begin{blocLivre}{cor-17-5}\begin{corrige}{5}{Doubler plusieurs fois}\demarche{Distinguer additions et multiplications répétées.}
Les trois tours donnent 4, 8, 16. La sortie est $\res{16=2\times2^3}$. Le script multiplie à chaque tour ; il n’ajoute pas 2 trois fois.
\label{fig-cor-17-5}\figureGeo[corrige]{script-05}
\end{corrige}\end{blocLivre}
\begin{blocLivre}{cor-17-6}\begin{corrige}{6}{Deux variables}\demarche{Faire un tableau des deux variables.}
Au départ $(s;n)=(0;1)$. Après les cinq tours : $(1;2),(3;3),(6;4),(10;5),(15;6)$. La sortie est $\res{15}$ : le script calcule la somme $1+2+3+4+5$. L’ajout à s précède l’augmentation de n.
\label{fig-cor-17-6}\figureGeo[corrige]{script-06}
\end{corrige}\end{blocLivre}
\begin{blocLivre}{cor-17-7}\begin{corrige}{7}{Répéter jusqu’à}\demarche{Vérifier le test avant chaque nouveau tour.}
En partant de 0, on obtient 6, 12, 18 puis 24. Le test $x\ge20$ devient vrai à 24 : $\res{4\text{ tours}}$ et sortie $\res{24}$. La boucle ne s’arrête pas à 20, qui n’est jamais atteint.
\label{fig-cor-17-7}\figureGeo[corrige]{script-07}
\end{corrige}\end{blocLivre}
\begin{blocLivre}{cor-17-8}\begin{corrige}{8}{Une condition composée}\demarche{Évaluer chaque condition simple, puis les combiner.}
La condition est vraie seulement si $x\ge10$ et $x\le20$ sont vraies toutes les deux. Les affichages sont $\res{\text{hors, dans, dans, hors de la plage}}$ ; les bornes 10 et 20 sont incluses. Avec « ou », tout nombre vérifierait l’une des deux conditions : le script afficherait toujours « dans la plage ».
\label{fig-cor-17-8}\figureGeo[corrige]{script-08}
\end{corrige}\end{blocLivre}
\begin{blocLivre}{cor-17-9}\begin{corrige}{9}{Énumérer des diviseurs}\demarche{Comprendre l’intervalle parcouru par la boucle.}
Le script teste d de 1 à 12 et affiche $\res{1,2,3,4,6,12}$. Ce sont tous les diviseurs positifs de 12. Commencer par 0 rendrait le calcul du reste impossible, car on ne divise pas par zéro.
\label{fig-cor-17-9}\figureGeo[corrige]{script-09}
\end{corrige}\end{blocLivre}
\begin{blocLivre}{cor-17-10}\begin{corrige}{10}{Le reste par soustractions}\demarche{Suivre r et q à chaque tour, en testant la condition avant le tour.}
Pour 23 et 5, r prend les valeurs 23, 18, 13, 8, 3 et q compte 4 soustractions : affichage $\res{4\text{ puis }3}$. Pour 3 et 5, le test $r<b$ est vrai dès le départ : affichage $\res{0\text{ puis }3}$. Le script calcule le quotient q et le reste r de la division euclidienne de a par b : $23=5\times4+3$.
\label{fig-cor-17-10}\figureGeo[corrige]{script-10}
\end{corrige}\end{blocLivre}
\begin{blocLivre}{cor-17-11}\begin{corrige}{11}{Une boucle qui dépasse}\demarche{Prouver que la variable progresse vers la condition d’arrêt.}
Les valeurs sont 1, 2, 4, 8, 16, 32, 64. Il faut $\res6$ tours avant $x>50$. Avec « ajouter 0 », x reste 1 ; le test ne devient jamais vrai et la boucle ne termine pas.
\label{fig-cor-17-11}\figureGeo[corrige]{script-11}
\end{corrige}\end{blocLivre}
\begin{blocLivre}{cor-17-12}\begin{corrige}{12}{Un bloc personnalisé}\demarche{Séparer le motif répété et le déplacement entre motifs.}
Le programme trace trois carrés de côté 30, en relevant le stylo pour avancer de 40 entre deux carrés : leurs points de départ sont espacés de 40 unités. La longueur tracée vaut $3\times4\times30=\res{360}$ unités. Le déplacement n’est pas dans le bloc carré, sinon il déplacerait aussi les coins.
\label{fig-cor-17-12}\figureGeo[corrige]{script-carres}
\end{corrige}\end{blocLivre}
\begin{blocLivre}{cor-17-13}\begin{corrige}{13}{Simuler un dé}\demarche{Initialiser une fois, compter à chaque tour, diviser à la fin.}
On met succès à $\res0$ avant la boucle ; on ajoute $\res1$ à succès quand le dé vaut 6 ; on affiche succès/$\res{600}$ après la boucle. La fréquence devrait approcher la probabilité $1/6\approx0{,}17$, sans être forcément égale.
\label{fig-cor-17-13}\figureGeo[corrige]{script-simulation}
\end{corrige}\end{blocLivre}
\begin{blocLivre}{cor-17-14}\begin{corrige}{14}{Deux dés et un « ou »}\demarche{Traduire la condition composée en évènement, puis utiliser la relation de la réunion.}
Il compte les lancers où le premier dé ou le second donne 6 : l’évènement « au moins un six ». Sa probabilité vaut $1/6+1/6-1/36=\res{11/36}\approx0{,}31$. L’affichage plausible est $\res{0{,}31}$ ; 0,17 correspondrait à un seul dé et 0,50 serait très éloigné pour 1000 lancers.
\label{fig-cor-17-14}\figureGeo[corrige]{script-deux-des}
\end{corrige}\end{blocLivre}
\begin{blocLivre}{cor-17-15}\begin{corrige}{15}{Des polygones réguliers}\demarche{Exécuter le bloc pour chaque valeur du paramètre.}
On obtient un $\res{\text{triangle équilatéral, un carré, un pentagone et un hexagone réguliers}}$. Pour l’hexagone, le lutin tourne de $360/6=\res{60^\circ}$ à chaque sommet. La longueur totale vaut $30\times(3+4+5+6)=\res{540}$ unités. Après chaque polygone, le lutin a tourné de 360° et revient à son point de départ.
\label{fig-cor-17-15}\figureGeo[corrige]{script-polygones}
\end{corrige}\end{blocLivre}

\clearpage\pagestyle{otMain}\fontsize{11}{15}\selectfont\corrigesChapitre{18}{Raisonner et préparer le brevet}
\begin{blocLivre}{cor-18-1}\begin{corrige}{1}{Un contre-exemple}\demarche{Chercher un cas qui contredit exactement l’affirmation.}
Avec $x=1/2$, $x^2=1/4<1/2$. Ce $\res{\text{seul contre-exemple suffit}}$. Avec 0 ou 1, l’inégalité stricte serait également fausse, par égalité.
\end{corrige}\end{blocLivre}
\begin{blocLivre}{cor-18-2}\begin{corrige}{2}{Deux entiers impairs}\demarche{Traduire l’hypothèse par une écriture générale.}
Écrivons les deux impairs $2a+1$ et $2b+1$, avec a, b entiers. Leur somme vaut $2(a+b+1)$, donc c’est $\res{\text{un multiple de 2}}$. La preuve vaut pour tous les entiers concernés.
\end{corrige}\end{blocLivre}
\begin{blocLivre}{cor-18-3}\begin{corrige}{3}{Réciproque et contraposée}\demarche{Échanger hypothèse et conclusion, puis les nier.}
Réciproque : « si un entier est pair, alors il est multiple de 6 » ; elle est $\res{\text{fausse}}$ : 4 est pair mais n’est pas multiple de 6. Contraposée : « si un entier n’est pas pair, alors il n’est pas multiple de 6 » ; elle est $\res{\text{vraie}}$, comme toute contraposée d’une propriété vraie.
\end{corrige}\end{blocLivre}
\begin{blocLivre}{cor-18-4}\begin{corrige}{4}{Un arrondi utile}\demarche{Interpréter la quantité entière dans le contexte.}
$7{,}2/3=2{,}4$. Deux bidons n’apportent que 6 L ; trois apportent 9 L. Il faut $\res3$ bidons. L’arrondi à l’entier le plus proche ne couvre pas forcément le besoin.
\end{corrige}\end{blocLivre}
\begin{blocLivre}{cor-18-5}\begin{corrige}{5}{Une hypothèse manquante}\demarche{Identifier l’hypothèse nécessaire.}
$\res{\text{Non}}$ : il faudrait savoir que les côtés de 3 et 4 cm sont perpendiculaires pour appliquer Pythagore ainsi. Un triangle de côtés 3, 4 et 6 cm existe aussi, car $3+4>6$, et son troisième côté n’est pas 5.
\end{corrige}\end{blocLivre}
\begin{blocLivre}{cor-18-6}\begin{corrige}{6}{Un programme constant}\demarche{Passer des essais à l’expression générale.}
Le résultat est $2(x+5)-2x=2x+10-2x=\res{10}$. Des essais suggèrent ce résultat ; la simplification le prouve pour tout x.
\end{corrige}\end{blocLivre}
\begin{blocLivre}{cor-18-7}\begin{corrige}{7}{Un raisonnement par l’absurde}\demarche{Supposer le contraire et chercher une contradiction.}
Supposons qu’un triangle rectangle ait aussi un angle obtus, donc supérieur à $90^\circ$. La somme de ces deux angles dépasserait déjà $180^\circ$, alors que la somme des trois angles d’un triangle vaut $180^\circ$ : contradiction. La supposition est fausse : $\res{\text{aucun triangle rectangle n’a d’angle obtus}}$.
\end{corrige}\end{blocLivre}
\begin{blocLivre}{cor-18-8}\begin{corrige}{8}{Une conjecture de divisibilité}\demarche{Distinguer les deux parités possibles.}
Deux entiers consécutifs comprennent toujours un pair : si n est pair, le produit contient ce facteur ; sinon $n+1$ est pair. Dans les deux cas, $\res{n(n+1)\text{ est multiple de 2}}$. On a distingué tous les cas possibles.
\end{corrige}\end{blocLivre}
\begin{blocLivre}{cor-18-9}\begin{corrige}{9}{Vérifier une équation}\demarche{Une substitution rapide détecte une erreur.}
Pour 4, les membres valent 6 et 8 : ce n’est pas une solution. $3x-6=x+4$ équivaut à $2x=10$, donc $\res{x=5}$. Le contrôle donne 9 de chaque côté.
\end{corrige}\end{blocLivre}
\begin{blocLivre}{cor-18-10}\begin{corrige}{10}{Une moyenne mal interprétée}\demarche{Construire un contre-exemple satisfaisant la donnée.}
$\res{\text{Non}}$. Considérons quatre notes $10\,;10\,;10\,;18$ : leur moyenne vaut 12, mais une seule est au moins 12. La propriété de partage en deux moitiés concerne la médiane, pas la moyenne.
\end{corrige}\end{blocLivre}
\begin{blocLivre}{cor-18-11}\begin{corrige}{11}{Une commande au meilleur prix}\demarche{Inclure conditionnement et frais fixes.}
A coûte $24\times2{,}80+6=\res{73{,}20\text{ €}}$. Il faut trois lots chez B : $\res{63\text{ €}}$. B économise $\res{10{,}20\text{ €}}$. On compare les coûts totaux pour la même quantité.
\end{corrige}\end{blocLivre}
\begin{blocLivre}{cor-18-12}\begin{corrige}{12}{Un modèle de bassin}\demarche{Relier volume, aire de base et hauteur.}
6000 L valent 6 m$^3$. La base vaut 10 m$^2$, donc la hauteur est $6/10=\res{0{,}6\text{ m}}$. Le volume maximal vaut 8 m$^3$ ; il reste $\res{2000\text{ L}}$.
\label{fig-cor-18-12}\figureGeo[corrige]{synthese-bassin}
\end{corrige}\end{blocLivre}
\begin{blocLivre}{cor-18-13}\begin{corrige}{13}{Un hasard et un budget}\demarche{Séparer modèle aléatoire et bilan comptable.}
La probabilité de gain vaut $20/200=\res{1/10}$. Les recettes sont 600 € et les lots coûtent 300 €. Le solde vaut $\res{300\text{ €}}$. Le résultat financier global ne décrit pas la chance d’un billet particulier.
\end{corrige}\end{blocLivre}
\begin{blocLivre}{cor-18-14}\begin{corrige}{14}{Un maximum expliqué}\demarche{Transformer l’expression avec une identité remarquable.}
L’autre côté vaut $20-x$, avec $0<x<20$. L’aire vaut $20x-x^2$, et $100-(x-10)^2=100-x^2+20x-100=20x-x^2$. Un carré étant positif ou nul, l’aire est au plus $\res{100\text{ cm}^2}$, valeur atteinte pour $x=10$, donc pour un carré.
\label{fig-cor-18-14}\figureGeo[corrige]{rectangle-fonction}
\end{corrige}\end{blocLivre}
\begin{blocLivre}{cor-18-15}\begin{corrige}{15}{Un plan de révision}\demarche{Relier le bilan à une action précise.}
Les automatismes rapportent $8\times0{,}5=4$ points. Le total vaut $\res{13/20}$. Reprendre le cours de fractions, refaire les questions ratées sans regarder la solution, puis contrôler sur des exercices nouveaux cible la difficulté identifiée. Ce bilan guide la suite sans résumer toutes les compétences à une note.
\end{corrige}\end{blocLivre}

\editionPage[\pdfbookmark[0]{Corrigés des défis}{bm-defi-cor-1-parent}]{defi-cor-1}{D01 · Corrigé}{Une fraction entre deux autres}{\hyperref[defi-ex-1]{Énoncé p.\,\pageref*{defi-ex-1}}}
\begin{blocEdition}{defi-cor-1}
$5/12=25/60$ et $7/15=28/60$. Le numérateur doit être strictement entre 25 et 28 : on obtient $\res{26/60\text{ et }27/60}$, soit exactement deux fractions demandées. Au dénominateur 120, les bornes sont 50/120 et 56/120 ; par exemple $\res{53/120}$ convient. Il existe cinq choix de numérateur, de 51 à 55. Attention : compter des écritures de dénominateur imposé ne signifie pas qu’elles soient irréductibles.
\end{blocEdition}
\editionPage{defi-cor-2}{D02 · Corrigé}{Comparer sans grands calculs}{\hyperref[defi-ex-2]{Énoncé p.\,\pageref*{defi-ex-2}}}
\begin{blocEdition}{defi-cor-2}
$2^{10}=1024>1000=10^3$. En élevant au cube des nombres positifs, $2^{30}=1024^3>1000^3=10^9$. En élevant cette dernière inégalité au carré, $\res{2^{60}>10^{18}}$. Une valeur approchée n’est pas nécessaire. La comparaison utilise seulement les règles des puissances et la croissance du carré et du cube sur les nombres positifs.
\end{blocEdition}
\editionPage{defi-cor-3}{D03 · Corrigé}{Une recherche entière exhaustive}{\hyperref[defi-ex-3]{Énoncé p.\,\pageref*{defi-ex-3}}}
\begin{blocEdition}{defi-cor-3}
$180=2^2\times3^2\times5$. Un diviseur est de la forme $2^a3^b5^c$, avec $0\le a,b\le2$ et $0\le c\le1$. Pour être un carré, les exposants doivent être pairs : $a,b$ valent 0 ou 2, et c vaut 0. Les quatre choix donnent $\res{1,4,9,36}$. Cette analyse épuise tous les exposants possibles, elle prouve donc la complétude de la liste.
\end{blocEdition}
\editionPage{defi-cor-4}{D04 · Corrigé}{L’identité de Sophie Germain}{\hyperref[defi-ex-4]{Énoncé p.\,\pageref*{defi-ex-4}}}
\begin{blocEdition}{defi-cor-4}
$(a^2+2b^2)^2-(2ab)^2=a^4+4a^2b^2+4b^4-4a^2b^2=a^4+4b^4$. Avec $a=n$ et $b=1$ : $n^4+4=(n^2+2)^2-(2n)^2=\res{(n^2-2n+2)(n^2+2n+2)}$. Or $n^2-2n+2=(n-1)^2+1\ge2$ pour $n\ge2$, et l’autre facteur est plus grand encore. Le nombre est le produit de deux entiers supérieurs ou égaux à 2 : $\res{\text{il n’est pas premier}}$. Pour $n=1$, on obtient 5, qui est premier. Cette identité porte le nom de Sophie Germain, mathématicienne française du XIX\textsuperscript{e} siècle.
\end{blocEdition}
\editionPage{defi-cor-5}{D05 · Corrigé}{Deux solutions et un contexte}{\hyperref[defi-ex-5]{Énoncé p.\,\pageref*{defi-ex-5}}}
\begin{blocEdition}{defi-cor-5}
La différence carré moins rectangle vaut $(x+1)^2-x(x+3)=1-x$. Une différence de 1 en grandeur signifie $1-x=1$ ou $1-x=-1$. On obtient x = 0 ou x = 2. La valeur 0 ne forme pas un rectangle de côté positif. Seule $\res{x=2}$ convient : rectangle d’aire 10 cm$^2$, carré d’aire 9 cm$^2$. Le contexte élimine une solution algébrique.
\end{blocEdition}
\editionPage{defi-cor-6}{D06 · Corrigé}{Promotions dans quel ordre ?}{\hyperref[defi-ex-6]{Énoncé p.\,\pageref*{defi-ex-6}}}
\begin{blocEdition}{defi-cor-6}
Réduction puis bon : $A=0{,}8P-10$. Bon puis réduction : $B=0{,}8(P-10)=0{,}8P-8$. Alors $B-A=\res{2\text{ €}}$. Le premier ordre est toujours plus avantageux de 2 €, pour les prix autorisés. Le bon de 10 € réduit avant pourcentage ne diminue finalement le prix que de 8 €. Les deux opérations ne commutent pas.
\end{blocEdition}
\editionPage{defi-cor-7}{D07 · Corrigé}{Une rencontre de débits}{\hyperref[defi-ex-7]{Énoncé p.\,\pageref*{defi-ex-7}}}
\begin{blocEdition}{defi-cor-7}
Pendant 12 minutes, le débit net vaut $18-3=15$ L/min. Le volume atteint est $15\times12=180$ L. La fuite vide ensuite 180 L à 3 L/min : il faut 60 minutes supplémentaires. Le temps total est $\res{72\text{ min}}$. On suppose la cuve assez grande et les débits constants ; aucun débordement n’est prévu.
\end{blocEdition}
\editionPage{defi-cor-8}{D08 · Corrigé}{Une fonction cachée derrière des essais}{\hyperref[defi-ex-8]{Énoncé p.\,\pageref*{defi-ex-8}}}
\begin{blocEdition}{defi-cor-8}
La fonction $g(x)=x^2+2+(x+1)x(x-1)$ possède bien les trois images indiquées, car le produit ajouté s’annule en -1, 0 et 1. Mais $f_0(2)=6$ pour la formule proposée et $g(2)=6+3\times2\times1=\res{12}$. Le tableau ne détermine donc pas une fonction unique. Une hypothèse supplémentaire sur sa famille serait nécessaire.
\end{blocEdition}
\editionPage{defi-cor-9}{D09 · Corrigé}{Trois tarifs, un choix argumenté}{\hyperref[defi-ex-9]{Énoncé p.\,\pageref*{defi-ex-9}}}
\begin{blocEdition}{defi-cor-9}
A = B pour n = 10 ; B = C pour $n=40/3\approx13{,}33$ ; A = C pour $n=35/3\approx11{,}67$. Pour 1 à 9 séances, A est la moins chère. Pour 10, A et B sont à égalité à 80 € et C est plus chère. Pour 11 à 13, B gagne. Pour 14 à 20, C gagne. Il faut comparer les trois coûts : l’intersection A-C ne signifie pas que ces deux offres soient les meilleures à cet endroit.
\end{blocEdition}
\editionPage{defi-cor-10}{D10 · Corrigé}{Une série à reconstruire}{\hyperref[defi-ex-10]{Énoncé p.\,\pageref*{defi-ex-10}}}
\begin{blocEdition}{defi-cor-10}
On impose $v_1=2$, $v_2=5$, $v_6=12$, $v_8=20$ et $v_4+v_5=16$. Pour une moyenne de 9, la somme vaut 72, donc $v_3+v_7=72-55=17$. La série $\res{2\,;5\,;5\,;7\,;9\,;12\,;12\,;20}$ convient : $Q_1=5$, médiane $(7+9)/2=8$, $Q_3=12$. Pour la seconde question, $v_3\le v_4\le8$, car $v_4\le v_5$ et $v_4+v_5=16$, et $v_7\le20$ : la somme vaut au plus $2+5+8+16+12+20+20=83$, donc la moyenne au plus $10{,}375$. Une moyenne de 15 est $\res{\text{impossible}}$.
\end{blocEdition}
\editionPage{defi-cor-11}{D11 · Corrigé}{Un jeu de somme équitable ?}{\hyperref[defi-ex-11]{Énoncé p.\,\pageref*{defi-ex-11}}}
\begin{blocEdition}{defi-cor-11}
Une somme 6 possède 5 couples, 8 en possède aussi 5, et 7 en possède 6. A gagne sur 10 couples, B sur 6. Parmi les 16 couples décisifs équiprobables, les parts sont $10/16=\res{5/8}$ et $6/16=\res{3/8}$. Le jeu avantage A. Recommencer les lancers indécis ne rend pas les trois sommes équiprobables.
\end{blocEdition}
\editionPage{defi-cor-12}{D12 · Corrigé}{Une corde dans un cercle}{\hyperref[defi-ex-12]{Énoncé p.\,\pageref*{defi-ex-12}}}
\begin{blocEdition}{defi-cor-12}
Le triangle OHA est rectangle en H. $HA^2=OA^2-OH^2=169-25=144$, donc HA = 12 cm. H étant le milieu de AB, $\res{AB=24\text{ cm}}$. L’aire vaut $24\times5/2=\res{60\text{ cm}^2}$. Le rayon n’est pas la hauteur relative à AB : c’est OH qui lui est perpendiculaire.\label{fig-defi-cor-12}\figureGeo[corrige]{defi-corde}

\end{blocEdition}
\editionPage{defi-cor-13}{D13 · Corrigé}{Les diagonales d’un trapèze}{\hyperref[defi-ex-13]{Énoncé p.\,\pageref*{defi-ex-13}}}
\begin{blocEdition}{defi-cor-13}
Les points A, O, C sont alignés, ainsi que B, O, D, et $(AB)\parallel(CD)$ : c’est un papillon de sommet O. Le théorème de Thalès donne $OA/OC=OB/OD=AB/CD=\res{6/10=3/5}$. Si AC = 12 cm, on partage 12 cm selon le ratio $3:5$ : $\res{OA=4{,}5\text{ cm}}$ et $\res{OC=7{,}5\text{ cm}}$. Contrôle : $4{,}5/7{,}5=0{,}6$.\label{fig-defi-cor-13}\figureGeo[corrige]{defi-trapeze}

\end{blocEdition}
\editionPage{defi-cor-14}{D14 · Corrigé}{Deux angles de visée}{\hyperref[defi-ex-14]{Énoncé p.\,\pageref*{defi-ex-14}}}
\begin{blocEdition}{defi-cor-14}
Notons x la distance entre B et le pied. L’angle de 45 degrés impose une hauteur h = x. Depuis A, $h/(x+10)=\tan30^\circ$. Donc $x(1-\tan30^\circ)=10\tan30^\circ$, puis $x=10\tan30^\circ/(1-\tan30^\circ)\approx13{,}6603$. La hauteur est $\res{13{,}7\text{ m}}$. Les distances et la hauteur obtenues sont positives.\label{fig-defi-cor-14}\figureGeo[corrige]{defi-visees}

\end{blocEdition}
\editionPage{defi-cor-15}{D15 · Corrigé}{Une somme de deux vecteurs}{\hyperref[defi-ex-15]{Énoncé p.\,\pageref*{defi-ex-15}}}
\begin{blocEdition}{defi-cor-15}
Par Chasles, $\vect{AE}=\vect{AB}+\vect{BE}=\vect{AB}+\vect{AC}=\vect{AD}$, donc $\res{E=D}$. Comme $\vect{BD}=\vect{AC}$, ABDC est $\res{\text{un parallélogramme}}$. Ses diagonales [AD] et [BC] ont le même milieu I, donc $\vect{AD}=2\vect{AI}$ : $\res{\vect{AB}+\vect{AC}=2\vect{AI}}$. Cette égalité est le point de départ de nombreuses démonstrations au lycée.\label{fig-defi-cor-15}\figureGeo[corrige]{defi-vecteurs}

\end{blocEdition}
\editionPage{defi-cor-16}{D16 · Corrigé}{Une section d’aire imposée}{\hyperref[defi-ex-16]{Énoncé p.\,\pageref*{defi-ex-16}}}
\begin{blocEdition}{defi-cor-16}
Le rayon du disque vaut 8 cm car $\pi r^2=64\pi$ et r est positif. Pythagore donne $OH^2=100-64=36$, donc $\res{OH=6\text{ cm}}$. Pour une direction de coupe fixée, deux plans parallèles conviennent, à 6 cm de part et d’autre du centre. Sans direction imposée, il existe de nombreuses orientations de plans.\label{fig-defi-cor-16}\figureGeo[corrige]{defi-section}

\end{blocEdition}
\editionPage{defi-cor-17}{D17 · Corrigé}{Un script de recherche et sa preuve}{\hyperref[defi-ex-17]{Énoncé p.\,\pageref*{defi-ex-17}}}
\begin{blocEdition}{defi-cor-17}
Après k tours, s est la somme de 1 à k et n vaut k+1. Les sommes passent par 91 pour k = 13 puis 105 pour k = 14. Le test devient faux après le quatorzième tour ; la sortie contient $\res{n=15}$ et s = 105. Une écriture par blocs utilise « répéter jusqu’à s $\ge$ 100 ». La somme augmente d’au moins 1 à chaque tour, ce qui garantit l’arrêt.
\end{blocEdition}
\editionPage{defi-cor-18}{D18 · Corrigé}{Le meilleur rectangle entier}{\hyperref[defi-ex-18]{Énoncé p.\,\pageref*{defi-ex-18}}}
\begin{blocEdition}{defi-cor-18}
L’aire vaut $x(23-x)$ avec x entier entre 1 et 22. On calcule $132-x(23-x)=(x-11)(x-12)$. Pour un entier x, les facteurs sont tous deux négatifs ou nuls si $x\le11$, tous deux positifs ou nuls si $x\ge12$. Leur produit est donc positif ou nul : l’aire ne dépasse pas 132. L’égalité est atteinte pour x = 11 ou 12, soit le même rectangle de $\res{11\text{ cm sur }12\text{ cm}}$. L’argument utilise le caractère entier : aucun entier n’est strictement entre 11 et 12.
\end{blocEdition}

\editionPage[\pdfbookmark[0]{Corrigés des brevets blancs}{bm-bilan-A-cor-auto-parent}]{bilan-A-cor-auto}{Brevet A · Corrigé}{Automatismes expliqués}{\hyperref[bilan-A-auto]{Énoncé p.\,\pageref*{bilan-A-auto}}}
\begin{blocEdition}{bilan-A-cor-auto}
\begin{questions}
\item $\res8$ : la multiplication est prioritaire.
\item $4/6+1/6=\res{5/6}$.
\item $\res{4{,}5\times10^5}$.
\item $\res{x=-9\text{ ou }x=9}$ : deux solutions opposées.
\item $\res{x^2+6x+9}$.
\item $\res{x\le-4}$ : diviser par $-2$ change le sens.
\item $0{,}15\times80=\res{12}$.
\item $-4+1=\res{-3}$.
\item 7 parts de 5 € : $\res{10\text{ € et }25\text{ €}}$.
\item La troisième valeur est $\res5$.
\item $0{,}5+0{,}4-0{,}2=\res{0{,}7}$.
\item $\res{\vect{AB}=\vect{DC}}$ : côtés opposés de même sens.\end{questions}
\end{blocEdition}
\editionPage{bilan-A-cor-1}{Brevet A · Corrigé 1}{Un atelier de carnets}{\hyperref[bilan-A-ex-1]{Énoncé p.\,\pageref*{bilan-A-ex-1}}}
\begin{blocEdition}{bilan-A-cor-1}
\editionEtape{1 · Partager sans reste}Le nombre de carnets divise 168 et 126. $168=2^3\times3\times7$ et $126=2\times3^2\times7$. Le plus grand diviseur commun est $2\times3\times7=42$. On prépare $\res{42\text{ carnets}}$, chacun contenant $\res4$ feuilles bleues et $\res3$ blanches.\editionEtape{2 · Recettes et solde}Les recettes sont $42\times4{,}50=189$ €. Le solde, hors autres frais, est $189-78=\res{111\text{ €}}$.\editionEtape{3 · Une comparaison de prix}Diminuer de 10 \% revient à multiplier par 0,9 : $5\times0{,}9=\res{4{,}50\text{ €}}$. Les prix unitaires sont égaux.\editionRetenir{Le maximum de lots identiques correspond au plus grand diviseur commun, pas au plus petit multiple commun.}
\end{blocEdition}
\editionPage{bilan-A-cor-2}{Brevet A · Corrigé 2}{Une voile divisée}{\hyperref[bilan-A-ex-2]{Énoncé p.\,\pageref*{bilan-A-ex-2}}}
\begin{blocEdition}{bilan-A-cor-2}
\editionEtape{1 · Le triangle rectangle}Le théorème de Pythagore donne $BC^2=12^2+9^2=225$, donc $\res{BC=15\text{ m}}$.\editionEtape{2 · La réduction}Les alignements et le parallélisme permettent d’appliquer le théorème de Thalès : $AM/AB=AN/AC=MN/BC=4/12=1/3$. Ainsi $\res{AN=3\text{ m}}$ et $\res{MN=5\text{ m}}$.\editionEtape{3 · Une translation}Par définition de la translation de vecteur $\vect{MB}$, $\vect{ND}=\vect{MB}$ : $\res{MBDN\text{ est un parallélogramme}}$. Donc $(BD)\parallel(MN)$. Comme $(MN)\parallel(BC)$, les droites (BD) et (BC) sont parallèles et ont le point B en commun : elles sont confondues, et $\res{D\in(BC)}$.\editionRetenir{Deux droites parallèles ayant un point commun sont confondues : c’est un argument classique pour prouver un alignement.}\label{fig-bilan-A-cor-2}\figureGeo[corrige]{brevet-A-geo}

\end{blocEdition}
\editionPage{bilan-A-cor-3}{Brevet A · Corrigé 3}{Deux locations de matériel}{\hyperref[bilan-A-ex-3]{Énoncé p.\,\pageref*{bilan-A-ex-3}}}
\begin{blocEdition}{bilan-A-cor-3}
\editionEtape{1 · Quatre journées}A coûte $18+28=\res{46\text{ €}}$ ; B coûte $\res{40\text{ €}}$. B est moins chère de 6 €.\editionEtape{2 · Le seuil}$18+7n=10n$ donne $3n=18$, donc $\res{n=6}$ ; les deux prix valent 60 €. L’inéquation $18+7n<10n$ équivaut à $18<3n$, soit $n>6$ : $\res{A\text{ est moins chère à partir de 7 jours}}$, B l’est pour moins de 6 jours.\editionEtape{3 · Un budget}Chez A, $18+7n\le90$ équivaut à $n\le72/7\approx10{,}3$ : maximum $\res{10\text{ jours}}$. Chez B, $10n\le90$ équivaut à $n\le9$ : maximum $\res{9\text{ jours}}$. On ne garde que des nombres entiers de jours.\label{fig-bilan-A-cor-3}\figureGeo[corrige]{brevet-A-tarifs}

\end{blocEdition}
\editionPage{bilan-A-cor-4}{Brevet A · Corrigé 4}{Des scores et un tirage}{\hyperref[bilan-A-ex-4]{Énoncé p.\,\pageref*{bilan-A-ex-4}}}
\begin{blocEdition}{bilan-A-cor-4}
\editionEtape{1 · La série}La somme pondérée vaut $0+5+16+15=36$, donc la moyenne vaut $\res{1{,}8}$. Les effectifs cumulés sont 2, 7, 15, 20. Les rangs 10 et 11 valent 2 : médiane $\res2$. $Q_1$ est la 5\textsuperscript{e} valeur : $\res{Q_1=1}$ ; $Q_3$ est la 15\textsuperscript{e} : $\res{Q_3=2}$.\editionEtape{2 · Deux rouges}Le tableau compte $4\times4=16$ couples équiprobables, dont $3\times3=9$ formés de deux rouges : $\res{9/16}$.\editionEtape{3 · Une bleue, au moins une rouge}Exactement une bleue : $3+3=6$ couples, soit $\res{6/16=3/8}$. Au moins une rouge est le contraire de « deux bleues », qui ne contient que le couple $(B\,;B)$ : $\res{1-1/16=15/16}$.\label{fig-bilan-A-cor-4}\figureGeo[corrige]{brevet-A-tableau}

\end{blocEdition}
\editionPage{bilan-B-cor-auto}{Brevet B · Corrigé}{Automatismes expliqués}{\hyperref[bilan-B-auto]{Énoncé p.\,\pageref*{bilan-B-auto}}}
\begin{blocEdition}{bilan-B-cor-auto}
\begin{questions}
\item $\res6$.
\item $3/5\times10/9=\res{2/3}$.
\item $\res{10^{-2}}$.
\item $\res{2^2\times3\times5}$.
\item $\res{(x-7)(x+7)}$.
\item $\res{x=3\text{ ou }x=-1/2}$.
\item $50\times1{,}12=\res{56}$.
\item $-3x=12$, donc $\res{-4}$.
\item $\frac43\pi\times27=\res{36\pi\text{ cm}^3}$.
\item $8\times0{,}25=2$ : la 2\textsuperscript{e} valeur, $\res3$.
\item $1-0{,}3=\res{0{,}7}$.
\item $\res{\vect{AC}}$, d’après la relation de Chasles.\end{questions}
\end{blocEdition}
\editionPage{bilan-B-cor-1}{Brevet B · Corrigé 1}{Un abri et sa pente}{\hyperref[bilan-B-ex-1]{Énoncé p.\,\pageref*{bilan-B-ex-1}}}
\begin{blocEdition}{bilan-B-cor-1}
\editionEtape{1 · Longueur du versant}$BC^2=4^2+2{,}5^2=22{,}25$, donc $BC=\sqrt{22{,}25}\approx\res{4{,}72\text{ m}}$.\editionEtape{2 · Angle de pente}$\tan\widehat B=2{,}5/4=0{,}625$. La fonction inverse donne $\widehat B\approx32{,}005^\circ$, soit $\res{32^\circ}$.\editionEtape{3 · Achat entier}L’aire utile vaut $6\sqrt{22{,}25}\approx28{,}3019$ m$^2$. Avec les découpes, on multiplie par 1,08 : environ 30,5660 m$^2$. Le nombre de plaques doit être au moins $30{,}5660/2\approx15{,}283$. Il faut $\res{16\text{ plaques}}$. Quinze ne couvriraient que 30 m$^2$.\label{fig-bilan-B-cor-1}\figureGeo[corrige]{brevet-B-geo}

\end{blocEdition}
\editionPage{bilan-B-cor-2}{Brevet B · Corrigé 2}{Une bouteille et une boule}{\hyperref[bilan-B-ex-2]{Énoncé p.\,\pageref*{bilan-B-ex-2}}}
\begin{blocEdition}{bilan-B-cor-2}
\editionEtape{1 · Capacité}Le volume vaut $\pi\times4^2\times20=\res{320\pi\text{ cm}^3}$. Un cm$^3$ vaut un mL : capacité $\res{1005\text{ mL}}$ au mL.\editionEtape{2 · Remplissage}Le volume visé vaut $0{,}9\times320\pi=288\pi$ mL. La durée vaut $288\pi/50\approx\res{18{,}1\text{ s}}$.\editionEtape{3 · Une boule}Le récipient a un volume de $\frac43\pi\times7^3=\frac{1372\pi}3\approx1436{,}8$ cm$^3$, supérieur à $1005$ cm$^3$ : $\res{\text{il est assez grand}}$. La part remplie vaut $\dfrac{320\pi}{1372\pi/3}=\dfrac{240}{343}\approx\res{70\,\%}$.\label{fig-bilan-B-cor-2}\figureGeo[corrige]{brevet-B-cylindre}

\end{blocEdition}
\editionPage{bilan-B-cor-3}{Brevet B · Corrigé 3}{Un script de figures}{\hyperref[bilan-B-ex-3]{Énoncé p.\,\pageref*{bilan-B-ex-3}}}
\begin{blocEdition}{bilan-B-cor-3}
\editionEtape{1 · Une variable à chaque tour}Les carrés ont pour côtés $\res{10,15,20,25}$ unités : l’ajout de 5 vient après le tracé. Après le quatrième tour, c vaut 30 et la condition $c>25$ devient vraie : $\res{4\text{ passages}}$.\editionEtape{2 · Longueur dessinée}La somme des côtés vaut 70. Les périmètres totalisent $4\times70=\res{280}$ unités.\editionEtape{3 · État final}Le script affiche $\res{30}$. Les déplacements sans tracer valent $c+5$ à chaque tour, soit 15, 20, 25, 30 : leur somme vaut $\res{90}$ unités. La boucle s’arrête dès que $c>25$ ; aucun appel carré(30) n’est exécuté.\label{fig-bilan-B-cor-3}\figureGeo[corrige]{brevet-B-script}

\end{blocEdition}
\editionPage{bilan-B-cor-4}{Brevet B · Corrigé 4}{Un jeu à deux étapes}{\hyperref[bilan-B-ex-4]{Énoncé p.\,\pageref*{bilan-B-ex-4}}}
\begin{blocEdition}{bilan-B-cor-4}
\editionEtape{1 · Dénombrer}Le tableau donne $4\times2=\res{8\text{ issues équiprobables}}$. Les issues gagnantes sont A1-rouge, A2-rouge, A3-rouge et B-bleue : $P(G)=4/8=\res{1/2}$.\editionEtape{2 · Fréquence et probabilité}$\res{\text{Non}}$ : sur une série de 100 parties, le nombre de gains varie d’une série à l’autre. La probabilité $1/2$ indique que la fréquence des gains se stabilise autour de 0,5 quand le nombre de parties devient très grand.\editionEtape{3 · Réunion}$P(E)=6/8=3/4$, $P(F)=4/8=1/2$ et $E\cap F$ contient 3 issues : $P(E\cap F)=3/8$. La relation du cours donne $P(E\cup F)=3/4+1/2-3/8=\res{7/8}$. Contrôle : seule l’issue B-bleue n’est ni dans E ni dans F.\label{fig-bilan-B-cor-4}\figureGeo[corrige]{brevet-B-tableau}

\end{blocEdition}

\editionPage[\pdfbookmark[0]{Fiches outils}{bm-outils-debut-parent}]{outils-debut}{Fiche outil}{Mesures et formules utiles}{À consulter au fil des chapitres.}
\begin{blocEdition}{outils-debut}

\editionEtape{Contours et aires}
\begin{tabularx}{\linewidth}{@{}lYY@{}}\toprule
Figure & Périmètre & Aire\\\midrule
Rectangle & $2(L+\ell)$ & $L\ell$\\[2mm]
Carré & $4c$ & $c^2$\\[2mm]
Triangle & somme des côtés & base $\times$ hauteur associée / 2\\[2mm]
Parallélogramme & somme des côtés & base $\times$ hauteur associée\\[2mm]
Disque / cercle & $2\pi r$ pour le cercle & $\pi r^2$ pour le disque\\\bottomrule
\end{tabularx}
\editionEtape{Volumes}
Pavé : $L\ell h$ ; prisme droit et cylindre : aire de la base $\times$ hauteur ;
pyramide et cône : aire de la base $\times$ hauteur / 3 ; boule de rayon $r$ : $\frac43\pi r^3$.
\editionEtape{Unités}
$1$ km = $1000$ m ; $1$ m = $100$ cm ; $1$ m$^2$ = $10000$ cm$^2$ ;
$1$ m$^3$ = $1000$ L ; $1$ L = $1000$ mL ; $1$ mL = $1$ cm$^3$.
Une vitesse en km/h se divise par 3,6 pour obtenir des m/s.
\editionRetenir{La hauteur d’un triangle est perpendiculaire à la base choisie. Un côté oblique n’est pas automatiquement une hauteur.}

\end{blocEdition}
\editionPage{outils-2}{Fiche outil}{Données : quartiles et boîtes à moustaches}{À consulter au fil des chapitres.}
\begin{blocEdition}{outils-2}

\editionEtape{Un exemple complet}
La série de pointures $36\,;37\,;37\,;38\,;38\,;38\,;39\,;40\,;40\,;42\,;43\,;44$ compte
$N=12$ valeurs. La moyenne vaut $472/12\approx39{,}3$. La médiane est la moyenne des
6\textsuperscript{e} et 7\textsuperscript{e} valeurs : $(38+39)/2=38{,}5$.
\editionEtape{Quartiles}
$12\times0{,}25=3$ : $Q_1$ est la 3\textsuperscript{e} valeur, soit 37.
$12\times0{,}75=9$ : $Q_3$ est la 9\textsuperscript{e} valeur, soit 40.
Au moins 25 \% des pointures sont inférieures ou égales à 37, au moins 75 \%
inférieures ou égales à 40.
\editionEtape{Boîte à moustaches}
\begin{center}\begin{tikzpicture}[x=.62cm,font=\small]
\draw[otMuted,->](35,0)--(45.5,0);
\foreach\t in{35,...,45}{\draw[otMuted](\t,.07)--(\t,-.07)node[below,font=\tiny]{\t};}
\draw[geoTrait](36,.7)--(37,.7) (40,.7)--(44,.7) (36,.55)--(36,.85) (44,.55)--(44,.85);
\draw[geoTrait](37,.4)rectangle(40,1);\draw[geoTrait,line width=1.4pt](38.5,.4)--(38.5,1);
\end{tikzpicture}\end{center}
Minimum 36, $Q_1=37$, médiane 38,5, $Q_3=40$, maximum 44 ; l’étendue vaut 8.
\editionRetenir{La boîte contient environ la moitié centrale des valeurs ; sa longueur $Q_3-Q_1$ mesure leur dispersion.}

\end{blocEdition}
\editionPage{outils-3}{Fiche outil}{Décomposer en facteurs premiers avec un algorithme}{À consulter au fil des chapitres.}
\begin{blocEdition}{outils-3}

\editionEtape{Une procédure en français structuré}
Choisir un entier $N\ge2$. Mettre reste à N et diviseur à 2.
Répéter jusqu’à ce que reste soit égal à 1 :\par
\begin{itemize}
\item si le reste de la division de reste par diviseur vaut 0, afficher
diviseur et remplacer reste par reste / diviseur ;
\item sinon, augmenter diviseur de 1.
\end{itemize}
La boucle conditionnelle teste donc de nouveau le même diviseur après une division réussie.
En blocs, employer « répéter jusqu’à reste = 1 » et une condition à deux branches.
\editionEtape{Exemple pour 84}
On affiche 2, puis 2, puis 3, puis 7. Les quotients successifs après
division sont 42, 21, 7 et 1. La décomposition vaut
$84=2^2\times3\times7$. Un diviseur composé ne sera pas affiché :
ses facteurs premiers auraient été retirés auparavant.
\editionEtape{Contrôles}
Multiplier les facteurs affichés pour retrouver N.
Tester un nombre premier, par exemple 13 : un seul facteur est affiché.
La procédure doit refuser 0, 1 et les entrées non entières dans ce cadre.
Elle termine, car soit le reste diminue par une division, soit le diviseur
augmente jusqu’à atteindre un diviseur du reste.

\end{blocEdition}
\editionPage{outils-4}{Fiche outil}{Rédiger un raisonnement géométrique}{À consulter au fil des chapitres.}
\begin{blocEdition}{outils-4}

\editionEtape{Une longueur avec Pythagore}
« Le triangle ABC est rectangle en A. D’après le théorème de Pythagore,
$BC^2=AB^2+AC^2$. Donc… Comme BC est une longueur positive, BC = … ».
Pour prouver qu’un triangle n’est pas rectangle, citer la contraposée.
\editionEtape{Un parallélisme avec Thalès}
« A, M, B sont alignés dans cet ordre, ainsi que A, N, C.
On compare $AM/AB$ et $AN/AC$ : … Les deux rapports sont égaux.
D’après la réciproque du théorème de Thalès, $(MN)\parallel(BC)$. »
Si les rapports diffèrent, conclure avec la contraposée.
\editionEtape{Un angle avec la trigonométrie}
« Dans ABC rectangle en A, pour l’angle $\widehat B$, le côté AC est opposé et
AB adjacent. Donc $\tan\widehat B=AC/AB$. La fonction inverse donne
$\widehat B\approx\ldots^\circ$. » Nommer l’angle évite un rapport inversé.
\editionEtape{Une égalité de vecteurs}
« ABCD est un parallélogramme, donc $\vect{AB}=\vect{DC}$. » Pour démontrer
qu’un quadrilatère est un parallélogramme, établir une égalité de vecteurs
en respectant l’ordre des lettres.
\editionRetenir{Une rédaction courte suffit si elle contient les hypothèses utiles, la propriété, le calcul et la conclusion.}

\end{blocEdition}

\clearpage\phantomsection\label{couverture-programme}\pdfbookmark[0]{Correspondance avec le programme}{bm-couverture-programme}
\gdef\otChapNum{0}\pagestyle{otFront}\otPageTitre{Programme et chapitres}{Troisième · Programme 2026 applicable à la rentrée 2028}
Chaque rubrique de Troisième de l’annexe 2 de l’arrêté du 18 février 2026 est
reliée aux chapitres qui la traitent. Les pages suivantes associent ensuite quatre
repères par chapitre à des exercices précis ; les séries citées ne sont pas exhaustives.\par\medskip
{\fontsize{9.5}{12.5}\selectfont
\begin{tabularx}{\linewidth}{@{}>{\raggedright\arraybackslash}p{32mm}Yr@{}}\toprule
Domaine & Rubrique de Troisième & Chapitres\\\midrule
Nombres et calculs & Nombres rationnels : forme irréductible, problèmes ; notation des ensembles de nombres (prolongement) & 1\\[1.6mm]
 & Puissances : exposants négatifs, produits et quotients, notation scientifique & 2, 7\\[1.6mm]
 & Racine carrée : équations $x^2=a$, résolution par le calcul et graphique, problèmes & 2, 5, 8\\[1.6mm]
 & Multiples et diviseurs : décomposition, critères de divisibilité, fractions (automatismes) & 3, 1\\[1.6mm]
 & Calcul littéral : produits et rapports, double distributivité, identités remarquables, produit nul, inéquation $ax\ge b$, analyse-synthèse & 4, 5\\[1.6mm]
Espace et géométrie & Repérage sur une droite graduée et dans le plan (automatismes) & 5, 6, 8, 9\\[1.6mm]
 & Représentation de l’espace : solides et volumes, boule, sphère, grands cercles, sections planes & 16\\[1.6mm]
 & Triangles : Thalès, réciproque et contraposée ; cosinus, sinus, tangente ; Pythagore, cercle circonscrit et droite des milieux (automatismes) & 12, 13, 14\\[1.6mm]
 & Translations et vecteurs : vecteurs égaux, nul, opposé, somme, relation de Chasles & 15\\[1.6mm]
Données et probabilités & Statistiques : effectifs cumulés, médiane, quartiles, boîtes à moustaches, tableur & 10\\[1.6mm]
 & Probabilités : relation $P(A\cup B)+P(A\cap B)=P(A)+P(B)$, expériences indépendantes, simulation et stabilisation des fréquences & 11, 17\\[1.6mm]
Proportionnalité, fonctions & Proportionnalité : ratios, pourcentages, échelles, coefficient multiplicateur, lien avec Thalès, fonctions linéaires & 6, 7, 9\\[1.6mm]
 & Fonctions : représentations, image et antécédents, fonctions linéaires et affines, inéquations, fonction carré & 8, 9\\[1.6mm]
Pensée informatique & Variables, conditions composées, boucle conditionnelle, blocs personnalisés, écriture de programmes & 17\\[1.6mm]
Principes du programme & Raisonnement : conjecture, preuve, contre-exemple, réciproque, contraposée, absurde ; automatismes et résolution de problèmes & 18 et tous\\[1.6mm]
\bottomrule\end{tabularx}}\par\medskip
\begin{attention}Les automatismes de Troisième s’appuient sur des notions étudiées en Quatrième sans y être automatisées. L’épreuve évalue les connaissances et compétences attendues en fin de cycle 4, déclinées par le programme de Troisième. Les défis prolongent ces notions sans ajouter d’exigence à l’examen.\end{attention}
\clearpage
\otPageTitre{Repères et exercices}{Chapitres 1 à 3}\fontsize{10}{13}\selectfont
\otIntertitre{1. Nombres rationnels et fractions}\begin{tabularx}{\linewidth}{@{}Yl@{}}\toprule
Priorités, signes et vocabulaire & \hyperref[exo-1-1]{1.1}, \hyperref[exo-1-2]{1.2}, \hyperref[exo-1-10]{1.10}\\[2mm]
Opérations et fractions irréductibles & \hyperref[exo-1-3]{1.3}, \hyperref[exo-1-4]{1.4}, \hyperref[exo-1-5]{1.5}, \hyperref[exo-1-6]{1.6}, \hyperref[exo-1-8]{1.8}\\[2mm]
Comparaisons et ensembles de nombres & \hyperref[exo-1-7]{1.7}, \hyperref[exo-1-9]{1.9}, \hyperref[exo-1-15]{1.15}\\[2mm]
Fractions dans les problèmes & \hyperref[exo-1-11]{1.11}, \hyperref[exo-1-12]{1.12}, \hyperref[exo-1-13]{1.13}, \hyperref[exo-1-14]{1.14}\\[2mm]
\bottomrule\end{tabularx}\par\vspace{4mm}
\otIntertitre{2. Puissances, racines et notation scientifique}\begin{tabularx}{\linewidth}{@{}Yl@{}}\toprule
Puissances d’exposant entier & \hyperref[exo-2-1]{2.1}, \hyperref[exo-2-2]{2.2}, \hyperref[exo-2-3]{2.3}, \hyperref[exo-2-11]{2.11}\\[2mm]
Notation scientifique et préfixes & \hyperref[exo-2-4]{2.4}, \hyperref[exo-2-6]{2.6}, \hyperref[exo-2-7]{2.7}, \hyperref[exo-2-8]{2.8}, \hyperref[exo-2-12]{2.12}, \hyperref[exo-2-13]{2.13}\\[2mm]
Racine carrée et équation $x^2=a$ & \hyperref[exo-2-5]{2.5}, \hyperref[exo-2-9]{2.9}, \hyperref[exo-2-10]{2.10}, \hyperref[exo-2-15]{2.15}\\[2mm]
Problèmes de grandeurs et de racines & \hyperref[exo-2-11]{2.11}, \hyperref[exo-2-12]{2.12}, \hyperref[exo-2-13]{2.13}, \hyperref[exo-2-14]{2.14}, \hyperref[exo-2-15]{2.15}\\[2mm]
\bottomrule\end{tabularx}\par\vspace{4mm}
\otIntertitre{3. Multiples, diviseurs et nombres premiers}\begin{tabularx}{\linewidth}{@{}Yl@{}}\toprule
Divisibilité et nombres premiers & \hyperref[exo-3-1]{3.1}, \hyperref[exo-3-2]{3.2}, \hyperref[exo-3-5]{3.5}, \hyperref[exo-3-6]{3.6}, \hyperref[exo-3-7]{3.7}, \hyperref[exo-3-10]{3.10}\\[2mm]
Décomposition et fractions irréductibles & \hyperref[exo-3-3]{3.3}, \hyperref[exo-3-4]{3.4}, \hyperref[exo-3-8]{3.8}, \hyperref[exo-3-9]{3.9}\\[2mm]
Diviseurs communs et partages & \hyperref[exo-3-11]{3.11}, \hyperref[exo-3-13]{3.13}\\[2mm]
Multiples communs et périodicité & \hyperref[exo-3-12]{3.12}, \hyperref[exo-3-14]{3.14}, \hyperref[exo-3-15]{3.15}\\[2mm]
\bottomrule\end{tabularx}\par\vspace{4mm}
\clearpage
\otPageTitre{Repères et exercices}{Chapitres 4 à 6}\fontsize{10}{13}\selectfont
\otIntertitre{4. Calcul littéral et identités remarquables}\begin{tabularx}{\linewidth}{@{}Yl@{}}\toprule
Nature, opposé et substitution & \hyperref[exo-4-1]{4.1}, \hyperref[exo-4-2]{4.2}, \hyperref[exo-4-3]{4.3}\\[2mm]
Simplifier et développer & \hyperref[exo-4-4]{4.4}, \hyperref[exo-4-6]{4.6}, \hyperref[exo-4-8]{4.8}, \hyperref[exo-4-12]{4.12}\\[2mm]
Identités remarquables et factorisation & \hyperref[exo-4-5]{4.5}, \hyperref[exo-4-7]{4.7}, \hyperref[exo-4-9]{4.9}, \hyperref[exo-4-10]{4.10}\\[2mm]
Preuves par le calcul littéral & \hyperref[exo-4-11]{4.11}, \hyperref[exo-4-13]{4.13}, \hyperref[exo-4-14]{4.14}, \hyperref[exo-4-15]{4.15}\\[2mm]
\bottomrule\end{tabularx}\par\vspace{4mm}
\otIntertitre{5. Équations, inéquations et produits nuls}\begin{tabularx}{\linewidth}{@{}Yl@{}}\toprule
Équations du premier degré & \hyperref[exo-5-1]{5.1}, \hyperref[exo-5-2]{5.2}, \hyperref[exo-5-3]{5.3}, \hyperref[exo-5-4]{5.4}, \hyperref[exo-5-6]{5.6}, \hyperref[exo-5-7]{5.7}\\[2mm]
Inéquations et représentation des solutions & \hyperref[exo-5-5]{5.5}, \hyperref[exo-5-10]{5.10}, \hyperref[exo-5-13]{5.13}\\[2mm]
Produits nuls et factorisations & \hyperref[exo-5-8]{5.8}, \hyperref[exo-5-9]{5.9}, \hyperref[exo-5-15]{5.15}\\[2mm]
Mise en équation, analyse et synthèse & \hyperref[exo-5-11]{5.11}, \hyperref[exo-5-12]{5.12}, \hyperref[exo-5-13]{5.13}, \hyperref[exo-5-14]{5.14}, \hyperref[exo-5-15]{5.15}\\[2mm]
\bottomrule\end{tabularx}\par\vspace{4mm}
\otIntertitre{6. Proportionnalité, ratios et pourcentages}\begin{tabularx}{\linewidth}{@{}Yl@{}}\toprule
Proportionnalité, coefficient et alignement & \hyperref[exo-6-1]{6.1}, \hyperref[exo-6-2]{6.2}, \hyperref[exo-6-15]{6.15}\\[2mm]
Ratios, partages et échelles & \hyperref[exo-6-5]{6.5}, \hyperref[exo-6-8]{6.8}, \hyperref[exo-6-9]{6.9}, \hyperref[exo-6-10]{6.10}\\[2mm]
Évolutions et coefficient multiplicateur & \hyperref[exo-6-3]{6.3}, \hyperref[exo-6-4]{6.4}, \hyperref[exo-6-6]{6.6}, \hyperref[exo-6-7]{6.7}\\[2mm]
Comparaisons et pourcentages emboîtés & \hyperref[exo-6-11]{6.11}, \hyperref[exo-6-12]{6.12}, \hyperref[exo-6-13]{6.13}, \hyperref[exo-6-14]{6.14}\\[2mm]
\bottomrule\end{tabularx}\par\vspace{4mm}
\clearpage
\otPageTitre{Repères et exercices}{Chapitres 7 à 9}\fontsize{10}{13}\selectfont
\otIntertitre{7. Grandeurs quotients et conversions}\begin{tabularx}{\linewidth}{@{}Yl@{}}\toprule
Conversions de longueurs, aires et volumes & \hyperref[exo-7-1]{7.1}, \hyperref[exo-7-2]{7.2}, \hyperref[exo-7-13]{7.13}, \hyperref[exo-7-15]{7.15}\\[2mm]
Durées et vitesses moyennes & \hyperref[exo-7-3]{7.3}, \hyperref[exo-7-4]{7.4}, \hyperref[exo-7-6]{7.6}, \hyperref[exo-7-7]{7.7}, \hyperref[exo-7-11]{7.11}\\[2mm]
Débits et masses volumiques & \hyperref[exo-7-5]{7.5}, \hyperref[exo-7-8]{7.8}, \hyperref[exo-7-12]{7.12}, \hyperref[exo-7-14]{7.14}, \hyperref[exo-7-15]{7.15}\\[2mm]
Grandeurs quotients et contrôle des unités & \hyperref[exo-7-9]{7.9}, \hyperref[exo-7-10]{7.10}, \hyperref[exo-7-13]{7.13}, \hyperref[exo-7-14]{7.14}\\[2mm]
\bottomrule\end{tabularx}\par\vspace{4mm}
\otIntertitre{8. Fonctions, images et fonction carré}\begin{tabularx}{\linewidth}{@{}Yl@{}}\toprule
Vocabulaire, notations et calcul d’images & \hyperref[exo-8-1]{8.1}, \hyperref[exo-8-3]{8.3}, \hyperref[exo-8-5]{8.5}, \hyperref[exo-8-6]{8.6}\\[2mm]
Tableaux, courbes et antécédents & \hyperref[exo-8-2]{8.2}, \hyperref[exo-8-4]{8.4}, \hyperref[exo-8-7]{8.7}, \hyperref[exo-8-9]{8.9}, \hyperref[exo-8-14]{8.14}\\[2mm]
Fonction carré et équation $x^2=a$ & \hyperref[exo-8-8]{8.8}, \hyperref[exo-8-10]{8.10}, \hyperref[exo-8-12]{8.12}\\[2mm]
Modèles et domaines de validité & \hyperref[exo-8-11]{8.11}, \hyperref[exo-8-13]{8.13}, \hyperref[exo-8-14]{8.14}, \hyperref[exo-8-15]{8.15}\\[2mm]
\bottomrule\end{tabularx}\par\vspace{4mm}
\otIntertitre{9. Fonctions linéaires et affines}\begin{tabularx}{\linewidth}{@{}Yl@{}}\toprule
Fonctions linéaires et proportionnalité & \hyperref[exo-9-1]{9.1}, \hyperref[exo-9-2]{9.2}, \hyperref[exo-9-4]{9.4}, \hyperref[exo-9-6]{9.6}, \hyperref[exo-9-13]{9.13}\\[2mm]
Fonctions affines et tracés & \hyperref[exo-9-5]{9.5}, \hyperref[exo-9-10]{9.10}, \hyperref[exo-9-12]{9.12}, \hyperref[exo-9-14]{9.14}\\[2mm]
Coefficients lus et calculés & \hyperref[exo-9-3]{9.3}, \hyperref[exo-9-7]{9.7}, \hyperref[exo-9-8]{9.8}, \hyperref[exo-9-9]{9.9}\\[2mm]
Inéquations, comparaisons et seuils & \hyperref[exo-9-11]{9.11}, \hyperref[exo-9-13]{9.13}, \hyperref[exo-9-15]{9.15}\\[2mm]
\bottomrule\end{tabularx}\par\vspace{4mm}
\clearpage
\otPageTitre{Repères et exercices}{Chapitres 10 à 12}\fontsize{10}{13}\selectfont
\otIntertitre{10. Statistiques et boîtes à moustaches}\begin{tabularx}{\linewidth}{@{}Yl@{}}\toprule
Moyenne, médiane et étendue & \hyperref[exo-10-1]{10.1}, \hyperref[exo-10-2]{10.2}, \hyperref[exo-10-4]{10.4}, \hyperref[exo-10-10]{10.10}, \hyperref[exo-10-11]{10.11}\\[2mm]
Effectifs cumulés et quartiles & \hyperref[exo-10-3]{10.3}, \hyperref[exo-10-6]{10.6}, \hyperref[exo-10-7]{10.7}, \hyperref[exo-10-9]{10.9}\\[2mm]
Boîtes à moustaches & \hyperref[exo-10-5]{10.5}, \hyperref[exo-10-8]{10.8}, \hyperref[exo-10-12]{10.12}\\[2mm]
Comparer et interpréter des séries & \hyperref[exo-10-12]{10.12}, \hyperref[exo-10-13]{10.13}, \hyperref[exo-10-14]{10.14}, \hyperref[exo-10-15]{10.15}\\[2mm]
\bottomrule\end{tabularx}\par\vspace{4mm}
\otIntertitre{11. Probabilités, évènements et simulations}\begin{tabularx}{\linewidth}{@{}Yl@{}}\toprule
Évènements, contraire et notations & \hyperref[exo-11-1]{11.1}, \hyperref[exo-11-2]{11.2}, \hyperref[exo-11-3]{11.3}, \hyperref[exo-11-4]{11.4}, \hyperref[exo-11-13]{11.13}\\[2mm]
Réunion, intersection et « ou » exclusif & \hyperref[exo-11-6]{11.6}, \hyperref[exo-11-9]{11.9}, \hyperref[exo-11-10]{11.10}, \hyperref[exo-11-14]{11.14}\\[2mm]
Dénombrer deux épreuves indépendantes & \hyperref[exo-11-7]{11.7}, \hyperref[exo-11-8]{11.8}, \hyperref[exo-11-10]{11.10}, \hyperref[exo-11-11]{11.11}, \hyperref[exo-11-14]{11.14}\\[2mm]
Fréquences et simulations & \hyperref[exo-11-5]{11.5}, \hyperref[exo-11-12]{11.12}, \hyperref[exo-11-15]{11.15}\\[2mm]
\bottomrule\end{tabularx}\par\vspace{4mm}
\otIntertitre{12. Pythagore, cercle et droite des milieux}\begin{tabularx}{\linewidth}{@{}Yl@{}}\toprule
Pythagore : calculer une longueur & \hyperref[exo-12-1]{12.1}, \hyperref[exo-12-2]{12.2}, \hyperref[exo-12-3]{12.3}, \hyperref[exo-12-4]{12.4}, \hyperref[exo-12-11]{12.11}, \hyperref[exo-12-14]{12.14}\\[2mm]
Réciproque et contraposée & \hyperref[exo-12-6]{12.6}, \hyperref[exo-12-7]{12.7}, \hyperref[exo-12-10]{12.10}\\[2mm]
Cercle circonscrit au triangle rectangle & \hyperref[exo-12-5]{12.5}, \hyperref[exo-12-12]{12.12}, \hyperref[exo-12-13]{12.13}\\[2mm]
Théorèmes des milieux & \hyperref[exo-12-8]{12.8}, \hyperref[exo-12-9]{12.9}, \hyperref[exo-12-15]{12.15}\\[2mm]
\bottomrule\end{tabularx}\par\vspace{4mm}
\clearpage
\otPageTitre{Repères et exercices}{Chapitres 13 à 15}\fontsize{10}{13}\selectfont
\otIntertitre{13. Thalès, réciproque et contraposée}\begin{tabularx}{\linewidth}{@{}Yl@{}}\toprule
Thalès emboîté et en papillon & \hyperref[exo-13-1]{13.1}, \hyperref[exo-13-2]{13.2}, \hyperref[exo-13-3]{13.3}, \hyperref[exo-13-4]{13.4}, \hyperref[exo-13-5]{13.5}, \hyperref[exo-13-8]{13.8}, \hyperref[exo-13-9]{13.9}, \hyperref[exo-13-13]{13.13}\\[2mm]
Réciproque du théorème de Thalès & \hyperref[exo-13-6]{13.6}, \hyperref[exo-13-10]{13.10}, \hyperref[exo-13-14]{13.14}\\[2mm]
Contraposée et non-parallélisme & \hyperref[exo-13-7]{13.7}, \hyperref[exo-13-15]{13.15}\\[2mm]
Modélisation et preuve par les aires & \hyperref[exo-13-11]{13.11}, \hyperref[exo-13-12]{13.12}, \hyperref[exo-13-14]{13.14}\\[2mm]
\bottomrule\end{tabularx}\par\vspace{4mm}
\otIntertitre{14. Trigonométrie du triangle rectangle}\begin{tabularx}{\linewidth}{@{}Yl@{}}\toprule
Choisir cosinus, sinus ou tangente & \hyperref[exo-14-1]{14.1}, \hyperref[exo-14-2]{14.2}, \hyperref[exo-14-3]{14.3}, \hyperref[exo-14-4]{14.4}, \hyperref[exo-14-15]{14.15}\\[2mm]
Calculer une longueur & \hyperref[exo-14-5]{14.5}, \hyperref[exo-14-6]{14.6}, \hyperref[exo-14-7]{14.7}, \hyperref[exo-14-13]{14.13}, \hyperref[exo-14-14]{14.14}\\[2mm]
Calculer un angle aigu & \hyperref[exo-14-8]{14.8}, \hyperref[exo-14-9]{14.9}, \hyperref[exo-14-12]{14.12}\\[2mm]
Relation $\cos^2+\sin^2=1$ et mesures indirectes & \hyperref[exo-14-10]{14.10}, \hyperref[exo-14-11]{14.11}, \hyperref[exo-14-14]{14.14}\\[2mm]
\bottomrule\end{tabularx}\par\vspace{4mm}
\otIntertitre{15. Translations et vecteurs}\begin{tabularx}{\linewidth}{@{}Yl@{}}\toprule
Symétrie, demi-tour et translation & \hyperref[exo-15-1]{15.1}, \hyperref[exo-15-4]{15.4}, \hyperref[exo-15-10]{15.10}, \hyperref[exo-15-14]{15.14}\\[2mm]
Vecteurs égaux, nul et opposé & \hyperref[exo-15-2]{15.2}, \hyperref[exo-15-3]{15.3}, \hyperref[exo-15-5]{15.5}, \hyperref[exo-15-12]{15.12}\\[2mm]
Parallélogramme et égalité de vecteurs & \hyperref[exo-15-4]{15.4}, \hyperref[exo-15-6]{15.6}, \hyperref[exo-15-7]{15.7}, \hyperref[exo-15-11]{15.11}, \hyperref[exo-15-15]{15.15}\\[2mm]
Somme de vecteurs et relation de Chasles & \hyperref[exo-15-8]{15.8}, \hyperref[exo-15-9]{15.9}, \hyperref[exo-15-10]{15.10}, \hyperref[exo-15-11]{15.11}, \hyperref[exo-15-13]{15.13}\\[2mm]
\bottomrule\end{tabularx}\par\vspace{4mm}
\clearpage
\otPageTitre{Repères et exercices}{Chapitres 16 à 18}\fontsize{10}{13}\selectfont
\otIntertitre{16. Espace : sphère, boule et sections}\begin{tabularx}{\linewidth}{@{}Yl@{}}\toprule
Volumes usuels et boule & \hyperref[exo-16-1]{16.1}, \hyperref[exo-16-2]{16.2}, \hyperref[exo-16-3]{16.3}, \hyperref[exo-16-4]{16.4}, \hyperref[exo-16-5]{16.5}, \hyperref[exo-16-7]{16.7}\\[2mm]
Sphère, boule et grands cercles & \hyperref[exo-16-5]{16.5}, \hyperref[exo-16-6]{16.6}, \hyperref[exo-16-10]{16.10}\\[2mm]
Sections planes de solides & \hyperref[exo-16-8]{16.8}, \hyperref[exo-16-9]{16.9}, \hyperref[exo-16-10]{16.10}, \hyperref[exo-16-13]{16.13}, \hyperref[exo-16-14]{16.14}\\[2mm]
Assemblages et problèmes de volumes & \hyperref[exo-16-11]{16.11}, \hyperref[exo-16-12]{16.12}, \hyperref[exo-16-13]{16.13}, \hyperref[exo-16-15]{16.15}\\[2mm]
\bottomrule\end{tabularx}\par\vspace{4mm}
\otIntertitre{17. Pensée informatique et algorithmique}\begin{tabularx}{\linewidth}{@{}Yl@{}}\toprule
Variables et exécution à la main & \hyperref[exo-17-1]{17.1}, \hyperref[exo-17-2]{17.2}, \hyperref[exo-17-6]{17.6}, \hyperref[exo-17-10]{17.10}\\[2mm]
Conditions simples et composées & \hyperref[exo-17-4]{17.4}, \hyperref[exo-17-8]{17.8}, \hyperref[exo-17-9]{17.9}, \hyperref[exo-17-14]{17.14}\\[2mm]
Boucles bornées et conditionnelles & \hyperref[exo-17-3]{17.3}, \hyperref[exo-17-5]{17.5}, \hyperref[exo-17-7]{17.7}, \hyperref[exo-17-10]{17.10}, \hyperref[exo-17-11]{17.11}\\[2mm]
Blocs personnalisés et simulations & \hyperref[exo-17-12]{17.12}, \hyperref[exo-17-13]{17.13}, \hyperref[exo-17-14]{17.14}, \hyperref[exo-17-15]{17.15}\\[2mm]
\bottomrule\end{tabularx}\par\vspace{4mm}
\otIntertitre{18. Raisonner et préparer le brevet}\begin{tabularx}{\linewidth}{@{}Yl@{}}\toprule
Conjecture, preuve et contre-exemple & \hyperref[exo-18-1]{18.1}, \hyperref[exo-18-2]{18.2}, \hyperref[exo-18-6]{18.6}, \hyperref[exo-18-8]{18.8}, \hyperref[exo-18-10]{18.10}\\[2mm]
Réciproque, contraposée et absurde & \hyperref[exo-18-3]{18.3}, \hyperref[exo-18-5]{18.5}, \hyperref[exo-18-7]{18.7}\\[2mm]
Problèmes de synthèse & \hyperref[exo-18-11]{18.11}, \hyperref[exo-18-12]{18.12}, \hyperref[exo-18-13]{18.13}, \hyperref[exo-18-14]{18.14}\\[2mm]
Méthode et format du brevet 2029 & \hyperref[exo-18-4]{18.4}, \hyperref[exo-18-9]{18.9}, \hyperref[exo-18-15]{18.15}\\[2mm]
\bottomrule\end{tabularx}\par\vspace{4mm}
\clearpage

\clearpage\thispagestyle{otFront}\phantomsection\label{references}\pdfbookmark[0]{Sources et édition}{bm-references}\otPageTitre{Sources et édition}{Références officielles · France}\fontsize{11}{15}\selectfont
\par\medskip\textbf{1. Arrêté du 18 février 2026 : calendrier}\par « Programmes d’enseignement de français et de mathématiques du cycle des approfondissements (cycle 4) », arrêté du 18-2-2026, JO du 4-3-2026, NOR MENE2602912A, Bulletin officiel n° 10 du 5 mars 2026. Article 3 : application en classe de troisième à la rentrée de l’année scolaire 2028-2029.\par\href{https://www.education.gouv.fr/sites/default/files/document/Bulletin\%20officiel\%20n\%C2\%B0\%2010\%20du\%205\%20mars\%202026-480719.pdf}{\textcolor{otGreenDark}{Consulter la source officielle}}\par
\par\medskip\textbf{2. Programme de mathématiques du cycle 4}\par Annexe 2 de cet arrêté, « Programme de mathématiques pour le cycle 4 » : principes, puis rubriques de Troisième (automatismes, objectifs d’apprentissage et prolongements possibles).\par\href{https://www.education.gouv.fr/sites/default/files/document/Annexe\%202\%20\%E2\%80\%93\%20Programme\%20de\%20math\%C3\%A9matiques\%20pour\%20le\%20cycle\%204-480716.pdf}{\textcolor{otGreenDark}{Consulter la source officielle}}\par
\par\medskip\textbf{3. Exemples pour la mise en œuvre}\par Éduscol, « Mathématiques. Exemples pour la mise en œuvre des programmes 2026. Cycle 4 » : exemples de réussite et conseils de mise en œuvre pour chaque objectif de Troisième.\par\href{https://eduscol.education.gouv.fr/sites/default/files/document/exemple-mise-en-oeuvre-c4-mathematiques0pdf-126080.pdf}{\textcolor{otGreenDark}{Consulter la source officielle}}\par
\par\medskip\textbf{4. Définition de l’épreuve du brevet}\par Note de service du 11-9-2026, NOR MENE2623228N, « Définition des épreuves de l’examen du diplôme national du brevet à compter de la session 2027 », Bulletin officiel spécial n° 4 du 17 septembre 2026 ; partie 2 : épreuve écrite de mathématiques.\par\href{https://www.education.gouv.fr/sites/default/files/document/bulletin-officiel-special-ndeg-4-du-17-septembre-2026-521113.pdf}{\textcolor{otGreenDark}{Consulter la source officielle}}\par
\par\medskip\textbf{5. Maîtrise de la langue aux examens}\par Note de service du 11-9-2026, NOR MENE2623195N, « Exigences en matière de maîtrise de la langue à tous les examens à compter de la session 2027 », même Bulletin officiel spécial n° 4.\par\href{https://www.education.gouv.fr/sites/default/files/document/bulletin-officiel-special-ndeg-4-du-17-septembre-2026-521113.pdf}{\textcolor{otGreenDark}{Consulter la source officielle}}\par
\par\bigskip Sources consultées le 23 octobre 2026. Le programme et les exemples
de mise en œuvre ont guidé le périmètre et la matrice de correspondance. Les
exercices sont originaux, non reproduits des documents officiels. Les défis et
les sujets du livre ne sont pas des sujets officiels.\par\medskip
À cette date, aucun texte propre à la session 2029 du brevet n’était publié.
La note du 11 septembre 2026 s’applique « à compter de la session 2027 » et
précise que les sujets portent sur le programme de Troisième s’il est en vigueur :
c’est la définition retenue pour les brevets blancs. Un texte ultérieur pourrait
la modifier.

\clearpage

%% =====================================================================
%%  4e de couverture
%% =====================================================================
\clearpage
\thispagestyle{empty}
\begin{tikzpicture}[remember picture,overlay]
\begin{scope}[shift={(current page.south west)}]
  % Fond
  \fill[otNavy] (0,0) rectangle (210mm,297mm);
  \shade[bottom color=otNavyDark,top color=otNavy] (0,0) rectangle (210mm,160mm);
  \draw[white,opacity=0.06,line width=0.3pt,step=7mm] (0,0) grid (210mm,297mm);

  % Accroche et texte
  \node[anchor=north west,inner sep=0pt,text width=150mm,align=left] at (24mm,270mm)
    {{\otSemi\fontsize{9}{11}\selectfont\color{otGreen}%
       \otCaps[18]{\otTitreLivre\ · \otNiveau}}\par\vspace{7mm}
     {\otSemi\fontsize{25}{31}\selectfont\color{white}%
       Réussir sa Troisième,\\exercice après exercice.}\par\vspace{9mm}
     {\fontsize{11.5}{16}\selectfont\color{white!88!otNavy}%
       Cette édition suit le nouveau programme de mathématiques du cycle 4,
       applicable en Troisième à la rentrée 2028. Elle réunit les cours et
       288 exercices corrigés : 270 exercices progressifs dans 18 chapitres
       et 18 défis facultatifs. Vecteurs, quartiles, identités remarquables,
       inéquations, sphère et sections sont traités en détail. Deux brevets
       blancs originaux préparent la session 2029.\par}\vspace{9mm}
     {\otRegular\fontsize{10.5}{14}\selectfont\color{white}%
       \otCheck\hspace{3mm}Des bases accessibles aux défis très exigeants\par\vspace{3.4mm}
       \otCheck\hspace{3mm}Des corrigés rédigés, avec la méthode à retenir\par\vspace{3.4mm}
       \otCheck\hspace{3mm}Des figures, tableaux et scripts par blocs soignés\par\vspace{3.4mm}
       \otCheck\hspace{3mm}Un tableau de bord pour suivre sa progression\par}};

  % Bloc « en ligne »
  \fill[white,rounded corners=3mm] (24mm,62mm) rectangle (186mm,98mm);
  \node[anchor=west,inner sep=0pt] at (30mm,80mm)
    {\color{otNavy}\qrcode[height=26mm,nolink]{\otURL}};
  \node[anchor=west,inner sep=0pt,text width=110mm] at (64mm,80mm)
    {{\otSemi\fontsize{12}{15}\selectfont\color{otNavy}Collection OptiTAB}\par\vspace{2mm}
     {\fontsize{10.5}{14}\selectfont\color{otInk}Du collège au lycée : une même présentation pour\par apprendre, chercher et progresser.}};

  % Logo, prix et code-barres
  \node[anchor=south west,inner sep=0pt] at (24mm,20mm)
    {\includegraphics[width=46mm]{logo-blanc}};
  \fill[white] (138mm,16mm) rectangle (186mm,44mm);
  \node[text=otMuted,align=center,font=\otLight\fontsize{7}{9}\selectfont]
    at (162mm,30mm) {NOUVEAU PROGRAMME\\18 chapitres\\Édition 2028-2029};
\end{scope}
\end{tikzpicture}

\end{document}
